\documentclass[11pt,a4paper]{article}
% preamble.tex - shared house preamble for /explain-projects documents.
% Every package here is VERIFIED present in this machine's TeX Live install.
% Do not add packages without running: kpsewhich <pkg>.sty
% Notably ABSENT (do not \usepackage them, the build will fail):
%   algorithm2e, algorithm, algorithmic, algpseudocode  (texlive-science is not installed)
%   -> use the \begin{pseudocode} tcolorbox environment defined below instead.
%   lmodern, charter, mathptmx, times, mathpazo, courier, newtx*, tgtermes
%   -> their .sty files may resolve but the FONT FILES are absent, so the build
%      dies with "I can't find file `bchr8t'" or similar. Verified 2026-07-16 by
%      test-compiling each one. Only Computer Modern (the default) and helvet
%      actually render on this box. Do not "upgrade" the font stack without
%      test-compiling first: kpsewhich lying about a font is the exact trap here.
%
% To get a nicer font stack + algorithm2e, install:
%   sudo apt install texlive-fonts-recommended texlive-science
% Until then, this preamble deliberately uses the default font.
%
% Usage from a document:
%   \documentclass[11pt,a4paper]{article}
%   % preamble.tex - shared house preamble for /explain-projects documents.
% Every package here is VERIFIED present in this machine's TeX Live install.
% Do not add packages without running: kpsewhich <pkg>.sty
% Notably ABSENT (do not \usepackage them, the build will fail):
%   algorithm2e, algorithm, algorithmic, algpseudocode  (texlive-science is not installed)
%   -> use the \begin{pseudocode} tcolorbox environment defined below instead.
%   lmodern, charter, mathptmx, times, mathpazo, courier, newtx*, tgtermes
%   -> their .sty files may resolve but the FONT FILES are absent, so the build
%      dies with "I can't find file `bchr8t'" or similar. Verified 2026-07-16 by
%      test-compiling each one. Only Computer Modern (the default) and helvet
%      actually render on this box. Do not "upgrade" the font stack without
%      test-compiling first: kpsewhich lying about a font is the exact trap here.
%
% To get a nicer font stack + algorithm2e, install:
%   sudo apt install texlive-fonts-recommended texlive-science
% Until then, this preamble deliberately uses the default font.
%
% Usage from a document:
%   \documentclass[11pt,a4paper]{article}
%   % preamble.tex - shared house preamble for /explain-projects documents.
% Every package here is VERIFIED present in this machine's TeX Live install.
% Do not add packages without running: kpsewhich <pkg>.sty
% Notably ABSENT (do not \usepackage them, the build will fail):
%   algorithm2e, algorithm, algorithmic, algpseudocode  (texlive-science is not installed)
%   -> use the \begin{pseudocode} tcolorbox environment defined below instead.
%   lmodern, charter, mathptmx, times, mathpazo, courier, newtx*, tgtermes
%   -> their .sty files may resolve but the FONT FILES are absent, so the build
%      dies with "I can't find file `bchr8t'" or similar. Verified 2026-07-16 by
%      test-compiling each one. Only Computer Modern (the default) and helvet
%      actually render on this box. Do not "upgrade" the font stack without
%      test-compiling first: kpsewhich lying about a font is the exact trap here.
%
% To get a nicer font stack + algorithm2e, install:
%   sudo apt install texlive-fonts-recommended texlive-science
% Until then, this preamble deliberately uses the default font.
%
% Usage from a document:
%   \documentclass[11pt,a4paper]{article}
%   \input{preamble}          % this file is copied next to the .tex

\usepackage[T1]{fontenc}
\usepackage[utf8]{inputenc}
% Font: Computer Modern (default). See the note above before changing this.
\usepackage[a4paper,margin=2.4cm,headheight=14pt]{geometry}
% microtype WITHOUT font expansion. This is load-bearing, do not "restore" it:
% cm-super is not installed, so T1 Computer Modern resolves to BITMAP fonts, and
% pdfTeX's auto-expansion needs scalable ones. With expansion on, any document
% long enough to reach a second page dies with
%   "pdfTeX error (font expansion): auto expansion is only possible with scalable fonts"
% Protrusion (the bigger visual win) still works and stays on.
% `sudo apt install cm-super` would lift this; until then, leave expansion off.
\usepackage[protrusion=true,expansion=false]{microtype}
\usepackage{graphicx}
\usepackage{xcolor}
\usepackage{booktabs}
\usepackage{enumitem}
\usepackage{fancyhdr}
\usepackage{titlesec}
\usepackage{tcolorbox}
\usepackage{listings}
\usepackage{float}
\usepackage{caption}
\usepackage{tikz}
\usepackage{forest}
\usepackage{pgfplots}
\pgfplotsset{compat=1.18}
\usetikzlibrary{arrows.meta,positioning,shapes.geometric,fit,backgrounds,calc,decorations.pathreplacing}
\tcbuselibrary{skins,breakable}
\usepackage[hidelinks,bookmarks=true,bookmarksnumbered=true]{hyperref}

% ===== palette ===========================================================
\definecolor{inkblue}{HTML}{1F3A5F}
\definecolor{accent}{HTML}{2E7DAF}
\definecolor{softgrey}{HTML}{F4F6F8}
\definecolor{linegrey}{HTML}{C8D0D8}
\definecolor{codebg}{HTML}{F7F7F4}
\definecolor{warnred}{HTML}{A33B3B}
\definecolor{okgreen}{HTML}{2E7D53}

% ===== headings ==========================================================
\titleformat{\section}{\Large\bfseries\color{inkblue}}{\thesection}{0.7em}{}
\titleformat{\subsection}{\large\bfseries\color{inkblue}}{\thesubsection}{0.6em}{}
\titleformat{\subsubsection}{\normalsize\bfseries\color{accent}}{\thesubsubsection}{0.5em}{}

\pagestyle{fancy}
\fancyhf{}
\fancyhead[L]{\small\color{inkblue}\nouppercase{\leftmark}}
\fancyhead[R]{\small\color{inkblue}\thepage}
\renewcommand{\headrulewidth}{0.4pt}
\renewcommand{\headrule}{\hbox to\headwidth{\color{linegrey}\leaders\hrule height \headrulewidth\hfill}}

% ===== code ==============================================================
\lstdefinestyle{house}{
  backgroundcolor=\color{codebg},
  basicstyle=\ttfamily\footnotesize,
  keywordstyle=\color{inkblue}\bfseries,
  commentstyle=\color{gray}\itshape,
  stringstyle=\color{okgreen},
  numbers=left, numberstyle=\tiny\color{gray}, numbersep=8pt,
  frame=single, rulecolor=\color{linegrey},
  breaklines=true, breakatwhitespace=false,
  showstringspaces=false, tabsize=2,
  captionpos=b, xleftmargin=14pt, framexleftmargin=14pt,
  columns=fullflexible, keepspaces=true,
  upquote=true  % render ' and " as straight quotes, not curly
}
\lstset{style=house}

% ===== callout boxes =====================================================
% \begin{keyidea}{Title} ... \end{keyidea}   - a concept a beginner must grasp
\newtcolorbox{keyidea}[1]{
  enhanced, breakable, colback=softgrey, colframe=accent,
  boxrule=0.4pt, arc=2pt, left=8pt, right=8pt, top=6pt, bottom=6pt,
  title={\bfseries #1}, coltitle=white, colbacktitle=accent, fonttitle=\small
}

% \begin{jargon}{Term} definition \end{jargon}  - define a term at first use
\newtcolorbox{jargon}[1]{
  enhanced, breakable, colback=white, colframe=linegrey,
  boxrule=0.4pt, arc=2pt, left=8pt, right=8pt, top=6pt, bottom=6pt,
  title={\bfseries Jargon: #1}, coltitle=inkblue, colbacktitle=softgrey, fonttitle=\small
}

% \begin{honest}{Title} ... \end{honest}  - brutally honest finding (rule 18)
\newtcolorbox{honest}[1]{
  enhanced, breakable, colback=white, colframe=warnred,
  boxrule=0.6pt, arc=2pt, left=8pt, right=8pt, top=6pt, bottom=6pt,
  title={\bfseries #1}, coltitle=white, colbacktitle=warnred, fonttitle=\small
}

% \begin{provenance}{path/to/file} ... \end{provenance}
% Where a file came from: hand-written, generated by a tool, scaffolded by a
% framework, copied or adapted from an upstream project, ported from an older
% module. State the EVIDENCE (a header comment, a git log entry, an upstream
% URL, a generator command) and label the confidence. If no evidence exists,
% say so in the box. Never guess a lineage into existence.
\newtcolorbox{provenance}[1]{
  enhanced, breakable, colback=white, colframe=okgreen,
  boxrule=0.4pt, arc=2pt, left=8pt, right=8pt, top=6pt, bottom=6pt,
  title={\bfseries Where this came from: #1}, coltitle=white,
  colbacktitle=okgreen, fonttitle=\small
}

% \begin{pseudocode}{Caption} ... \end{pseudocode}
% Stands in for algorithm2e, which is NOT installed. Use \Step / \Indent inside.
\newtcolorbox{pseudocode}[1]{
  enhanced, breakable, colback=codebg, colframe=inkblue,
  boxrule=0.4pt, arc=2pt, left=10pt, right=8pt, top=6pt, bottom=6pt,
  fontupper=\ttfamily\small,
  title={\bfseries\rmfamily #1}, coltitle=white, colbacktitle=inkblue, fonttitle=\small
}
\newcommand{\Step}[1]{\par\noindent #1}
\newcommand{\Indent}[1]{\par\noindent\hspace*{2em}#1}

% ===== tikz house styles =================================================
\tikzset{
  box/.style   = {rectangle, rounded corners=2pt, draw=inkblue, fill=softgrey,
                  text=inkblue, align=center, inner sep=6pt, minimum height=9mm,
                  font=\small},
  extbox/.style= {rectangle, rounded corners=2pt, draw=linegrey, fill=white,
                  text=black!70, align=center, inner sep=6pt, minimum height=9mm,
                  font=\small, dashed},
  store/.style = {cylinder, shape border rotate=90, aspect=0.25, draw=accent,
                  fill=white, text=inkblue, align=center, inner sep=5pt,
                  font=\small},
  dec/.style   = {diamond, aspect=1.6, draw=accent, fill=white, text=inkblue,
                  align=center, inner sep=2pt, font=\scriptsize},
  flow/.style  = {-{Stealth[length=2.4mm]}, draw=inkblue, thick},
  dflow/.style = {-{Stealth[length=2.4mm]}, draw=accent, thick, dashed},
  % align=center is load-bearing: without it, a \\ inside an edge label throws
  % the misleading "perhaps a missing \item" error.
  lbl/.style   = {font=\scriptsize\itshape, text=black!65, fill=white,
                  inner sep=2pt, align=center}
}

% forest default for directory trees
\forestset{
  dirtree/.style={
    for tree={
      font=\ttfamily\small, grow'=0, child anchor=west, parent anchor=south,
      anchor=west, calign=first, inner sep=2pt,
      edge path={\noexpand\path [draw, \forestoption{edge}]
        (!u.south west) +(9pt,0) |- (.child anchor) \forestoption{edge label};},
      before typesetting nodes={if n=1{insert before={[,phantom]}}{}},
      fit=band, before computing xy={l=15pt}
    }
  }
}

% ===== misc ==============================================================
\captionsetup{font=small, labelfont={bf,color=inkblue}}
\setlist[itemize]{leftmargin=1.2em, itemsep=2pt, topsep=3pt}
\setlist[enumerate]{leftmargin=1.4em, itemsep=2pt, topsep=3pt}
\newcommand{\file}[1]{\texttt{\small #1}}
\newcommand{\code}[1]{\texttt{\small #1}}
\setcounter{tocdepth}{2}
\setcounter{secnumdepth}{3}
          % this file is copied next to the .tex

\usepackage[T1]{fontenc}
\usepackage[utf8]{inputenc}
% Font: Computer Modern (default). See the note above before changing this.
\usepackage[a4paper,margin=2.4cm,headheight=14pt]{geometry}
% microtype WITHOUT font expansion. This is load-bearing, do not "restore" it:
% cm-super is not installed, so T1 Computer Modern resolves to BITMAP fonts, and
% pdfTeX's auto-expansion needs scalable ones. With expansion on, any document
% long enough to reach a second page dies with
%   "pdfTeX error (font expansion): auto expansion is only possible with scalable fonts"
% Protrusion (the bigger visual win) still works and stays on.
% `sudo apt install cm-super` would lift this; until then, leave expansion off.
\usepackage[protrusion=true,expansion=false]{microtype}
\usepackage{graphicx}
\usepackage{xcolor}
\usepackage{booktabs}
\usepackage{enumitem}
\usepackage{fancyhdr}
\usepackage{titlesec}
\usepackage{tcolorbox}
\usepackage{listings}
\usepackage{float}
\usepackage{caption}
\usepackage{tikz}
\usepackage{forest}
\usepackage{pgfplots}
\pgfplotsset{compat=1.18}
\usetikzlibrary{arrows.meta,positioning,shapes.geometric,fit,backgrounds,calc,decorations.pathreplacing}
\tcbuselibrary{skins,breakable}
\usepackage[hidelinks,bookmarks=true,bookmarksnumbered=true]{hyperref}

% ===== palette ===========================================================
\definecolor{inkblue}{HTML}{1F3A5F}
\definecolor{accent}{HTML}{2E7DAF}
\definecolor{softgrey}{HTML}{F4F6F8}
\definecolor{linegrey}{HTML}{C8D0D8}
\definecolor{codebg}{HTML}{F7F7F4}
\definecolor{warnred}{HTML}{A33B3B}
\definecolor{okgreen}{HTML}{2E7D53}

% ===== headings ==========================================================
\titleformat{\section}{\Large\bfseries\color{inkblue}}{\thesection}{0.7em}{}
\titleformat{\subsection}{\large\bfseries\color{inkblue}}{\thesubsection}{0.6em}{}
\titleformat{\subsubsection}{\normalsize\bfseries\color{accent}}{\thesubsubsection}{0.5em}{}

\pagestyle{fancy}
\fancyhf{}
\fancyhead[L]{\small\color{inkblue}\nouppercase{\leftmark}}
\fancyhead[R]{\small\color{inkblue}\thepage}
\renewcommand{\headrulewidth}{0.4pt}
\renewcommand{\headrule}{\hbox to\headwidth{\color{linegrey}\leaders\hrule height \headrulewidth\hfill}}

% ===== code ==============================================================
\lstdefinestyle{house}{
  backgroundcolor=\color{codebg},
  basicstyle=\ttfamily\footnotesize,
  keywordstyle=\color{inkblue}\bfseries,
  commentstyle=\color{gray}\itshape,
  stringstyle=\color{okgreen},
  numbers=left, numberstyle=\tiny\color{gray}, numbersep=8pt,
  frame=single, rulecolor=\color{linegrey},
  breaklines=true, breakatwhitespace=false,
  showstringspaces=false, tabsize=2,
  captionpos=b, xleftmargin=14pt, framexleftmargin=14pt,
  columns=fullflexible, keepspaces=true,
  upquote=true  % render ' and " as straight quotes, not curly
}
\lstset{style=house}

% ===== callout boxes =====================================================
% \begin{keyidea}{Title} ... \end{keyidea}   - a concept a beginner must grasp
\newtcolorbox{keyidea}[1]{
  enhanced, breakable, colback=softgrey, colframe=accent,
  boxrule=0.4pt, arc=2pt, left=8pt, right=8pt, top=6pt, bottom=6pt,
  title={\bfseries #1}, coltitle=white, colbacktitle=accent, fonttitle=\small
}

% \begin{jargon}{Term} definition \end{jargon}  - define a term at first use
\newtcolorbox{jargon}[1]{
  enhanced, breakable, colback=white, colframe=linegrey,
  boxrule=0.4pt, arc=2pt, left=8pt, right=8pt, top=6pt, bottom=6pt,
  title={\bfseries Jargon: #1}, coltitle=inkblue, colbacktitle=softgrey, fonttitle=\small
}

% \begin{honest}{Title} ... \end{honest}  - brutally honest finding (rule 18)
\newtcolorbox{honest}[1]{
  enhanced, breakable, colback=white, colframe=warnred,
  boxrule=0.6pt, arc=2pt, left=8pt, right=8pt, top=6pt, bottom=6pt,
  title={\bfseries #1}, coltitle=white, colbacktitle=warnred, fonttitle=\small
}

% \begin{provenance}{path/to/file} ... \end{provenance}
% Where a file came from: hand-written, generated by a tool, scaffolded by a
% framework, copied or adapted from an upstream project, ported from an older
% module. State the EVIDENCE (a header comment, a git log entry, an upstream
% URL, a generator command) and label the confidence. If no evidence exists,
% say so in the box. Never guess a lineage into existence.
\newtcolorbox{provenance}[1]{
  enhanced, breakable, colback=white, colframe=okgreen,
  boxrule=0.4pt, arc=2pt, left=8pt, right=8pt, top=6pt, bottom=6pt,
  title={\bfseries Where this came from: #1}, coltitle=white,
  colbacktitle=okgreen, fonttitle=\small
}

% \begin{pseudocode}{Caption} ... \end{pseudocode}
% Stands in for algorithm2e, which is NOT installed. Use \Step / \Indent inside.
\newtcolorbox{pseudocode}[1]{
  enhanced, breakable, colback=codebg, colframe=inkblue,
  boxrule=0.4pt, arc=2pt, left=10pt, right=8pt, top=6pt, bottom=6pt,
  fontupper=\ttfamily\small,
  title={\bfseries\rmfamily #1}, coltitle=white, colbacktitle=inkblue, fonttitle=\small
}
\newcommand{\Step}[1]{\par\noindent #1}
\newcommand{\Indent}[1]{\par\noindent\hspace*{2em}#1}

% ===== tikz house styles =================================================
\tikzset{
  box/.style   = {rectangle, rounded corners=2pt, draw=inkblue, fill=softgrey,
                  text=inkblue, align=center, inner sep=6pt, minimum height=9mm,
                  font=\small},
  extbox/.style= {rectangle, rounded corners=2pt, draw=linegrey, fill=white,
                  text=black!70, align=center, inner sep=6pt, minimum height=9mm,
                  font=\small, dashed},
  store/.style = {cylinder, shape border rotate=90, aspect=0.25, draw=accent,
                  fill=white, text=inkblue, align=center, inner sep=5pt,
                  font=\small},
  dec/.style   = {diamond, aspect=1.6, draw=accent, fill=white, text=inkblue,
                  align=center, inner sep=2pt, font=\scriptsize},
  flow/.style  = {-{Stealth[length=2.4mm]}, draw=inkblue, thick},
  dflow/.style = {-{Stealth[length=2.4mm]}, draw=accent, thick, dashed},
  % align=center is load-bearing: without it, a \\ inside an edge label throws
  % the misleading "perhaps a missing \item" error.
  lbl/.style   = {font=\scriptsize\itshape, text=black!65, fill=white,
                  inner sep=2pt, align=center}
}

% forest default for directory trees
\forestset{
  dirtree/.style={
    for tree={
      font=\ttfamily\small, grow'=0, child anchor=west, parent anchor=south,
      anchor=west, calign=first, inner sep=2pt,
      edge path={\noexpand\path [draw, \forestoption{edge}]
        (!u.south west) +(9pt,0) |- (.child anchor) \forestoption{edge label};},
      before typesetting nodes={if n=1{insert before={[,phantom]}}{}},
      fit=band, before computing xy={l=15pt}
    }
  }
}

% ===== misc ==============================================================
\captionsetup{font=small, labelfont={bf,color=inkblue}}
\setlist[itemize]{leftmargin=1.2em, itemsep=2pt, topsep=3pt}
\setlist[enumerate]{leftmargin=1.4em, itemsep=2pt, topsep=3pt}
\newcommand{\file}[1]{\texttt{\small #1}}
\newcommand{\code}[1]{\texttt{\small #1}}
\setcounter{tocdepth}{2}
\setcounter{secnumdepth}{3}
          % this file is copied next to the .tex

\usepackage[T1]{fontenc}
\usepackage[utf8]{inputenc}
% Font: Computer Modern (default). See the note above before changing this.
\usepackage[a4paper,margin=2.4cm,headheight=14pt]{geometry}
% microtype WITHOUT font expansion. This is load-bearing, do not "restore" it:
% cm-super is not installed, so T1 Computer Modern resolves to BITMAP fonts, and
% pdfTeX's auto-expansion needs scalable ones. With expansion on, any document
% long enough to reach a second page dies with
%   "pdfTeX error (font expansion): auto expansion is only possible with scalable fonts"
% Protrusion (the bigger visual win) still works and stays on.
% `sudo apt install cm-super` would lift this; until then, leave expansion off.
\usepackage[protrusion=true,expansion=false]{microtype}
\usepackage{graphicx}
\usepackage{xcolor}
\usepackage{booktabs}
\usepackage{enumitem}
\usepackage{fancyhdr}
\usepackage{titlesec}
\usepackage{tcolorbox}
\usepackage{listings}
\usepackage{float}
\usepackage{caption}
\usepackage{tikz}
\usepackage{forest}
\usepackage{pgfplots}
\pgfplotsset{compat=1.18}
\usetikzlibrary{arrows.meta,positioning,shapes.geometric,fit,backgrounds,calc,decorations.pathreplacing}
\tcbuselibrary{skins,breakable}
\usepackage[hidelinks,bookmarks=true,bookmarksnumbered=true]{hyperref}

% ===== palette ===========================================================
\definecolor{inkblue}{HTML}{1F3A5F}
\definecolor{accent}{HTML}{2E7DAF}
\definecolor{softgrey}{HTML}{F4F6F8}
\definecolor{linegrey}{HTML}{C8D0D8}
\definecolor{codebg}{HTML}{F7F7F4}
\definecolor{warnred}{HTML}{A33B3B}
\definecolor{okgreen}{HTML}{2E7D53}

% ===== headings ==========================================================
\titleformat{\section}{\Large\bfseries\color{inkblue}}{\thesection}{0.7em}{}
\titleformat{\subsection}{\large\bfseries\color{inkblue}}{\thesubsection}{0.6em}{}
\titleformat{\subsubsection}{\normalsize\bfseries\color{accent}}{\thesubsubsection}{0.5em}{}

\pagestyle{fancy}
\fancyhf{}
\fancyhead[L]{\small\color{inkblue}\nouppercase{\leftmark}}
\fancyhead[R]{\small\color{inkblue}\thepage}
\renewcommand{\headrulewidth}{0.4pt}
\renewcommand{\headrule}{\hbox to\headwidth{\color{linegrey}\leaders\hrule height \headrulewidth\hfill}}

% ===== code ==============================================================
\lstdefinestyle{house}{
  backgroundcolor=\color{codebg},
  basicstyle=\ttfamily\footnotesize,
  keywordstyle=\color{inkblue}\bfseries,
  commentstyle=\color{gray}\itshape,
  stringstyle=\color{okgreen},
  numbers=left, numberstyle=\tiny\color{gray}, numbersep=8pt,
  frame=single, rulecolor=\color{linegrey},
  breaklines=true, breakatwhitespace=false,
  showstringspaces=false, tabsize=2,
  captionpos=b, xleftmargin=14pt, framexleftmargin=14pt,
  columns=fullflexible, keepspaces=true,
  upquote=true  % render ' and " as straight quotes, not curly
}
\lstset{style=house}

% ===== callout boxes =====================================================
% \begin{keyidea}{Title} ... \end{keyidea}   - a concept a beginner must grasp
\newtcolorbox{keyidea}[1]{
  enhanced, breakable, colback=softgrey, colframe=accent,
  boxrule=0.4pt, arc=2pt, left=8pt, right=8pt, top=6pt, bottom=6pt,
  title={\bfseries #1}, coltitle=white, colbacktitle=accent, fonttitle=\small
}

% \begin{jargon}{Term} definition \end{jargon}  - define a term at first use
\newtcolorbox{jargon}[1]{
  enhanced, breakable, colback=white, colframe=linegrey,
  boxrule=0.4pt, arc=2pt, left=8pt, right=8pt, top=6pt, bottom=6pt,
  title={\bfseries Jargon: #1}, coltitle=inkblue, colbacktitle=softgrey, fonttitle=\small
}

% \begin{honest}{Title} ... \end{honest}  - brutally honest finding (rule 18)
\newtcolorbox{honest}[1]{
  enhanced, breakable, colback=white, colframe=warnred,
  boxrule=0.6pt, arc=2pt, left=8pt, right=8pt, top=6pt, bottom=6pt,
  title={\bfseries #1}, coltitle=white, colbacktitle=warnred, fonttitle=\small
}

% \begin{provenance}{path/to/file} ... \end{provenance}
% Where a file came from: hand-written, generated by a tool, scaffolded by a
% framework, copied or adapted from an upstream project, ported from an older
% module. State the EVIDENCE (a header comment, a git log entry, an upstream
% URL, a generator command) and label the confidence. If no evidence exists,
% say so in the box. Never guess a lineage into existence.
\newtcolorbox{provenance}[1]{
  enhanced, breakable, colback=white, colframe=okgreen,
  boxrule=0.4pt, arc=2pt, left=8pt, right=8pt, top=6pt, bottom=6pt,
  title={\bfseries Where this came from: #1}, coltitle=white,
  colbacktitle=okgreen, fonttitle=\small
}

% \begin{pseudocode}{Caption} ... \end{pseudocode}
% Stands in for algorithm2e, which is NOT installed. Use \Step / \Indent inside.
\newtcolorbox{pseudocode}[1]{
  enhanced, breakable, colback=codebg, colframe=inkblue,
  boxrule=0.4pt, arc=2pt, left=10pt, right=8pt, top=6pt, bottom=6pt,
  fontupper=\ttfamily\small,
  title={\bfseries\rmfamily #1}, coltitle=white, colbacktitle=inkblue, fonttitle=\small
}
\newcommand{\Step}[1]{\par\noindent #1}
\newcommand{\Indent}[1]{\par\noindent\hspace*{2em}#1}

% ===== tikz house styles =================================================
\tikzset{
  box/.style   = {rectangle, rounded corners=2pt, draw=inkblue, fill=softgrey,
                  text=inkblue, align=center, inner sep=6pt, minimum height=9mm,
                  font=\small},
  extbox/.style= {rectangle, rounded corners=2pt, draw=linegrey, fill=white,
                  text=black!70, align=center, inner sep=6pt, minimum height=9mm,
                  font=\small, dashed},
  store/.style = {cylinder, shape border rotate=90, aspect=0.25, draw=accent,
                  fill=white, text=inkblue, align=center, inner sep=5pt,
                  font=\small},
  dec/.style   = {diamond, aspect=1.6, draw=accent, fill=white, text=inkblue,
                  align=center, inner sep=2pt, font=\scriptsize},
  flow/.style  = {-{Stealth[length=2.4mm]}, draw=inkblue, thick},
  dflow/.style = {-{Stealth[length=2.4mm]}, draw=accent, thick, dashed},
  % align=center is load-bearing: without it, a \\ inside an edge label throws
  % the misleading "perhaps a missing \item" error.
  lbl/.style   = {font=\scriptsize\itshape, text=black!65, fill=white,
                  inner sep=2pt, align=center}
}

% forest default for directory trees
\forestset{
  dirtree/.style={
    for tree={
      font=\ttfamily\small, grow'=0, child anchor=west, parent anchor=south,
      anchor=west, calign=first, inner sep=2pt,
      edge path={\noexpand\path [draw, \forestoption{edge}]
        (!u.south west) +(9pt,0) |- (.child anchor) \forestoption{edge label};},
      before typesetting nodes={if n=1{insert before={[,phantom]}}{}},
      fit=band, before computing xy={l=15pt}
    }
  }
}

% ===== misc ==============================================================
\captionsetup{font=small, labelfont={bf,color=inkblue}}
\setlist[itemize]{leftmargin=1.2em, itemsep=2pt, topsep=3pt}
\setlist[enumerate]{leftmargin=1.4em, itemsep=2pt, topsep=3pt}
\newcommand{\file}[1]{\texttt{\small #1}}
\newcommand{\code}[1]{\texttt{\small #1}}
\setcounter{tocdepth}{2}
\setcounter{secnumdepth}{3}


% ===========================================================================
% Two characters in this repository's source files cannot be typeset by the
% font this document uses, so they are mapped to the closest faithful stand-in.
% The mapping is stated openly in section 1.4 so no reader is misled by it.
%   U+2500  horizontal box-drawing rule  ->  hyphen
%   U+2550  double box-drawing rule      ->  equals sign
% U+00B7 (middle dot) renders exactly and needs no substitution.
% ===========================================================================
\DeclareUnicodeCharacter{2500}{-}
\DeclareUnicodeCharacter{2550}{=}
\lstset{extendedchars=true,
  literate={·}{{\textperiodcentered}}1
           {─}{{-}}1
           {═}{{=}}1
}

\title{Salman Adnan Site\\[2mm]\large Deep Dive: every folder, every file, every line}
\author{Digitalise Agency}
\date{2026-08-11}

\begin{document}
\maketitle
\tableofcontents
\newpage

% ===========================================================================
\section{What this project is}
% ===========================================================================

This repository holds a personal brand website for \textbf{salmanadnan.com}: a
single page introducing Salman Adnan, showing six client projects, and asking the
visitor to book a thirty-minute call.

It is six files and 1{,}180 lines. No framework, no build step, no package
manager, no database, no server-side code. A browser can open the page straight
off a disk.

\begin{keyidea}{The one thing to understand before anything else}
This repository does \emph{not} contain the page currently served at
salmanadnan.com. The design here was finished, then deliberately parked, and the
machinery that published it was switched off in the same commit. Something else
serves that address today. Section 1.2 has the evidence.
\end{keyidea}

\begin{keyidea}{Its twin, and why you should read both}
There is a near-identical sibling repository in this workspace,
\code{ai-digitalise}, built in the same week from the same skeleton. It has
exactly the same six file paths, the same three-file architecture, the same
publishing workflow and the same fate.

The two are not copies: every substantive file differs. This one is a personal
brand page in orange; that one is an agency service page in blue. Where a
difference between them is instructive, this document points it out, because
comparing the pair teaches more than reading either alone.
\end{keyidea}

\subsection{The problem it was built to solve}

A consultant needs a page to send a prospective client to. It has to say who he
is, prove he has done this before, and make booking a conversation one click.
That is the whole purpose.

\begin{figure}[H]\centering
\begin{tikzpicture}[node distance=10mm and 18mm]
  \node[box]                (v)  {A visitor\\\scriptsize from a link, a profile or a referral};
  \node[box, right=of v]    (p)  {The page\\\scriptsize \file{index.html}};
  \node[extbox, right=of p] (cal){cal.com\\\scriptsize \code{salman-adnan/meeting}};
  \node[box, below=of p]    (out){A booked call\\\scriptsize the only outcome that counts};
  \draw[flow]  (v) -- node[lbl]{opens} (p);
  \draw[dflow] (p) -- node[lbl]{Book a call} (cal);
  \draw[flow]  (cal) |- node[lbl, pos=0.25]{picks a time} (out);
\end{tikzpicture}
\caption{The whole project. The dashed box is a third-party service. Everything in
these six files exists to move a visitor from the left box to the bottom box.}
\end{figure}

\subsection{The split between this repository and the live site}

The final commit, \code{0dd839e} on 2026-07-29, moved every file down into a
folder called \file{new-design-reference/} and renamed the publishing workflow so
it could no longer run. Its message says why:

\begin{quote}
\emph{``The live site is served from the portfolio repo, not this one. Parking
this design for reference and turning the workflow off so it cannot overwrite the
real site.''}
\end{quote}

Checked against the live address on 2026-08-11, while this document was being
written:

\begin{table}[H]\centering\small
\begin{tabular}{@{}ll@{}}
\toprule
What was checked & What came back \\
\midrule
\code{https://salmanadnan.com/} & \code{200}, the site is up \\
\code{http://salmanadnan.com/} & \code{308}, redirects to the secure version \\
Address it resolves to & \code{46.202.194.85}, the VPS named in the workflow \\
Title served & \emph{Salman Adnan} then \emph{AI \& Full-Stack Engineer} \\
Title in this repository & \emph{Salman Adnan: AI Systems for Entrepreneurs} \\
Markers unique to this repository & none present in the live page \\
Booking link on the live page & \code{cal.com/salman-adnan/meeting}, the same one \\
\bottomrule
\end{tabular}
\caption{Verified by fetching the live address on 2026-08-11. The live page is a
different build that happens to share this project's booking link.}
\end{table}

\begin{keyidea}{This repository holds the answer to a bug in its sibling}
The \code{ai-digitalise} Deep Dive records a finding it could not fully explain:
the page live at ai.digitalise.agency declares its canonical URL to be
\code{https://salmanadnan.com/}, which tells search engines to index this domain
instead of that one.

Fetching both addresses on 2026-08-11 completes the picture. The live
salmanadnan.com declares its own canonical correctly and carries a proper
\code{og:image}; the live ai.digitalise.agency declares salmanadnan.com's. Both
are served from the same VPS. So the agency page was produced from this personal
site's build and its canonical link was never changed.

Neither repository can fix that, because neither contains the live build. It is
recorded in both documents so that whoever maintains the portfolio repository
finds it.
\end{keyidea}

\begin{jargon}{Canonical URL}
A line in a page's head saying ``the authoritative home of this content is at
this address''. Search engines use it to avoid indexing the same content twice.
Pointing it at the wrong domain hands your search ranking to that domain.
\end{jargon}

\subsection{Reading order}

Files are explained in the order that makes each understandable when you reach
it, not alphabetically:

\begin{enumerate}
\item \file{new-design-reference/index.html}, the page itself, which every other
  file serves.
\item \file{new-design-reference/css/site.css}, which the page names and cannot
  be read without.
\item \file{new-design-reference/js/main.js}, which animates what the first two
  have already drawn.
\item \file{new-design-reference/netlify.toml}, a setting for a host never used.
\item \file{new-design-reference/assets/ASSETS.md}, a checklist never completed.
\item \file{.github/workflows-disabled/deploy.yml.disabled}, about publishing,
  which only makes sense once you know what is published.
\end{enumerate}

\subsection{A note on two characters}

Two characters in this repository's sources cannot be printed by this document's
typeface. They are printed as the closest stand-in, listed so no quoted line
misleads you:

\begin{table}[H]\centering\small
\begin{tabular}{@{}lll@{}}
\toprule
In the real file & Printed here as & Where it appears \\
\midrule
A long horizontal rule (U+2500) & a hyphen & comment banners in the CSS and JavaScript \\
A double horizontal rule (U+2550) & an equals sign & comment banners in the HTML \\
\bottomrule
\end{tabular}
\caption{Character substitutions used in every quoted listing in this document.}
\end{table}

The middle dot (U+00B7) prints exactly and is not substituted.

% ===========================================================================
\section{The ten-minute mental model}
% ===========================================================================

A web page is made of three kinds of file, and this project has one of each.

\begin{figure}[H]\centering
\begin{tikzpicture}[node distance=11mm and 14mm]
  \node[box]                 (html) {\file{index.html}\\\scriptsize 395 lines: the words and structure};
  \node[box, right=of html]  (css)  {\file{css/site.css}\\\scriptsize 575 lines: what it looks like};
  \node[box, right=of css]   (js)   {\file{js/main.js}\\\scriptsize 64 lines: what moves};
  \node[extbox, below=of css](br)   {The visitor's browser\\\scriptsize assembles all three};
  \draw[flow] (html) -- node[lbl]{names} (css);
  \draw[flow] (css)  -- node[lbl]{then} (js);
  \draw[flow] (html) |- (br);
  \draw[flow] (css)  -- (br);
  \draw[flow] (js)   |- (br);
\end{tikzpicture}
\caption{Nothing runs on a server. Everything happens on the visitor's machine.}
\end{figure}

\begin{jargon}{HTML}
HyperText Markup Language: the file holding the page's actual content, wrapped in
labels called tags that say what each piece of content \emph{is}. A tag looks like
\code{<h1>} to open and \code{</h1>} to close.
\end{jargon}

\begin{jargon}{CSS}
Cascading Style Sheets: the file deciding appearance. It names a kind of element
(``every heading'', ``anything labelled \code{btn}'') and lists the rules that
apply. ``Cascading'' means there is a defined pecking order when two rules
disagree.
\end{jargon}

\begin{jargon}{JavaScript}
A programming language running inside the browser after the page loads. It is
what makes a page react to a scroll or a click.
\end{jargon}

\subsection{The one paragraph that matters}

\file{index.html} stacks six bands: a fixed navigation bar, a two-column hero with
a large photo on the left and the name, tagline and three small statistics on the
right, a grid of six client projects, a band of four large counting numbers, two
case-study cards, a booking panel and a footer. \file{site.css} gives all of it a
near-black background with an orange accent, and rearranges it into one column
below 960 pixels. \file{main.js} does three things: the navigation bar gains a
frosted background past 40 pixels of scroll, the four large numbers count up from
zero when they scroll into view, and most blocks fade upward as they arrive. Every
``Book a call'' button is a link: two jump down the page, and the one in the
booking panel opens a cal.com calendar in a new tab.

\begin{keyidea}{The hero photo is a placeholder, and it says so on screen}
The left half of the hero, which is half the visible page when it loads, contains
no photograph. It contains a large letter \code{S} over an orange gradient and the
words \emph{``Drop Salman's photo here / assets/salman.jpg''}.

That text is not a code comment. It is content, styled and positioned, and any
visitor would read it. This is the single most consequential fact about the
design's state, and it is examined in section 7.1.
\end{keyidea}

% ===========================================================================
\section{Background a beginner needs}
% ===========================================================================

\subsection{What happens when someone opens the page}

\begin{figure}[H]\centering
\begin{tikzpicture}[node distance=26mm]
  \node[box] (b) {Browser};
  \node[box, right=of b] (s) {Web server\\\scriptsize (Caddy on the VPS)};
  \foreach \n in {b,s} \draw[dashed, linegrey] (\n.south) -- ++(0,-4.0);
  \draw[flow]  ($(b.south)+(0,-0.6)$) -- node[lbl,above]{give me \file{/}} ($(s.south)+(0,-0.6)$);
  \draw[flow]  ($(s.south)+(0,-1.3)$) -- node[lbl,above]{\file{index.html}} ($(b.south)+(0,-1.3)$);
  \draw[flow]  ($(b.south)+(0,-2.0)$) -- node[lbl,above]{give me \file{/css/site.css}} ($(s.south)+(0,-2.0)$);
  \draw[flow]  ($(s.south)+(0,-2.7)$) -- node[lbl,above]{the stylesheet} ($(b.south)+(0,-2.7)$);
  \draw[flow]  ($(b.south)+(0,-3.4)$) -- node[lbl,above]{give me \file{/js/main.js}} ($(s.south)+(0,-3.4)$);
\end{tikzpicture}
\caption{The browser reads \file{index.html} first and discovers the other two
files named inside it. The \code{<link>} and \code{<script>} tags are the only
reason it knows they exist.}
\end{figure}

\begin{jargon}{Web server}
A program on an always-on computer that waits for requests for files and sends
them back. This project is served by Caddy on the agency's rented machine, which
hands over files exactly as they sit on disk.
\end{jargon}

\begin{jargon}{VPS}
Virtual Private Server: a computer rented from a hosting company and controlled
remotely. The workflow names its address as \code{46.202.194.85}, which is the
same address salmanadnan.com resolves to today.
\end{jargon}

\subsection{Tags, attributes, classes and selectors}

A \textbf{tag} marks a piece of content: \code{<p>} a paragraph, \code{<section>}
a band of the page, \code{<a>} a link. Most come in pairs with content between.

An \textbf{attribute} is extra information inside the opening tag, written as a
name and a value: in \code{<a href="\#book">}, \code{href} has the value
\code{\#book}, meaning ``the thing named book on this same page''.

A \textbf{class} is a label you invent and attach with the \code{class}
attribute. It means nothing to the browser; it exists so CSS can find those
elements.

A \textbf{selector} is how CSS names what it wants: \code{.btn} means ``every
element labelled \code{btn}''.

\begin{jargon}{Element}
One thing on the page, created by a tag: a heading, a button, a box.
\end{jargon}

\subsection{Custom properties: values named once}

The stylesheet opens with a block giving names to values: \code{--accent:
\#ff6719;} means ``from now on \code{--accent} stands for this orange''. The rest
of the file writes \code{var(--accent)}.

\begin{jargon}{Hex colour}
A colour written as a hash and six characters, three pairs giving red, green and
blue from \code{00} to \code{ff}. \code{\#ff6719} is a strong orange: full red, a
little green, almost no blue.
\end{jargon}

\begin{jargon}{Custom property, also called a CSS variable}
A named value defined in CSS with two leading hyphens and read back with
\code{var(...)}. This project defines nineteen and reads fifteen of them.
\end{jargon}

% ===========================================================================
\section{Repository layout}
% ===========================================================================

\begin{figure}[H]\centering
\begin{forest}dirtree
[Salman Adnan Site
  [.github
    [workflows-disabled
      [deploy.yml.disabled]
    ]
  ]
  [new-design-reference
    [assets
      [ASSETS.md]
    ]
    [css
      [site.css]
    ]
    [js
      [main.js]
    ]
    [index.html]
    [netlify.toml]
  ]
]
\end{forest}
\caption{The complete repository. Nothing is folded away: these six files are
everything git tracks. Appendix A carries the same list with line counts.}
\end{figure}

Two things about this shape are deliberate and both are messages rather than
mechanics.

First, \textbf{every file lives under \file{new-design-reference/}}. Normally
\file{index.html} sits at the top so a web server pointed at the folder finds a
home page. Here it is one level down, so a server pointed at this folder finds
nothing. The folder name is a note to the next person: this is a reference, not a
site.

Second, \textbf{the automation folder is misspelled on purpose}. GitHub only runs
files in a folder named exactly \file{.github/workflows}. This one is
\file{.github/workflows-disabled}, and the file inside ends \file{.disabled}
rather than \file{.yml}. Either change alone would stop it. Both were made.

\begin{jargon}{Hidden folder}
A folder whose name begins with a dot, hidden from ordinary listings on Mac and
Linux. The dot is a convention meaning ``tooling lives here, not content''.
\end{jargon}

% ===========================================================================
\section{Architecture}
% ===========================================================================

There is no architecture in the usual sense, because nothing runs. What exists is
a set of one-way relationships between six files.

\begin{figure}[H]\centering
\begin{tikzpicture}[node distance=12mm and 20mm]
  \node[box]                      (html) {\file{index.html}\\\scriptsize 395 lines};
  \node[box, above right=of html] (css)  {\file{css/site.css}\\\scriptsize 575 lines};
  \node[box, below right=of html] (js)   {\file{js/main.js}\\\scriptsize 64 lines};
  \node[extbox, left=of html]     (fonts){Google Fonts\\\scriptsize two typefaces};
  \node[extbox, below=of html]    (cal)  {cal.com\\\scriptsize \code{salman-adnan/meeting}};
  \node[box, above=of css]        (as)   {\file{assets/ASSETS.md}\\\scriptsize a checklist, not code};
  \draw[flow]  (html) -- node[lbl]{\code{<link>}} (css);
  \draw[flow]  (html) -- node[lbl]{\code{<script>}} (js);
  \draw[dflow] (html) -- node[lbl]{\code{<link>}} (fonts);
  \draw[dflow] (html) -- node[lbl]{\code{href}} (cal);
  \draw[flow, dashed] (as) -- node[lbl]{describes} (css);
\end{tikzpicture}
\caption{What names what. Dashed boxes are services on other people's computers.
\file{netlify.toml} and the workflow are absent because nothing on the page refers
to them: they instruct the machines that publish the page, not the page.}
\end{figure}

The arrows go one way only. \file{index.html} knows the names of the other two
files; neither knows anything about it. The stylesheet styles class names it hopes
exist, and the JavaScript looks for elements that may not be there, with a guard
at every lookup so it does nothing quietly when they are missing.

\begin{honest}{Silent failure is the price of this simplicity}
Nothing checks that the class names in \file{index.html} and the selectors in
\file{site.css} agree. Section 7.2 documents a case where they do not: a pair of
stylesheet rules for the reversed case-study card target a class the HTML never
uses, so they have never once applied. The page looks right anyway, because the
HTML achieves the same layout by a different means. Nothing anywhere reported it.
\end{honest}

% ===========================================================================
\section{Data flow, end to end}
% ===========================================================================

\begin{figure}[H]\centering
\begin{tikzpicture}[node distance=8mm and 12mm]
  \node[box]             (a) {Browser downloads\\\file{index.html}};
  \node[box, right=of a] (b) {Reads \code{<head>},\\requests fonts and CSS};
  \node[box, right=of b] (c) {Draws the page\\top to bottom};
  \node[box, below=of c] (d) {Reaches \code{<script>}\\at the very end};
  \node[box, left=of d]  (e) {\file{main.js} runs,\\attaches watchers};
  \node[box, left=of e]  (f) {Visitor scrolls,\\watchers fire};
  \draw[flow] (a) -- (b); \draw[flow] (b) -- (c); \draw[flow] (c) -- (d);
  \draw[flow] (d) -- (e); \draw[flow] (e) -- (f);
\end{tikzpicture}
\caption{The full life of one page view. The JavaScript runs last because its
\code{<script>} tag is the last thing in the document, which is deliberate and is
explained in section 7.3.}
\end{figure}

There is no data flow beyond this. No form, nothing saved, nothing sent. The only
outbound step is following the cal.com link, at which point the visitor has left.

\begin{honest}{Nothing measures whether the page works}
No analytics of any kind: no counter, no tag manager, no event when somebody
presses ``Book a call''. The question this page exists to answer cannot be
answered from the page. Notably, the live salmanadnan.com \emph{does} carry a
full social-preview setup that this version lacks, so whoever built the live
version thought about distribution in a way this design did not.
\end{honest}

% ===========================================================================
\section{The file-by-file walkthrough}
% ===========================================================================

Every file git tracks gets its own subsection below, quoted from its first line to
its last. Nothing is summarized away and nothing is skipped.

% ===========================================================================
\subsection{\file{new-design-reference/index.html}}
% ===========================================================================

\textbf{HTML, 395 lines.} The page itself: every word a visitor reads and the
structure holding those words. Delete every other file and this one still opens in
a browser and still says everything the project wants to say. It would simply look
plain.

\begin{provenance}{new-design-reference/index.html}
Hand-written. It arrived complete in the first commit, \code{1f7ed12} on
2026-07-28, at the path \file{index.html}, and moved to its present path in
\code{0dd839e} on 2026-07-29 with no change to its contents. No generator banner,
no framework scaffold, and no \file{package.json} that could have produced it.

Its icons are a different matter and worth naming precisely. Every icon is an
inline SVG whose path data matches the Heroicons outline set, a widely used free
library: the shapes, the \code{24 24} view box, the \code{stroke-linecap} and
\code{stroke-linejoin} attributes and the \code{1.5}-based coordinates are that
library's house conventions. \textbf{The file carries no attribution comment}, so
this is an inference from the path data, not a documented fact.

Its skeleton is shared with the sibling repository \code{ai-digitalise}: identical
head block, identical navigation shape, identical metrics band, identical booking
panel and footer. Which was written first is not recorded in either repository and
is not guessed here; \code{ai-digitalise} was committed a few hours earlier on the
same day. \textbf{Confidence: certain for hand-written, inferred for the icon
origin, and the sibling relationship is stated as a resemblance with a shared
skeleton rather than as a direction of copying.}
\end{provenance}

\subsubsection*{Why it exists}

Without this file there is no page. The stylesheet would style class names
belonging to nothing, and the JavaScript would search for elements that do not
exist and quietly do nothing, which is exactly how it is written to behave.

\subsubsection*{Line by line}

\paragraph{Lines 1 to 17: the head, which the visitor never sees.}

\begin{lstlisting}[language=HTML,firstnumber=1,caption={\file{index.html}, lines 1--17}]
<!DOCTYPE html>
<html lang="en">
<head>
<meta charset="UTF-8" />
<meta name="viewport" content="width=device-width, initial-scale=1.0" />
<title>Salman Adnan: AI Systems for Entrepreneurs</title>
<meta name="description" content="I build AI systems and digital products that grow businesses without adding headcount. 17 projects, 1.6M+ organic views." />
<meta property="og:title" content="Salman Adnan" />
<meta property="og:description" content="AI systems that grow your business without adding headcount." />
<meta property="og:type" content="website" />
<meta name="theme-color" content="#05070d" />

<link rel="preconnect" href="https://fonts.googleapis.com" />
<link rel="preconnect" href="https://fonts.gstatic.com" crossorigin />
<link href="https://fonts.googleapis.com/css2?family=Space+Grotesk:wght@500;600;700&family=Inter:wght@400;500;600;700&display=swap" rel="stylesheet" />
<link rel="stylesheet" href="/css/site.css" />
</head>
\end{lstlisting}

Line 1 is not a tag. \code{<!DOCTYPE html>} declares ``interpret this as modern
HTML''. Without it, browsers fall back to a compatibility mode imitating software
from the 1990s. It must come first and has no closing partner.

Line 2 opens the document and declares its language English. That attribute does
real work: screen readers use it to choose a pronunciation and search engines use
it to decide who to show the page to.

Line 3 opens the \code{<head>}. Everything to line 17 is information \emph{about}
the page rather than content \emph{of} it, and none of it appears on screen.

Line 4 declares the character encoding UTF-8, the scheme letting a file contain any
character in any language. It matters concretely here because the page contains a
middle dot, which would otherwise render as nonsense. It should come as early as
possible, since the browser guesses until told.

\begin{jargon}{Meta tag}
A tag inside \code{<head>} carrying a fact about the page rather than content for
it. It has no closing partner and displays nothing. Different meta tags are read by
completely different audiences: browsers, search engines and social networks each
look for their own.
\end{jargon}

Line 5 is the viewport declaration, and it is the single line making the page work
on a phone. \code{width=device-width} tells the browser to treat the page width as
the real device width, and \code{initial-scale=1.0} says do not zoom out to begin
with. Without it a phone pretends to be about 980 pixels wide and shrinks
everything, which is why old sites look like tiny unreadable newspapers. Every
media query in the stylesheet depends on this line.

Line 6 is the title: the browser tab text and the blue headline in a search result.
Note it, because section 1.2 used it to prove this page is not the live one.

Line 7 is the description search engines often show as grey text under the title.
It is worth reading closely, because it makes two specific claims (``17 projects,
1.6M+ organic views'') that also appear on the page itself in two other places.

Lines 8 to 10 are Open Graph tags, a convention invented by Facebook and now read
by nearly every chat app to build the preview card when somebody pastes a link.

\begin{honest}{The Open Graph set is missing the tag that matters most}
Present: \code{og:title}, \code{og:description}, \code{og:type}. Missing:
\code{og:image} and \code{og:url}. Without an image, a pasted link produces a
preview card with no picture, which is the difference between a link that gets
clicked and one that does not.

This is more pointed here than in the sibling repository. The live salmanadnan.com
carries a complete set including \code{og:image} pointing at
\code{/assets/og.png}, verified on 2026-08-11. So whoever built the live version
supplied exactly the asset this one lists as required and never produced. See
section 7.5.
\end{honest}

Line 11 sets \code{theme-color} to \code{\#05070d}, the same near-black as the page
background, so a mobile browser paints its address bar to match and the browser
chrome blends into the page instead of framing it in white.

\begin{honest}{The theme colour was not rebranded}
\code{\#05070d} is the shared near-black, which is correct. But note it is
identical to the sibling's, as is the entire head block apart from the title and
description. The rebrand described in commit \code{05210af} touched the stylesheet
and left this file's shared parts alone, which is consistent with the five missed
colour literals documented in section 7.2.
\end{honest}

Lines 13 and 14 are \code{preconnect} hints, telling the browser it will shortly
need these two Google servers so it should start the handshake now. Establishing a
secure connection takes several network round trips, and doing that early typically
saves a few hundred milliseconds. The \code{crossorigin} attribute on line 14 is
required: font files are fetched under stricter rules than stylesheets, and a
preconnect made without it opens a connection the font request cannot reuse. That
subtlety catches out experienced developers constantly.

Line 15 fetches the fonts: Space Grotesk in three weights and Inter in four.
\code{display=swap} at the end is the important part. It tells the browser to render
text immediately in a fallback typeface and swap in the real one when it arrives,
rather than hiding text for up to three seconds. Text reflows slightly; the benefit
is that a visitor on a slow connection can read the page at all.

\begin{jargon}{Font weight}
How heavy a typeface looks, on a numeric scale where 400 is normal and 700 is bold.
Requesting \code{wght@500;600;700} downloads three separate files, which is why the
list is kept short.
\end{jargon}

Line 16 links the stylesheet, with the leading slash discussed in section 11: correct
when served by a web server, wrong when the file is opened off a disk.

Line 17 closes the head.

\paragraph{Lines 18 to 28: the navigation bar.}

\begin{lstlisting}[language=HTML,firstnumber=18,caption={\file{index.html}, lines 18--28}]
<body>

<!-- ══════════════════════════════════════════════
     NAV
══════════════════════════════════════════════ -->
<nav class="nav" aria-label="Site navigation">
  <div class="nav-inner">
    <span class="nav-logo">Salman Adnan</span>
    <a href="#book" class="btn btn-primary">Book a call</a>
  </div>
</nav>
\end{lstlisting}

Line 18 opens \code{<body>}. Everything to line 394 is what a visitor sees.

Lines 20 to 22 are a comment banner. Anything between \code{<!-{}-} and \code{-{}->} is
ignored entirely by the browser. This file uses six of these as signposts, one per
band, each drawn with a rule of repeated double-line characters. They are the
author's table of contents and the reason a 395-line file stays navigable.

Line 23 opens the navigation. \code{<nav>} is a \emph{semantic} tag: it looks no
different from a plain \code{<div>} but announces to screen readers and search
engines that this is the site's navigation. \code{aria-label} gives that region a
spoken name.

\begin{jargon}{ARIA}
Accessible Rich Internet Applications: attributes beginning \code{aria-} that
describe a page to assistive software. They change nothing visually.
\code{aria-label} supplies a name; \code{aria-hidden="true"} tells assistive
software to skip something because it is purely decorative.
\end{jargon}

Line 24 opens a \code{<div>} labelled \code{nav-inner}. A \code{<div>} is a generic
container with no meaning, existing purely to be styled. The reason for a second
container is a standard layout pattern: the outer element spans the full window so
its background stretches edge to edge, while the inner one is capped and centred so
the logo and button line up with the content below.

\begin{honest}{The logo is not a link}
Line 25 is a \code{<span>} containing the words ``Salman Adnan''. It is not
clickable, so there is no way back to the top of the page by clicking it, which is
the behaviour every visitor expects from a site logo.

On a single-page site this is defensible, and it differs from the sibling
repository, where the logo is an \code{<a>} pointing at \code{/} with a drawn mark
beside it. That divergence is one of the few places where the two skeletons
genuinely part company, and this is the weaker of the two.
\end{honest}

Line 26 is the navigation's only button, a link to \code{\#book}. It is a link and
not a \code{<button>}, which is correct: it navigates rather than performing an
action. It carries two classes, \code{btn} for the shared shape and
\code{btn-primary} for the filled variant. That two-class pattern repeats
throughout and is explained in section 7.2.

Lines 27 and 28 close the containers.

\paragraph{Lines 29 to 46: the hero opens, and the placeholder that defines this
project's state.}

\begin{lstlisting}[language=HTML,firstnumber=29,caption={\file{index.html}, lines 29--46}]


<!-- ══════════════════════════════════════════════
     HERO
     UPDATE: replace .hero-photo-placeholder with:
       <img src="assets/salman.jpg" alt="Salman Adnan" style="width:100%;height:100%;object-fit:cover;object-position:top;">
══════════════════════════════════════════════ -->
<section class="hero" aria-label="Hero">

  <!-- photo side -->
  <div class="hero-photo-side">
    <div class="hero-photo-frame">

      <!-- PLACEHOLDER: swap for Salman's actual photo (see comment above) -->
      <div class="hero-photo-placeholder" aria-hidden="true">
        <div class="placeholder-initial">S</div>
        <div class="placeholder-note">Drop Salman's photo here<br>assets/salman.jpg</div>
      </div>
\end{lstlisting}

Lines 31 to 35 are the second banner, and it carries an instruction rather than
just a label: replace the placeholder with an \code{<img>} tag, given in full,
including the exact styling needed (\code{object-fit: cover} to fill the frame
without distortion and \code{object-position: top} to keep a face in shot when the
frame is cropped). Whoever wrote this knew precisely what was needed and did not
have the file.

Line 36 opens the hero. Line 39 opens the left column and line 40 the frame that
holds the photograph and the badges floating over it.

Lines 42 to 46 are the most consequential nine lines in the repository.

\begin{honest}{Half the opening screen is a placeholder addressed to the developer}
The left column of the hero, which on a desktop is half of everything visible when
the page loads, contains no photograph. It contains a giant letter \code{S} in an
orange gradient and, beneath it, the words \emph{``Drop Salman's photo here''}
followed by \emph{``assets/salman.jpg''}.

That is not a code comment. It is content: it sits inside the page, is styled by
\file{site.css} lines 173 to 179, and renders on screen. The \code{aria-hidden} on
line 43 hides it from screen readers, which is correct for a decorative gradient
and means a screen-reader user is spared the instruction, while every sighted
visitor reads it.

For a personal brand page, whose entire proposition is a person, the missing
element is the person. This single fact is why the honest description of this
project is ``unfinished'' rather than ``finished and parked''. The sibling
\code{ai-digitalise} has no equivalent: its placeholders are gradients that merely
look like unfinished thumbnails, not text asking for a file.
\end{honest}

Note that the note names the path \code{assets/salman.jpg} while the instruction in
the banner comment above names \code{assets/salman.jpg} and \file{ASSETS.md} names
\code{/assets/salman.jpg} with a leading slash. The leading slash is the one that
would actually work, for the same reason as the stylesheet path.

\paragraph{Lines 47 to 59: the first floating badge.}

\begin{lstlisting}[language=HTML,firstnumber=47,caption={\file{index.html}, lines 47--59}]

      <!-- floating result badges over the photo -->
      <div class="hero-float-badge" aria-hidden="true">
        <div class="badge-icon badge-icon-blue">
          <svg viewBox="0 0 24 24" fill="none" stroke="#3f82ff" stroke-width="2" aria-hidden="true">
            <path stroke-linecap="round" stroke-linejoin="round" d="M3 13.125C3 12.504 3.504 12 4.125 12h2.25c.621 0 1.125.504 1.125 1.125v6.75C7.5 20.496 6.996 21 6.375 21h-2.25A1.125 1.125 0 013 19.875v-6.75zM9.75 8.625c0-.621.504-1.125 1.125-1.125h2.25c.621 0 1.125.504 1.125 1.125v11.25c0 .621-.504 1.125-1.125 1.125h-2.25a1.125 1.125 0 01-1.125-1.125V8.625zM16.5 4.125c0-.621.504-1.125 1.125-1.125h2.25C20.496 3 21 3.504 21 4.125v15.75c0 .621-.504 1.125-1.125 1.125h-2.25a1.125 1.125 0 01-1.125-1.125V4.125z"/>
          </svg>
        </div>
        <div>
          <div class="badge-text-main">17 projects</div>
          <div class="badge-text-sub">delivered</div>
        </div>
      </div>
\end{lstlisting}

Three small frosted cards drift slowly over the photograph. Each has an icon in a
tinted square and two lines of text: a number and a caption.

Line 49 is the container, marked \code{aria-hidden} because the same claims appear
as real text further down the page. That is the correct call: repeating them to a
screen reader would be duplication.

Lines 51 to 53 draw a bar-chart icon. Note \code{stroke="\#3f82ff"}, a literal
blue, rather than a token.

\begin{honest}{The badge icons are the one place blue is deliberate}
Three badges use three different colours: blue here, orange on line 63, green on
line 75, each matching a tinted background defined at \file{site.css} lines 230 to
232. Used as a set of three distinct signal colours this is a reasonable design
choice, and it is the only appearance of blue in this file that is plainly
intentional rather than left over from the sibling.

It is still written as three literals where two tokens exist (\code{-{}-accent2} is
exactly this blue and is otherwise never used, and \code{-{}-green} is exactly that
green). Section 7.2 covers what that costs.
\end{honest}

Lines 55 to 58 hold the text: ``17 projects'' and ``delivered''. That number will
appear twice more before the page ends.

\paragraph{Lines 60 to 71: the second badge.}

\begin{lstlisting}[language=HTML,firstnumber=60,caption={\file{index.html}, lines 60--71}]

      <div class="hero-float-badge" aria-hidden="true">
        <div class="badge-icon badge-icon-orange">
          <svg viewBox="0 0 24 24" fill="none" stroke="#ff6719" stroke-width="2" aria-hidden="true">
            <path stroke-linecap="round" stroke-linejoin="round" d="M2.25 18L9 11.25l4.306 4.307a11.95 11.95 0 015.814-5.519l2.74-1.22m0 0l-5.94-2.28m5.94 2.28l-2.28 5.941"/>
          </svg>
        </div>
        <div>
          <div class="badge-text-main">1.6M+ views</div>
          <div class="badge-text-sub">organic, for clients</div>
        </div>
      </div>
\end{lstlisting}

Structurally identical, with an upward-trending chart icon in orange and the claim
``1.6M+ views / organic, for clients''. The orange \code{\#ff6719} here is the
brand accent written literally rather than as \code{var(-{}-accent)}.

\paragraph{Lines 72 to 86: the third badge, and the end of the photo column.}

\begin{lstlisting}[language=HTML,firstnumber=72,caption={\file{index.html}, lines 72--86}]

      <div class="hero-float-badge" aria-hidden="true">
        <div class="badge-icon badge-icon-green">
          <svg viewBox="0 0 24 24" fill="none" stroke="#34d399" stroke-width="2" aria-hidden="true">
            <path stroke-linecap="round" stroke-linejoin="round" d="M9 12.75L11.25 15 15 9.75M21 12a9 9 0 11-18 0 9 9 0 0118 0z"/>
          </svg>
        </div>
        <div>
          <div class="badge-text-main">300% avg growth</div>
          <div class="badge-text-sub">across client sites</div>
        </div>
      </div>

    </div>
  </div><!-- /photo side -->
\end{lstlisting}

The third badge uses the tick-in-a-circle icon and the claim ``300\% avg growth /
across client sites''. Its green is \code{\#34d399}, which is exactly the value of
the \code{-{}-green} token defined at \file{site.css} line 15 and written literally
here.

Line 85 closes the frame and line 86 closes the column, with a trailing comment
naming what was closed. That habit is worth copying: in a file with this much
nesting a bare \code{</div>} tells the reader nothing.

\begin{keyidea}{Why these three badges are numbered from 2 in the stylesheet}
The stylesheet positions the badges with \code{.hero-float-badge:nth-child(2)},
\code{(3)} and \code{(4)}, at \file{site.css} lines 217 to 219. They are children
2, 3 and 4 of the frame because the placeholder on line 43 is child 1.

This survives the intended fix: replace the placeholder with an \code{<img>} as
the comment instructs and the image becomes child 1, leaving the badges where they
are. That is either careful thinking or luck; either way the numbers do not need
to change when the photograph arrives, which is worth knowing before somebody
"corrects" them.
\end{keyidea}

\paragraph{Lines 87 to 112: the right column, down to the buttons.}

\begin{lstlisting}[language=HTML,firstnumber=87,caption={\file{index.html}, lines 87--112}]

  <!-- content side -->
  <div class="hero-content-side">
    <div class="hero-tag">
      <svg width="14" height="14" viewBox="0 0 24 24" fill="none" stroke="currentColor" stroke-width="2.5" aria-hidden="true">
        <path stroke-linecap="round" stroke-linejoin="round" d="M9.813 15.904L9 18.75l-.813-2.846a4.5 4.5 0 00-3.09-3.09L2.25 12l2.846-.813a4.5 4.5 0 003.09-3.09L9 5.25l.813 2.846a4.5 4.5 0 003.09 3.09L15.75 12l-2.846.813a4.5 4.5 0 00-3.09 3.09z"/>
      </svg>
      Digitalise Agency
    </div>

    <h1 class="hero-h1">
      Salman<br>
      <em>Adnan.</em>
    </h1>

    <p class="hero-descriptor">I build AI systems that grow<br>your business without adding headcount.</p>

    <div class="hero-actions">
      <a href="#book" class="btn btn-primary">
        Book a call
        <svg width="16" height="16" viewBox="0 0 24 24" fill="none" stroke="currentColor" stroke-width="2.5" aria-hidden="true">
          <path stroke-linecap="round" stroke-linejoin="round" d="M13.5 4.5L21 12m0 0l-7.5 7.5M21 12H3"/>
        </svg>
      </a>
      <a href="#work" class="btn btn-ghost">See the work</a>
    </div>
\end{lstlisting}

Line 89 opens the content column. Lines 90 to 95 are the small pill above the name,
containing a sparkle icon and the words ``Digitalise Agency'', which ties this
personal page back to the agency.

Lines 97 to 100 are the headline in an \code{<h1>}. A page should have exactly one,
and this page does: it is the document's title as far as a screen reader or a
search engine is concerned, and here it is simply the person's name.
\code{<br>} forces the line break so the surname lands on its own line.

\begin{honest}{\code{<em>} is used for colour, not for emphasis}
Line 99 wraps ``Adnan.'' in \code{<em>}, and \file{site.css} line 272 then removes
the italics that tag normally implies and colours the word orange instead.
\code{<em>} means emphasis and a screen reader may change its intonation for it.
Using it as a colour hook works and slightly misrepresents the content to assistive
software; a \code{<span>} with a class would have been the honest choice. The
sibling repository does exactly the same thing, so it is a habit rather than a
slip. The intent is plainly visual.
\end{honest}

Line 102 is the tagline, and it repeats the page description from line 7 almost
word for word. \code{<br>} again controls where it breaks.

Lines 104 to 112 hold the two buttons: the primary one to \code{\#book} containing
both text and an arrow icon, and the ghost one to \code{\#work}. Both are ordinary
links, and both are intercepted by \file{main.js} line 55 to scroll smoothly rather
than jump.

\paragraph{Lines 113 to 130: the three small statistics, and the end of the hero.}

\begin{lstlisting}[language=HTML,firstnumber=113,caption={\file{index.html}, lines 113--130}]

    <div class="hero-mini-stats">
      <div class="hero-mini-stat">
        <div class="hero-mini-num">17+</div>
        <div class="hero-mini-lbl">Projects delivered</div>
      </div>
      <div class="hero-mini-stat">
        <div class="hero-mini-num">1.6M+</div>
        <div class="hero-mini-lbl">Organic views</div>
      </div>
      <div class="hero-mini-stat">
        <div class="hero-mini-num">10K+</div>
        <div class="hero-mini-lbl">Users served</div>
      </div>
    </div>
  </div>

</section>
\end{lstlisting}

Three statistics under a hairline: 17+ projects, 1.6M+ organic views, 10K+ users
served.

\begin{honest}{The same three numbers exist twice on this page, in two different
formats}
These three are written as plain visible text. The metrics band at lines 258 to 279
states the same three (plus a fourth) as \code{data-target} attributes with a
visible \code{0}, to be counted up by JavaScript.

So updating a number means editing it in two places, in two different ways, and
nothing checks that they agree. \file{ASSETS.md} line 24 knows this and says so
explicitly, which means it was a recognised cost rather than an oversight.

There is an upside worth stating. Because the hero copies are plain text, they are
the only version a search engine or a visitor without JavaScript ever sees, so this
duplication accidentally rescues the page from the ``claims zero of everything''
failure that the sibling repository has in full. The sibling has no hero
statistics, so its numbers vanish entirely without JavaScript.
\end{honest}

Line 130 closes the hero section.

\paragraph{Lines 131 to 162: the work grid, and the first project card.}

\begin{lstlisting}[language=HTML,firstnumber=131,caption={\file{index.html}, lines 131--162}]


<!-- ══════════════════════════════════════════════
     WORK PORTFOLIO
     UPDATE: replace gradient .wt* classes with real thumbnails.
             Add poster="assets/thumb-N.jpg" to each tile img.
══════════════════════════════════════════════ -->
<section class="work-section" id="work" aria-label="Work portfolio">
  <div class="container">
    <div class="reveal" style="margin-bottom:0;">
      <p class="section-kicker">The work</p>
      <h2 class="section-title">What I've built<br>for clients.</h2>
    </div>

    <div class="work-grid stagger">

      <div class="work-card" style="--i:0">
        <div class="work-thumb wt1">
          <div class="work-badge">300% traffic</div>
          <div class="work-thumb-overlay">
            <div class="work-play">
              <svg viewBox="0 0 24 24" fill="white" aria-hidden="true">
                <path d="M5.25 5.653c0-.856.917-1.398 1.667-.986l11.54 6.348a1.125 1.125 0 010 1.971l-11.54 6.347a1.125 1.125 0 01-1.667-.985V5.653z"/>
              </svg>
            </div>
          </div>
        </div>
        <div class="work-info">
          <div class="work-client">Hi-Tec Bearings</div>
          <div class="work-title">5,000+ SKUs, 300% organic traffic growth, 45% faster page loads</div>
        </div>
      </div>
\end{lstlisting}

Lines 133 to 137 are the third banner, again carrying an instruction: replace the
gradient classes with real thumbnails.

Line 138 carries \code{id="work"}, which is what makes the hero's ``See the work''
button work. An \code{id} names one element, and a link to \code{\#work} scrolls to
whatever carries it.

\begin{jargon}{id, and how it differs from class}
An \code{id} names exactly one element and must be unique. A \code{class} labels
any number. Use an \code{id} when something needs to be linked to; use a
\code{class} when you want to style a group. This page has exactly two ids,
\code{work} and \code{book}, and both exist to be jumped to.
\end{jargon}

Lines 140 to 143 establish a pattern used by the remaining sections: a wrapper with
the class \code{reveal} holding a small uppercase kicker and a large heading. The
\code{reveal} class is what \file{main.js} line 50 watches for, fading the pair
upward when they scroll into view. The inline \code{margin-bottom:0} on line 140
overrides the 52-pixel margin the stylesheet gives \code{.section-title}, because
this section wants its grid closer than the others do.

Line 142 is an \code{<h2>}. Heading levels form the document's outline: the single
\code{<h1>} in the hero, then an \code{<h2>} per major section. Getting that order
right is how a screen-reader user navigates a long page, and this file gets it
right.

Line 145 gives the grid the class \code{stagger}, and each card carries
\code{style="-{}-i:0"} and so on. That custom property is multiplied by 90
milliseconds at \file{site.css} line 86 to delay each child's fade-in, so the cards
appear in sequence rather than together. Setting a custom property inline is the
cleanest way to pass a per-element number to CSS.

Each card is two halves. The top, \code{work-thumb}, is where a real thumbnail
would go and is currently a gradient from the class \code{wt1}. Over it sit a badge
and an overlay containing a play button, which the stylesheet keeps invisible until
the mouse is over the card. The bottom, \code{work-info}, holds the client name and
the result.

\begin{honest}{The play buttons promise videos that do not exist}
Every card shows a play button on hover and none plays anything: the card is not a
link and has no click handler.

This is milder than the same fault in the sibling repository, which also sets
\code{cursor: pointer} so the mouse changes to a hand. Here it does not, so the
promise is visual only. The cheapest honest fix in both is to remove the play
overlay until the videos exist.
\end{honest}

\paragraph{Lines 163 to 196: project cards two and three.}

\begin{lstlisting}[language=HTML,firstnumber=163,caption={\file{index.html}, lines 163--196}]

      <div class="work-card" style="--i:1">
        <div class="work-thumb wt2">
          <div class="work-badge">3x leads</div>
          <div class="work-thumb-overlay">
            <div class="work-play">
              <svg viewBox="0 0 24 24" fill="white" aria-hidden="true">
                <path d="M5.25 5.653c0-.856.917-1.398 1.667-.986l11.54 6.348a1.125 1.125 0 010 1.971l-11.54 6.347a1.125 1.125 0 01-1.667-.985V5.653z"/>
              </svg>
            </div>
          </div>
        </div>
        <div class="work-info">
          <div class="work-client">DigiMarvixx</div>
          <div class="work-title">3x WhatsApp leads per month, 60% mobile traffic capture</div>
        </div>
      </div>

      <div class="work-card" style="--i:2">
        <div class="work-thumb wt3">
          <div class="work-badge">1.2s load</div>
          <div class="work-thumb-overlay">
            <div class="work-play">
              <svg viewBox="0 0 24 24" fill="white" aria-hidden="true">
                <path d="M5.25 5.653c0-.856.917-1.398 1.667-.986l11.54 6.348a1.125 1.125 0 010 1.971l-11.54 6.347a1.125 1.125 0 01-1.667-.985V5.653z"/>
              </svg>
            </div>
          </div>
        </div>
        <div class="work-info">
          <div class="work-client">Synergy Event</div>
          <div class="work-title">800 participants, 1.2s load, 98/100 Lighthouse, 100% mobile score</div>
        </div>
      </div>
\end{lstlisting}

Structurally identical to the first, differing in the stagger index, the gradient
class, the badge text, the client and the result. The six clients across this grid
are Hi-Tec Bearings, DigiMarvixx, Synergy Event, IBA GPT, Prestige Sleep and IBA
CED, the same six as in the sibling repository, described here in slightly more
detail (this version adds ``45\% faster page loads'' and ``100\% mobile score'').

These are specific, named, checkable claims, which is a meaningfully different kind
of statement from the floating badges above them.

\paragraph{Lines 197 to 251: project cards four, five and six.}

\begin{lstlisting}[language=HTML,firstnumber=197,caption={\file{index.html}, lines 197--251}]

      <div class="work-card" style="--i:3">
        <div class="work-thumb wt4">
          <div class="work-badge">AI chat</div>
          <div class="work-thumb-overlay">
            <div class="work-play">
              <svg viewBox="0 0 24 24" fill="white" aria-hidden="true">
                <path d="M5.25 5.653c0-.856.917-1.398 1.667-.986l11.54 6.348a1.125 1.125 0 010 1.971l-11.54 6.347a1.125 1.125 0 01-1.667-.985V5.653z"/>
              </svg>
            </div>
          </div>
        </div>
        <div class="work-info">
          <div class="work-client">IBA GPT</div>
          <div class="work-title">AI university assistant answering student questions around the clock</div>
        </div>
      </div>

      <div class="work-card" style="--i:4">
        <div class="work-thumb wt5">
          <div class="work-badge">24/7 sales</div>
          <div class="work-thumb-overlay">
            <div class="work-play">
              <svg viewBox="0 0 24 24" fill="white" aria-hidden="true">
                <path d="M5.25 5.653c0-.856.917-1.398 1.667-.986l11.54 6.348a1.125 1.125 0 010 1.971l-11.54 6.347a1.125 1.125 0 01-1.667-.985V5.653z"/>
              </svg>
            </div>
          </div>
        </div>
        <div class="work-info">
          <div class="work-client">Prestige Sleep</div>
          <div class="work-title">Clinic booking and product shop combined, sells after hours</div>
        </div>
      </div>

      <div class="work-card" style="--i:5">
        <div class="work-thumb wt6">
          <div class="work-badge">4 dashboards</div>
          <div class="work-thumb-overlay">
            <div class="work-play">
              <svg viewBox="0 0 24 24" fill="white" aria-hidden="true">
                <path d="M5.25 5.653c0-.856.917-1.398 1.667-.986l11.54 6.348a1.125 1.125 0 010 1.971l-11.54 6.347a1.125 1.125 0 01-1.667-.985V5.653z"/>
              </svg>
            </div>
          </div>
        </div>
        <div class="work-info">
          <div class="work-client">IBA CED</div>
          <div class="work-title">Analytics platform with 4 executive dashboards tracking entrepreneurship outcomes</div>
        </div>
      </div>

    </div>
  </div>
</section>
\end{lstlisting}

The final three complete the set, with indices 3, 4 and 5 and gradients \code{wt4}
through \code{wt6}. Two of the six are for the same institution (IBA GPT and IBA
CED), the only repetition among them.

Note that the same play-triangle path data appears six times in this section and
nowhere is it shared. That is the cost of having no templating system, and it is
also the correct trade for a project that deliberately has no build step.

Line 249 closes the grid and lines 250 and 251 the container and section.

\paragraph{Lines 252 to 279: the metrics band.}

\begin{lstlisting}[language=HTML,firstnumber=252,caption={\file{index.html}, lines 252--279}]


<!-- ══════════════════════════════════════════════
     METRICS
     UPDATE: replace data-target values with real numbers
══════════════════════════════════════════════ -->
<section class="metrics-section" aria-label="Results">
  <div class="container">
    <div class="metrics-grid stagger">
      <div class="metric-item" style="--i:0">
        <div class="metric-num" data-target="17" data-suffix="+">0</div>
        <div class="metric-label">Projects Delivered</div>
      </div>
      <div class="metric-item" style="--i:1">
        <div class="metric-num" data-target="1.6" data-suffix="M+">0</div>
        <div class="metric-label">Organic Views for Clients</div>
      </div>
      <div class="metric-item" style="--i:2">
        <div class="metric-num" data-target="10" data-suffix="K+">0</div>
        <div class="metric-label">Users Served</div>
      </div>
      <div class="metric-item" style="--i:3">
        <div class="metric-num" data-target="300" data-suffix="%">0</div>
        <div class="metric-label">Avg. Traffic Growth</div>
      </div>
    </div>
  </div>
</section>
\end{lstlisting}

This band is where the JavaScript does its most visible work, and it introduces the
mechanism the counter depends on.

Look at line 262: the element's visible content is the single character \code{0},
and the real number lives in \code{data-target="17"}, with \code{data-suffix="+"}
beside it.

\begin{jargon}{Data attribute}
Any attribute whose name starts with \code{data-}. HTML ignores it completely and
displays nothing; it exists so JavaScript can attach its own information to an
element. \file{main.js} reads these two at lines 17 and 18 to know what to count
towards and what to print after it.
\end{jargon}

Why is the visible content \code{0}? Because that is what a visitor sees before the
counter runs, and it is also what they see forever if JavaScript is disabled or
fails. The robust alternative is to write the real number in the HTML and count
towards it from zero, which fails gracefully. This file chose the fragile form, and
is partly rescued by the hero statistics repeating three of the four as plain text.

This band still carries an \code{UPDATE} comment on line 256 telling the reader to
replace the values with real numbers, even though commit \code{692e162} claims to
have replaced all placeholder content with real project data.

\paragraph{Lines 280 to 318: the case studies, and the first card.}

\begin{lstlisting}[language=HTML,firstnumber=280,caption={\file{index.html}, lines 280--318}]


<!-- ══════════════════════════════════════════════
     CASE STUDIES (visual-only, numbers speak)
     UPDATE: replace gradient backgrounds with real screenshots.
             Use background-image on .case-visual or swap the div for <img>.
══════════════════════════════════════════════ -->
<section class="cases-section" aria-label="Case studies">
  <div class="container">
    <div class="reveal" style="margin-bottom:52px;">
      <p class="section-kicker">Case studies</p>
      <h2 class="section-title">Results in numbers.</h2>
    </div>

    <!-- Case 1 -->
    <div class="case-card reveal">
      <div class="case-visual cv1">
        <div class="case-visual-inner">
          <div class="case-visual-metrics">
            <div class="cvm">
              <span class="cvm-num">300%</span>
              <span class="cvm-lbl">Traffic growth</span>
            </div>
            <div class="cvm">
              <span class="cvm-num">5K+</span>
              <span class="cvm-lbl">Products</span>
            </div>
            <div class="cvm">
              <span class="cvm-num">99.9%</span>
              <span class="cvm-lbl">Uptime</span>
            </div>
          </div>
        </div>
      </div>
      <div class="case-content">
        <div class="case-client-name">Hi-Tec Bearings</div>
        <div class="case-title">Industrial e-commerce with 5,000+ SKUs that tripled organic traffic</div>
      </div>
    </div>
\end{lstlisting}

Two large cards, each a two-column arrangement of a visual panel and a block of
text. The visual is a gradient standing in for a screenshot, with three small
frosted chips of metrics laid over its bottom edge.

Line 295 gives the card the classes \code{case-card} and \code{reveal}, so it fades
in as a whole rather than being staggered.

Lines 299 to 310 are the three metric chips: 300\% traffic growth, 5K+ products,
99.9\% uptime. Note that 99.9\% uptime is a claim appearing nowhere else on the page
and in no other repository in this pair.

\paragraph{Lines 319 to 347: the second case study, and a dead class.}

\begin{lstlisting}[language=HTML,firstnumber=319,caption={\file{index.html}, lines 319--347}]

    <!-- Case 2 -->
    <div class="case-card case-card-reverse reveal">
      <div class="case-visual cv2" style="order:2">
        <div class="case-visual-inner">
          <div class="case-visual-metrics">
            <div class="cvm">
              <span class="cvm-num">3x</span>
              <span class="cvm-lbl">Leads/month</span>
            </div>
            <div class="cvm">
              <span class="cvm-num">60%</span>
              <span class="cvm-lbl">Mobile traffic</span>
            </div>
            <div class="cvm">
              <span class="cvm-num">WhatsApp</span>
              <span class="cvm-lbl">Integration</span>
            </div>
          </div>
        </div>
      </div>
      <div class="case-content" style="order:1">
        <div class="case-client-name">DigiMarvixx</div>
        <div class="case-title">Digital marketing site that tripled monthly leads through WhatsApp</div>
      </div>
    </div>

  </div>
</section>
\end{lstlisting}

The second card is meant to mirror the first: text on the left, visual on the
right. Line 321 gives it the extra class \code{case-card-reverse}, and lines 322
and 340 set \code{order:2} and \code{order:1} inline.

\begin{honest}{The reversal class matches nothing, and the layout works anyway}
\file{site.css} lines 436 and 437 define \code{.case-card.reverse}, which means ``an
element carrying both the class \code{case-card} and the class \code{reverse}''.
This element carries \code{case-card} and \code{case-card-reverse}. A hyphen does
not create a second class, so \textbf{those rules have never once applied}, and
neither has the third at line 559 in the responsive block.

The reversal happens regardless, because of the inline \code{order} properties on
lines 322 and 340, which tell the grid to place the visual second and the text
first. So the design is correct, three stylesheet rules are dead, and nothing
anywhere reported it.

This is the clearest example in the pair of what ``silent failure'' means in
practice: a whole mechanism can be wrong while the page looks right, because a
second mechanism quietly does the same job.
\end{honest}

\begin{jargon}{order}
A CSS property setting an item's position inside a flex or grid container
independently of where it sits in the HTML. \code{order:2} places an item after
one with \code{order:1}, regardless of which is written first.
\end{jargon}

Note also that the inline \code{style="order:2"} on line 322 sits alongside the
class \code{cv2}, so one element is being positioned by an inline style and
coloured by a class, which is the same mixed habit seen throughout the pair.

\paragraph{Lines 348 to 377: the booking panel.}

\begin{lstlisting}[language=HTML,firstnumber=348,caption={\file{index.html}, lines 348--377}]


<!-- ══════════════════════════════════════════════
     BOOK A CALL
     UPDATE: replace YOUR_CAL_LINK with your Cal.com or Calendly URL.
══════════════════════════════════════════════ -->
<section class="book-section" id="book" aria-label="Book a call">
  <div class="container">
    <div class="book-wrap reveal">
      <p class="book-eyebrow">Let's talk</p>
      <h2 class="book-title">30 minutes.<br>No slides.</h2>
      <p class="book-sub">Tell me about your content goals. I'll tell you if I can help.</p>
      <div class="book-cta-row">
        <!-- UPDATE: replace href with your booking link -->
        <a href="https://cal.com/salman-adnan/meeting" target="_blank" rel="noopener" class="btn btn-primary" style="padding:16px 36px;font-size:17px;">
          Book a call
          <svg width="18" height="18" viewBox="0 0 24 24" fill="none" stroke="currentColor" stroke-width="2.5" aria-hidden="true">
            <path stroke-linecap="round" stroke-linejoin="round" d="M13.5 4.5L21 12m0 0l-7.5 7.5M21 12H3"/>
          </svg>
        </a>
      </div>
      <p class="book-detail">
        <svg viewBox="0 0 24 24" fill="none" stroke="currentColor" stroke-width="2" aria-hidden="true">
          <path stroke-linecap="round" stroke-linejoin="round" d="M12 6v6h4.5m4.5 0a9 9 0 11-18 0 9 9 0 0118 0z"/>
        </svg>
        Free · <span>30 minutes · Zoom</span>
      </p>
    </div>
  </div>
</section>
\end{lstlisting}

This is the destination of every ``Book a call'' button, via \code{id="book"} on
line 354. It is the point of the whole document.

Line 356 combines two classes, \code{book-wrap} for the panel's appearance and
\code{reveal} so it fades in, the only place in the file where a layout class and a
behaviour class share an element.

Line 362 is the booking link, and three of its attributes matter.
\code{target="\_blank"} opens it in a new tab so the visitor does not lose the page.
\code{rel="noopener"} is a security measure that must accompany it: without it the
newly opened page receives a reference back to this one and can navigate it
somewhere else, a trick known as tab-nabbing. Its presence is correct and is the
kind of detail usually forgotten.

\begin{keyidea}{This booking link is real, and it outlived the design}
\code{cal.com/salman-adnan/meeting} is a genuine personal calendar, set
deliberately in commit \code{014406e}. The sibling repository's equivalent points
at \code{cal.com/davonex/30min}, another project's calendar, and is a known
outstanding placeholder.

More than that, this link is the one part of this design that is demonstrably still
in service: fetching the live salmanadnan.com on 2026-08-11 found the same link on
the page that replaced this one.
\end{keyidea}

\begin{honest}{The comment above the link is stale}
Line 352 still tells the reader to search for \code{YOUR\_CAL\_LINK}, and line 361
repeats the instruction. That literal placeholder was replaced on 2026-07-28 and
appears nowhere in the file. \file{ASSETS.md} line 28 repeats the same stale
instruction a third time. Somebody following the written instructions would search,
find nothing, and have no way to tell whether the job was done or the instruction
was wrong.
\end{honest}

Lines 369 to 374 are the reassurance line under the button: a clock icon, then
``Free'', then the duration and platform. The middle dots are the U+00B7 character
discussed in section 1.4 and render exactly.

\paragraph{Lines 378 to 395: the footer, and the script tag.}

\begin{lstlisting}[language=HTML,firstnumber=378,caption={\file{index.html}, lines 378--395}]


<!-- ══════════════════════════════════════════════
     FOOTER
══════════════════════════════════════════════ -->
<footer>
  <div class="container">
    <p>
      <a href="mailto:salman@digitalise.agency">salman@digitalise.agency</a>
      &nbsp;·&nbsp;
      <a href="https://digitalise.agency" target="_blank" rel="noopener">digitalise.agency</a>
    </p>
  </div>
</footer>

<script src="/js/main.js"></script>
</body>
</html>
\end{lstlisting}

The footer is one line: an email address and a link to the agency, separated by a
middle dot padded with \code{\&nbsp;} on each side.

\begin{jargon}{Non-breaking space}
\code{\&nbsp;} is a space a browser will not use as a place to wrap a line, and
which is not collapsed when several sit together. Here it does the second job:
ordinary spaces around the dot would collapse to one, and these keep the separator
visually padded.
\end{jargon}

Line 386 is a \code{mailto:} link, which opens the visitor's email program with the
address filled in. Line 388 links to the agency site, correctly carrying both
\code{target="\_blank"} and \code{rel="noopener"}.

This footer is better than the sibling's, which links to \code{digitalise.agency}
as text but points the link at \code{/}, its own home page. Here the text and the
destination agree.

Line 393 loads the JavaScript, and \textbf{where it sits is the whole point}. It is
the last thing before \code{</body>}, so by the time the browser reaches it every
element already exists. \file{main.js} begins searching for elements immediately,
with no wrapper telling it to wait, and that only works because of this placement.
Move this line into the \code{<head>} and the script would run against an empty
document, find nothing, and the page would quietly lose all three behaviours with
no error message.

Lines 394 and 395 close the body and the document.

\subsubsection*{What it connects to}

\textbf{This file names:} \file{new-design-reference/css/site.css} (line 16),
\file{new-design-reference/js/main.js} (line 393), Google Fonts (lines 13 to 15),
cal.com (line 362), and the agency site (line 388).

\textbf{This file is named by:} nothing inside the repository. The workflow copies
the whole folder rather than referring to files individually, and
\file{netlify.toml} publishes the folder as a whole. \file{ASSETS.md} discusses it
at length and cannot affect it.

\textbf{Internal links:} the navigation and hero both point at \code{\#book} (lines
26 and 105), and the hero also points at \code{\#work} (line 111). Those two ids,
on lines 354 and 138, are the only link targets in the file.

\subsubsection*{How it breaks}

\begin{itemize}
\item \textbf{Rename a class here without renaming it in the stylesheet} and that
  element loses its appearance silently. Nothing warns you. Section 7.2 documents a
  live instance of the reverse: a stylesheet class that matches nothing here.
\item \textbf{Move the \code{<script>} tag into the \code{<head>}} and all three
  behaviours stop, with no error. The metrics band would then permanently show four
  zeros.
\item \textbf{Insert an element before the placeholder on line 43} and all three
  floating badges jump, because the stylesheet positions them by child number.
\item \textbf{Open the file directly off a disk} and it renders unstyled, because of
  the leading slash on line 16. This looks like catastrophic breakage and is not.
\item \textbf{Disable JavaScript} and the page is fully readable but static. The
  metrics band shows four zeros; the three hero statistics still show their real
  values, which is what saves it.
\item \textbf{Delete the id on line 354} and every ``Book a call'' button silently
  does nothing, because \file{main.js} returns early when it cannot find a link's
  target.
\end{itemize}

\subsubsection*{Honest findings for this file}

The visible photo placeholder on lines 43 to 46; the missing \code{og:image} and
favicon; seven stale comments (lines 33, 42, 135, 256, 284, 352, 361); a
non-clickable logo on line 25; \code{<em>} used as a colour hook on line 99; the
same three numbers stated twice in two formats; play buttons that play nothing;
colour literals used where tokens exist (lines 51, 63, 75); and a
\code{case-card-reverse} class that no stylesheet rule matches.

% ===========================================================================
\subsection{\file{new-design-reference/css/site.css}}
% ===========================================================================

\textbf{CSS, 575 lines.} Every visual decision: colours, typefaces, spacing,
layout, movement, and how all of it rearranges on a narrow screen. It is the
longest file here and, in the judgement of this document, the most valuable thing
in the repository.

\begin{provenance}{new-design-reference/css/site.css}
Hand-written, and it states its own lineage in its second line: \code{Personal
brand site. Shares design tokens with ai.digitalise.agency.} That is a direct
statement by the author that the naming scheme came from the sibling project, and
comparing the two token blocks confirms it: identical names, identical structure,
different values.

It arrived in the first commit \code{1f7ed12} and was substantially edited once, by
\code{05210af} (``Differentiate to orange brand''). No generator banner, no vendor
prefix block of the sort a preprocessor emits, and no minified section: it is
written to be read. \textbf{Confidence: certain for hand-written, certain for the
token lineage, which rests on the author's own comment and is corroborated by the
files.}
\end{provenance}

\subsubsection*{Why it exists}

Without this file the page contains every word it needs and looks like a 1994
document: black serif text on white, one column, no colour. Nothing would be
broken and nothing would be persuasive. The two-column hero would collapse into a
stack, and the three floating badges would appear as three small blocks of text
under the placeholder.

\subsubsection*{Line by line}

\paragraph{Lines 1 to 25: the header, and the token block.}

\begin{lstlisting}[firstnumber=1,caption={\file{css/site.css}, lines 1--25}]
/* salmanadnan.com · site.css
   Personal brand site. Shares design tokens with ai.digitalise.agency. */

/* ─── TOKENS ─── */
:root {
  --s0:    #05070d;
  --s1:    #0e1118;
  --s2:    #161b26;
  --ink1:  #f4f7fc;
  --ink2:  #c3ccdd;
  --ink3:  #93a0b8;
  --accent:  #ff6719;   /* orange, personal brand, distinct from agency blue */
  --accentH: #ff8c4a;
  --accent2: #1967ff;
  --green:   #34d399;
  --border:  #24201e;
  --borderS: #1c1814;
  --font-d: 'Space Grotesk', ui-sans-serif, system-ui, sans-serif;
  --font-b: 'Inter', ui-sans-serif, system-ui, sans-serif;
  --r:   16px;
  --rsm: 10px;
  --ease: cubic-bezier(.22,1,.36,1);
  --maxw: 1240px;
  --shadow: 0 8px 40px rgba(0,0,0,.45);
}
\end{lstlisting}

Lines 1 and 2 are a comment. In CSS, anything between \code{/*} and \code{*/} is
ignored, and unlike HTML there is only this one comment form. These two lines carry
the file's entire provenance.

Line 4 is the first of twelve section banners drawn with box-drawing characters,
the author's navigation aid.

Line 5 opens the token block. \code{:root} matches the document's outermost
element, \code{<html>}. Custom properties defined there are inherited by every
element on the page, which is exactly why token blocks are written against
\code{:root}: it is the one place a definition reaches everything.

The nineteen tokens fall into six groups:

\begin{table}[H]\centering\small
\begin{tabular}{@{}llp{6.2cm}@{}}
\toprule
Tokens & Line & What the group is for \\
\midrule
\code{--s0} to \code{--s2}      & 6--8   & Surfaces, darkest to lightest. \code{s0} is the page, \code{s1} is a card, \code{s2} would be a detail inside a card \\
\code{--ink1} to \code{--ink3}  & 9--11  & Text, brightest to dimmest: a heading, body text, a caption \\
\code{--accent}, \code{--accentH}, \code{--accent2} & 12--14 & The orange, its hover variant, and the agency blue kept as a secondary \\
\code{--green}                  & 15     & The one success colour \\
\code{--border}, \code{--borderS} & 16--17 & A visible hairline, and a subtler one \\
\code{--font-d}, \code{--font-b} & 18--19 & Display typeface and body typeface \\
\code{--r}, \code{--rsm}, \code{--ease}, \code{--maxw}, \code{--shadow} & 20--24 & Corner radius, small radius, motion curve, content width, drop shadow \\
\bottomrule
\end{tabular}
\caption{The token vocabulary, identical in naming to the sibling repository and
different in almost every value.}
\end{table}

\begin{keyidea}{Why a numbered scale beats named colours}
Calling a colour \code{--s1} instead of \code{--dark-navy} names what it is
\emph{for} (the surface one step up from the page) rather than how it looks. This
pair of repositories proves the point better than any argument could: the two files
share every token \emph{name} and differ in nearly every token \emph{value}, so the
same stylesheet vocabulary produced a blue agency site and an orange personal one
with no renaming at all.
\end{keyidea}

Line 12 carries the rebrand, and it carries its own explanation in an inline
comment: \code{orange, personal brand, distinct from agency blue}. Line 14 then
keeps the agency blue as \code{--accent2}.

\begin{honest}{Four tokens are defined and never read}
\code{--s2} (line 8), \code{--accent2} (line 14), \code{--rsm} (line 21) and
\code{--shadow} (line 24) are read zero times in this file. Counted mechanically:
\code{grep -c 'var(--accent2)'} returns \code{0}, and likewise for the other three.

\code{--accent2} is the most interesting of the four. It holds the agency blue,
deliberately kept as a secondary colour, and is then never used as one. Meanwhile
five other rules in this file hard-code that same blue as a literal rather than
reading this token. The colour is used and the token for it is not, which is the
worst of both worlds and is the subject of section 7.2's largest finding, below.
\end{honest}

Line 22 defines the motion curve. \code{cubic-bezier(.22,1,.36,1)} describes how an
animation's speed varies across its run, by giving two control points of a curve.
This one starts fast and settles gently, the shape usually described as ``ease
out''. Every transition in the file refers to it, so all movement shares one
personality.

Line 23 sets the content width to 1240 pixels, 80 narrower than the sibling's 1320,
which is the kind of small deliberate difference that shows these files were tuned
separately rather than copied wholesale.

\paragraph{Lines 26 to 43: the reset, and the one utility class.}

\begin{lstlisting}[firstnumber=26,caption={\file{css/site.css}, lines 26--43}]

*, *::before, *::after { box-sizing: border-box; margin: 0; padding: 0; }
html { color-scheme: dark; scroll-behavior: smooth; }
body {
  font-family: var(--font-b);
  background: var(--s0);
  color: var(--ink1);
  -webkit-font-smoothing: antialiased;
  text-rendering: optimizeLegibility;
  overflow-x: hidden;
  line-height: 1.5;
}
img, video, svg { max-width: 100%; display: block; }
a { color: inherit; text-decoration: none; }
button { cursor: pointer; border: none; font-family: inherit; }

/* ─── UTILITY ─── */
.container { max-width: var(--maxw); margin: 0 auto; padding: 0 28px; }
\end{lstlisting}

Line 27 is a reset doing three jobs. The selector \code{*} means every element;
\code{*::before} and \code{*::after} extend that to the two invisible
pseudo-elements CSS lets you attach to anything.

\code{box-sizing: border-box} is the most consequential declaration in the file. By
default, an element given a width of 200 pixels and 20 pixels of padding occupies
240: the padding is added outside the stated width. That default surprises
everyone. \code{border-box} makes the stated width final, with padding eating into
it. Every layout calculation here assumes it.

\code{margin: 0; padding: 0} removes the browser's built-in spacing on headings and
paragraphs, so all spacing in the design is deliberate rather than inherited from
defaults that vary between browsers.

Line 28 does two things to the document root. \code{color-scheme: dark} tells the
browser this page is dark, so scrollbars, selection highlights and form controls
are drawn in their dark variants rather than staying stubbornly white.
\code{scroll-behavior: smooth} makes every jump-link glide rather than snap.

\begin{honest}{Smooth scrolling is implemented twice}
Line 28 already makes anchor links glide. \file{main.js} lines 55 to 62 then
intercept every anchor link and call \code{scrollIntoView(\{behavior: 'smooth'\})}
to do the same again. Either alone suffices. The JavaScript copy is the redundant
one and it also carries a latent bug, described in section 7.3.
\end{honest}

Lines 29 to 37 set the page defaults: body typeface, darkest surface as background,
brightest ink as text, and a line height of 1.5, meaning each line occupies one and
a half times the font size. That last number is the most important typographic
decision in the file and the one most often set badly.

\code{-webkit-font-smoothing: antialiased} on line 33 thins text rendering on Mac
displays. It is a Safari and Chrome specific property with no standard equivalent,
and it is the conventional companion to light text on a dark background, which
otherwise looks heavier than intended.

\code{overflow-x: hidden} on line 35 forbids horizontal scrolling, and here it is
load-bearing rather than cosmetic: the hero is a two-column grid whose left column
is a full-height photo frame, and the floating badges are positioned by percentage
inside it.

Line 38 makes images, videos and SVGs never exceed their container and makes them
block-level. That second part fixes a genuinely confusing default: images are
inline by default, which reserves space beneath them for the descenders of text
that is not there, producing a mysterious few pixels of gap under every image.

Line 39 strips links of their blue colour and underline, making them inherit the
surrounding colour. That is a design choice and an accessibility trade: colour
alone now distinguishes a link from text, which fails the usual guidance. Every
link here is either a button or a footer item, so the practical harm is small.

Line 40 styles buttons, and is dead in the strict sense that this page contains no
\code{<button>} element at all. It costs nothing and would come alive the moment one
was added.

Line 43 is the only utility class: a centred column no wider than 1240 pixels with
28 pixels of breathing room each side. \code{margin: 0 auto} is the standard
centring idiom, meaning no vertical margin and equal automatic margins left and
right. Every section below the hero wraps its content in this; the hero
deliberately does not, because it spans the full window.

\paragraph{Lines 44 to 72: the button system, and the first missed rebrand.}

\begin{lstlisting}[firstnumber=44,caption={\file{css/site.css}, lines 44--72}]

.btn {
  display: inline-flex;
  align-items: center;
  gap: 8px;
  padding: 13px 28px;
  border-radius: 100px;
  font-family: var(--font-d);
  font-size: 15px;
  font-weight: 600;
  letter-spacing: -0.01em;
  transition: transform 0.18s var(--ease), box-shadow 0.18s var(--ease),
              background 0.18s var(--ease);
}
.btn-primary {
  background: var(--accent);
  color: #fff;
}
.btn-primary:hover {
  background: var(--accentH);
  transform: translateY(-2px);
  box-shadow: 0 8px 24px rgba(25,103,255,.4);
}
.btn-ghost {
  background: transparent;
  color: var(--ink2);
  border: 1px solid var(--border);
}
.btn-ghost:hover { background: var(--s1); color: var(--ink1); }
\end{lstlisting}

This is the two-class pattern the HTML uses: \code{.btn} holds everything both
buttons share and \code{.btn-primary} or \code{.btn-ghost} adds only what differs.
Add a third style and you write four lines, not thirty.

\code{display: inline-flex} on line 46 makes the button a flex container sitting in
the flow of text rather than on its own line. That matters because the primary
buttons contain both a word and an icon, and flex is what lines those up.

\begin{jargon}{Flexbox}
A layout system arranging an element's children in a row or a column and
controlling how they align and share leftover space. Set \code{display: flex} on a
parent and its children become flex items. \code{align-items: center} lines them up
on their centres, which is how the icon in a button sits level with the text rather
than on the baseline.
\end{jargon}

\code{gap: 8px} on line 48 puts eight pixels between word and icon without either
needing a margin. Before \code{gap} existed this required fiddly rules targeting
all-but-the-first child, and its arrival is one of the genuine improvements in
modern CSS.

\code{border-radius: 100px} on line 50 makes a pill. Any radius larger than half the
element's height gives a perfect semicircular end, so 100 is simply comfortably
larger than needed.

Lines 55 and 56 declare the transition: three named properties, each over 0.18
seconds, each using the shared curve. Listing properties explicitly rather than
writing \code{transition: all} is better practice, because \code{all} makes the
browser watch every property and can animate things you never intended.

\begin{honest}{An orange button that casts a blue shadow}
Line 65 sets the hover shadow to \code{rgba(25,103,255,.4)}. That is
\code{\#1967ff}, the agency blue, at 40\% opacity. The button itself is orange
(\code{--accent}, line 12).

Orange and blue sit close to opposite each other on the colour wheel, so this does
not read as a subtle mismatch; it reads as a mistake. It is the most visible of the
five surviving blue literals catalogued below, and it is a one-line fix: the value
should be \code{rgba(255,103,25,.4)}, or better still derived from the token.

The cause is legible from the history. Commit \code{05210af} changed the tokens to
orange; this line never used a token, so the change could not reach it.
\end{honest}

\begin{jargon}{rgba and opacity}
\code{rgba(25,103,255,.4)} names a colour by red, green and blue amounts from 0 to
255, plus an opacity from 0 (invisible) to 1 (solid). It is the same blue as
\code{\#1967ff}, expressed in the form that allows transparency.
\end{jargon}

\begin{honest}{Neither button has a focus state}
Neither \code{.btn-primary} nor \code{.btn-ghost} defines \code{:focus} or
\code{:focus-visible}. A visitor navigating by keyboard rather than mouse gets
whatever outline the browser draws by default, which on a dark background is often
nearly invisible. Since these are the page's only interactive elements, this is the
most consequential accessibility gap in the stylesheet, and it is three lines to
fix.
\end{honest}

\paragraph{Lines 73 to 88: the two reveal mechanisms.}

\begin{lstlisting}[firstnumber=73,caption={\file{css/site.css}, lines 73--88}]

/* ─── REVEAL ─── */
.reveal {
  opacity: 0;
  transform: translateY(28px);
  transition: opacity 0.7s var(--ease), transform 0.7s var(--ease);
}
.reveal.visible { opacity: 1; transform: none; }

.stagger > * {
  opacity: 0;
  transform: translateY(20px);
  transition: opacity 0.6s var(--ease), transform 0.6s var(--ease);
  transition-delay: calc(var(--i, 0) * 90ms);
}
.stagger.visible > * { opacity: 1; transform: none; }
\end{lstlisting}

These sixteen lines are the entire animation contract between the stylesheet and
the JavaScript: the CSS defines both states and the transition between them, and
the JavaScript's only job is to add the class \code{visible} at the right moment.

\code{.reveal} starts invisible and 28 pixels low. \code{.reveal.visible}, which
matches an element carrying \emph{both} classes, returns it to fully visible and
its natural position. Because a transition is declared, the browser animates
between the two over 0.7 seconds rather than snapping.

The stagger is the same idea with one addition. \code{.stagger > *} means every
direct child of a staggered container, and line 86 delays each child by
\code{calc(var(--i, 0) * 90ms)}: the number the HTML wrote inline, multiplied by 90
milliseconds. Child 0 starts immediately, child 5 starts 450 milliseconds later.

\begin{jargon}{calc, and a custom property's fallback}
\code{calc()} performs arithmetic in CSS and can mix units. \code{var(--i, 0)}
reads the custom property \code{--i} and falls back to \code{0} if it was never
set, which keeps this rule safe on any element whose inline \code{--i} was
forgotten: it gets no delay instead of breaking.
\end{jargon}

\begin{keyidea}{Why the animation lives in CSS and the trigger lives in JavaScript}
This split is why the page degrades well, and it is also its one real hazard. If
JavaScript never runs, the class \code{visible} is never added and every revealed
element stays at opacity zero, invisible forever. That would be catastrophic, and
it is exactly why line 573 exists: the reduced-motion block at the end of the file
forces \code{.reveal} and \code{.stagger > *} back to full opacity. That block is
doing double duty as an accessibility feature and as a safety net.
\end{keyidea}

\paragraph{Lines 89 to 122: the navigation bar.}

\begin{lstlisting}[firstnumber=89,caption={\file{css/site.css}, lines 89--122}]

/* ─── NAV ─── */
.nav {
  position: fixed;
  top: 0; left: 0; right: 0;
  z-index: 100;
  padding: 18px 28px;
  display: flex;
  align-items: center;
  justify-content: space-between;
  transition: background 0.3s, border-color 0.3s, box-shadow 0.3s;
}
.nav.scrolled {
  background: rgba(5,7,13,0.88);
  border-bottom: 1px solid var(--border);
  box-shadow: 0 4px 24px rgba(0,0,0,.3);
  backdrop-filter: blur(16px);
  -webkit-backdrop-filter: blur(16px);
}
.nav-inner {
  max-width: var(--maxw);
  width: 100%;
  margin: 0 auto;
  display: flex;
  align-items: center;
  justify-content: space-between;
}
.nav-logo {
  font-family: var(--font-d);
  font-size: 17px;
  font-weight: 700;
  letter-spacing: -0.02em;
  color: var(--ink1);
}
\end{lstlisting}

\code{position: fixed} on line 92 pins the bar to the window rather than the
document, so it stays put while the page scrolls beneath it. Line 93 pins it to the
top and stretches it across.

\code{z-index: 100} on line 94 controls stacking order: when elements overlap, the
higher number is drawn in front. 100 is arbitrary and generous, chosen so nothing
else can accidentally cover the bar.

Lines 101 to 107 are the scrolled state, the class \file{main.js} adds once the page
has moved more than 40 pixels: a background at 88\% opacity, a bottom hairline, and
a shadow.

\code{backdrop-filter: blur(16px)} on line 105 is the interesting one: it blurs
whatever is \emph{behind} the element, producing the frosted-glass effect that only
works because the background is slightly transparent. Line 106 repeats it with a
\code{-webkit-} prefix for Safari, which needed the prefixed form for years.

\begin{honest}{A doubled flex layout}
\code{.nav} sets \code{display: flex} with \code{justify-content: space-between}
(lines 96 to 98), and \code{.nav-inner} sets the same three properties again (lines
112 to 114). The outer one has a single child, so its \code{space-between} has
nothing to space. It is harmless and redundant, and it makes the layout harder to
reason about than it needs to be. The sibling repository has exactly the same
doubling, so it came from the shared skeleton.
\end{honest}

Lines 116 to 122 style the logo. Note there is no rule for a logo mark or an
accent-coloured fragment, unlike the sibling: this logo is one span of plain text,
matching the simpler markup described in section 7.1.

\paragraph{Lines 123 to 133: the hero as a two-column grid.}

\begin{lstlisting}[firstnumber=123,caption={\file{css/site.css}, lines 123--133}]

/* ─── HERO ─── */
.hero {
  min-height: 100vh;
  display: grid;
  grid-template-columns: 1fr 1fr;
  align-items: center;
  gap: 0;
  padding-top: 80px;
  overflow: hidden;
}
\end{lstlisting}

This is where the two siblings genuinely diverge. The other repository's hero is a
single centred column over a field of drifting cards; this one is a two-column
grid, photograph on the left and words on the right.

\begin{jargon}{CSS Grid}
A layout system for two dimensions, as opposed to flexbox's one.
\code{grid-template-columns: 1fr 1fr} means two columns each taking one equal
share of the space. \code{1fr} is a fraction of what remains after any fixed sizes
are accounted for.
\end{jargon}

\code{min-height: 100vh} on line 126 makes the hero at least as tall as the window.
\code{vh} means one percent of the viewport height. \code{padding-top: 80px} on
line 131 leaves room for the fixed navigation bar, which is removed from the normal
flow and would otherwise sit on top of the content.

\paragraph{Lines 134 to 148: the photo column.}

\begin{lstlisting}[firstnumber=134,caption={\file{css/site.css}, lines 134--148}]

/* photo side */
.hero-photo-side {
  position: relative;
  height: 100vh;
  display: flex;
  align-items: center;
  justify-content: center;
  overflow: hidden;
}
.hero-photo-frame {
  position: relative;
  width: 100%;
  height: 100%;
}
\end{lstlisting}

\code{position: relative} on line 137 looks like it does nothing and is essential.
It establishes this element as the reference point for any absolutely positioned
descendant, which here means the three floating badges. Without it they would
position themselves against the whole window instead of against the photo column.

\begin{jargon}{Positioning: relative and absolute}
\code{position: absolute} takes an element out of the normal flow and places it at
coordinates you give, measured from the nearest ancestor that is itself positioned,
which in practice means the nearest one with \code{position: relative}. This
pairing, a relative parent and absolute children, is the standard way to place
things inside a box.
\end{jargon}

Line 138 sets the column's height to \code{100vh}, the full window height, which is
what makes the photograph bleed to the top and bottom edges of the screen.

\paragraph{Lines 149 to 179: the placeholder, styled with real care.}

\begin{lstlisting}[firstnumber=149,caption={\file{css/site.css}, lines 149--179}]

/* Photo placeholder - replace with <img src="assets/salman.jpg"> */
.hero-photo-placeholder {
  width: 100%;
  height: 100%;
  background: linear-gradient(160deg, #1a0800 0%, #4a1800 40%, #05070d 100%);
  display: flex;
  align-items: center;
  justify-content: center;
  flex-direction: column;
  gap: 16px;
  color: rgba(255,255,255,.15);
  font-family: var(--font-d);
}
.hero-photo-placeholder .placeholder-initial {
  font-size: 120px;
  font-weight: 700;
  letter-spacing: -0.05em;
  line-height: 1;
  background: linear-gradient(135deg, rgba(255,103,25,.6), rgba(255,140,74,.3));
  -webkit-background-clip: text;
  -webkit-text-fill-color: transparent;
  background-clip: text;
}
.hero-photo-placeholder .placeholder-note {
  font-size: 13px;
  color: rgba(255,255,255,.25);
  text-align: center;
  max-width: 220px;
  line-height: 1.5;
}
\end{lstlisting}

These thirty lines style the placeholder that section 7.1 identified as the most
consequential thing in the repository, and they are worth reading for what they say
about intent.

Line 154 fills the frame with a three-stop gradient running from a near-black
brown, through a deep orange, to the page background. That is a considered stand-in
for a photograph: it establishes the brand colour and fades into the page rather
than ending at a hard edge.

Lines 163 to 172 style the giant letter, and they use a genuinely advanced
technique.

\begin{keyidea}{How the letter is filled with a gradient}
You cannot pour a gradient into text directly. The trick on lines 168 to 171 is
three steps: give the element a gradient background, tell the browser to clip that
background to the shape of the text (\code{background-clip: text}), and then make
the text itself transparent (\code{-webkit-text-fill-color: transparent}) so the
clipped background shows through the letterforms.

Line 171 sets the standard \code{background-clip: text} and line 169 the prefixed
\code{-webkit-} form, which is still needed for full support. Note the third line
has no unprefixed equivalent at all: \code{-webkit-text-fill-color} is a
non-standard property that every major browser implements, and this effect depends
on it.
\end{keyidea}

Lines 173 to 179 style the instruction text: 13 pixels, white at 25\% opacity,
centred, capped at 220 pixels wide.

\begin{honest}{The placeholder note is deliberately styled to be readable}
This is the detail that settles how to describe the placeholder. The note is not
raw unstyled debris: it has been given a size, a colour, an alignment, a maximum
width and a line height. Somebody laid it out to be read.

Read charitably, that is a designer making an unfinished state look tidy in a
review. Read plainly, it means a visitor to the published page would see a legible
instruction addressed to a developer. Both readings are consistent with the
evidence and the second is what a visitor experiences, which is why section 14
treats it as the project's defining unfinished element.
\end{honest}

\paragraph{Lines 180 to 200: the two edge fades.}

\begin{lstlisting}[firstnumber=180,caption={\file{css/site.css}, lines 180--200}]

/* fade gradient on the right edge so text side blends */
.hero-photo-frame::after {
  content: '';
  position: absolute;
  top: 0; right: 0;
  width: 40%;
  height: 100%;
  background: linear-gradient(to right, transparent, var(--s0));
  pointer-events: none;
}
/* fade at bottom */
.hero-photo-frame::before {
  content: '';
  position: absolute;
  bottom: 0; left: 0; right: 0;
  height: 30%;
  background: linear-gradient(to top, var(--s0), transparent);
  pointer-events: none;
  z-index: 1;
}
\end{lstlisting}

Two pseudo-elements, and they solve the problem a hard-edged half-page image
creates.

\begin{jargon}{Pseudo-element}
An extra box CSS creates inside an element with no corresponding tag in the HTML,
written \code{::before} or \code{::after}. It needs the \code{content} property to
exist at all, which is why lines 183 and 193 set \code{content: ''}: an empty
string is enough to bring it into being. It is the standard way to draw a
decorative shape without polluting the markup.
\end{jargon}

The \code{::after} (lines 182 to 190) covers the rightmost 40\% of the photo column
with a gradient running from transparent to the page background, so the photograph
dissolves into the text column instead of ending at a visible seam. The
\code{::before} (lines 192 to 200) does the same across the bottom 30\%.

Both carry \code{pointer-events: none}, which makes them transparent to the mouse
so clicks pass through. Without it these two invisible layers would sit over the
photo column and swallow any interaction beneath them.

\begin{keyidea}{Why the bottom fade needs a z-index and the right fade does not}
Line 199 gives the \code{::before} a \code{z-index: 1}. The reason is ordering:
without an explicit stacking position, a \code{::before} is painted \emph{behind}
its element's content, while an \code{::after} is painted in front. So the right
fade works by default and the bottom fade would be hidden behind the placeholder
until lifted.

The badges then sit at \code{z-index: 2} (line 214), one layer above the bottom
fade, which is what keeps them crisp while the photograph dims underneath. Three
elements, three deliberate layers.
\end{keyidea}

\paragraph{Lines 201 to 241: the floating badges.}

\begin{lstlisting}[firstnumber=201,caption={\file{css/site.css}, lines 201--241}]

/* floating badges on the photo */
.hero-float-badge {
  position: absolute;
  background: rgba(13,18,32,0.85);
  backdrop-filter: blur(12px);
  -webkit-backdrop-filter: blur(12px);
  border: 1px solid var(--border);
  border-radius: 12px;
  padding: 12px 16px;
  display: flex;
  align-items: center;
  gap: 10px;
  z-index: 2;
  animation: badgeFloat 6s ease-in-out infinite;
}
.hero-float-badge:nth-child(2) { bottom: 25%; left: 8%;  animation-delay: 0s;   animation-duration: 7s; }
.hero-float-badge:nth-child(3) { bottom: 40%; right: 12%; animation-delay: 2s;  animation-duration: 6s; }
.hero-float-badge:nth-child(4) { top: 22%;   left: 10%;  animation-delay: 1s;  animation-duration: 8s; }

@keyframes badgeFloat {
  0%, 100% { transform: translateY(0); }
  50%       { transform: translateY(-8px); }
}
.badge-icon {
  width: 36px; height: 36px;
  border-radius: 10px;
  display: flex; align-items: center; justify-content: center;
}
.badge-icon-blue   { background: rgba(25,103,255,.15); }
.badge-icon-orange { background: rgba(255,103,25,.15); }
.badge-icon-green  { background: rgba(52,211,153,.12); }
.badge-icon svg { width: 18px; height: 18px; }
.badge-text-main {
  font-family: var(--font-d);
  font-size: 14px;
  font-weight: 700;
  color: var(--ink1);
  white-space: nowrap;
}
.badge-text-sub { font-size: 11px; color: var(--ink3); white-space: nowrap; }
\end{lstlisting}

Lines 203 to 216 define a badge: absolutely positioned, dark and semi-transparent
with a 12-pixel blur behind it, bordered, rounded, laid out as a row with a gap,
lifted to \code{z-index: 2}, and animated forever.

Note line 205: \code{rgba(13,18,32,0.85)}.

\begin{honest}{A colour left behind by the rebrand, in a subtler way}
\code{rgb(13,18,32)} is \code{\#0d1220}, which is the \emph{sibling repository's}
\code{--s1} surface colour. This file's own \code{--s1} is \code{\#0e1118}, a
slightly different, less blue dark. So the badge background is the other project's
card colour, written as a literal, at 85\% opacity.

Nobody would notice: the two values differ by a few points of blue and the badge
sits over a gradient at 85\% opacity. It is included here because it is evidence
about how the rebrand was done, not because it is a visible defect. The tokens were
updated and every literal in the file was left exactly as inherited.
\end{honest}

Lines 217 to 219 place the three badges individually, and this is where the HTML's
child order becomes load-bearing. \code{:nth-child(2)} matches the second child of
its parent, so these three rules assume the placeholder occupies child 1, as section
7.1 explained. Each badge gets a different corner, a different animation delay and a
different duration (7, 6 and 8 seconds), which is what stops the three drifting in
visible unison.

Lines 221 to 224 define the drift itself: a gentle rise of 8 pixels at the halfway
point and back. Grouping \code{0\%, 100\%} on one line is the idiom for an
animation returning to where it started.

\begin{jargon}{Keyframes}
An animation defined by describing the element at particular points through its
run, given as percentages of the total duration. The browser fills in every moment
between. \code{0\%} is the start and \code{100\%} the end.
\end{jargon}

Lines 230 to 232 are the three tinted icon chips: blue, orange and green, each the
matching colour at 12 to 15\% opacity. As section 7.1 noted, these three are the one
place blue is plainly intentional.

Line 239 and line 241 both set \code{white-space: nowrap}, stopping ``1.6M+ views''
and ``organic, for clients'' wrapping onto two lines and making the badge grow
taller than its neighbours.

\paragraph{Lines 242 to 301: the content column.}

\begin{lstlisting}[firstnumber=242,caption={\file{css/site.css}, lines 242--301}]

/* content side */
.hero-content-side {
  padding: 0 48px 0 20px;
  display: flex;
  flex-direction: column;
  justify-content: center;
  gap: 32px;
}
.hero-tag {
  display: inline-flex;
  align-items: center;
  gap: 8px;
  background: rgba(25,103,255,.1);
  border: 1px solid rgba(25,103,255,.25);
  border-radius: 100px;
  padding: 6px 14px;
  font-size: 13px;
  font-weight: 600;
  color: var(--accentH);
  width: fit-content;
}
.hero-h1 {
  font-family: var(--font-d);
  font-size: clamp(40px, 5vw, 70px);
  font-weight: 700;
  letter-spacing: -0.03em;
  line-height: 1.04;
  color: var(--ink1);
}
.hero-h1 em { font-style: normal; color: var(--accent); }
.hero-descriptor {
  font-family: var(--font-d);
  font-size: 20px;
  font-weight: 600;
  color: var(--ink3);
  letter-spacing: -0.01em;
}
.hero-actions {
  display: flex;
  align-items: center;
  gap: 12px;
  flex-wrap: wrap;
}
.hero-mini-stats {
  display: flex;
  align-items: center;
  gap: 28px;
  padding-top: 4px;
  border-top: 1px solid var(--borderS);
}
.hero-mini-stat { text-align: left; }
.hero-mini-num {
  font-family: var(--font-d);
  font-size: 24px;
  font-weight: 700;
  color: var(--ink1);
  letter-spacing: -0.03em;
}
.hero-mini-lbl { font-size: 12px; color: var(--ink3); }
\end{lstlisting}

Line 244 lays the right column out as a vertical flex column with a 32-pixel gap,
so its five children space themselves without margins. The asymmetric padding
(\code{0 48px 0 20px}) pulls the text slightly left, toward the fading edge of the
photograph.

Lines 251 to 263 are the small pill above the name, and it contains the second
missed rebrand.

\begin{honest}{A blue pill containing orange text}
Lines 255 and 256 tint the pill with \code{rgba(25,103,255,...)}, the agency blue,
while line 261 colours its text \code{var(--accentH)}, which is the orange
\code{\#ff8c4a}.

So the element reads its foreground from a token and its background from a
literal, and after the rebrand the two no longer agree. On screen it is a
blue-tinted capsule with orange text inside it, which looks like neither brand.

\code{width: fit-content} on line 262 is worth noting separately and is correct: in
a flex column, a child stretches to full width by default, so without this the pill
would span the entire column instead of hugging its text.
\end{honest}

Line 266 is the most useful single line in the file to learn from.

\begin{jargon}{clamp}
\code{clamp(minimum, preferred, maximum)} takes the preferred value but never lets
it fall below the minimum or rise above the maximum. Here the preferred value is
\code{5vw}, 5\% of the window width, so the headline scales continuously between a
floor of 40 pixels and a ceiling of 70. One line replaces what used to need several
media queries.
\end{jargon}

Line 268 sets \code{letter-spacing: -0.03em}, pulling letters closer together. Large
display type looks loose at its default spacing, and negative tracking at large
sizes is standard practice. The unit \code{em} means ``relative to this element's
font size'', so the tightening scales with the text.

Line 272 is the rule flagged in section 7.1: it takes the \code{<em>} tag, removes
the italics it normally implies, and colours it with the accent instead.

Lines 286 to 301 build the three small statistics: a flex row with a 28-pixel gap
above a hairline, each statistic a 24-pixel number over a 12-pixel label.

\paragraph{Lines 302 to 321: the work section, and the shared headings.}

\begin{lstlisting}[firstnumber=302,caption={\file{css/site.css}, lines 302--321}]

/* ─── WORK GRID ─── */
.work-section { padding: 120px 0 80px; }
.section-kicker {
  font-size: 12px;
  font-weight: 700;
  text-transform: uppercase;
  letter-spacing: 0.1em;
  color: var(--accent);
  margin-bottom: 14px;
}
.section-title {
  font-family: var(--font-d);
  font-size: clamp(28px, 4vw, 48px);
  font-weight: 700;
  letter-spacing: -0.025em;
  line-height: 1.1;
  color: var(--ink1);
  margin-bottom: 52px;
}
\end{lstlisting}

Although they sit under the work banner, \code{.section-kicker} and
\code{.section-title} are used by three sections. Their placement here is a small
organisational untidiness: a reader looking for the heading styles would reasonably
expect them near the top with the other shared pieces.

The kicker is the small uppercase line (``The work'', ``Case studies'') and the
title is the large heading beneath it. Note the tracking is positive on the kicker
and negative on the title. That is not an inconsistency but a rule of typography:
large text needs tightening and small uppercase text needs opening up or it becomes
a smudge.

\begin{honest}{The kicker's contrast is worth measuring}
\code{--accent} is now \code{\#ff6719}, a saturated orange, set at 12 pixels
uppercase on a near-black background. Small text needs more contrast than large
text, not less. This document does not have a measured contrast ratio for that
exact pairing and so does not assert a failure, but it is the first thing to check
in an accessibility audit. Orange on near-black generally fares better than the
sibling's blue on near-black, so this is the less likely of the two to fail.
\end{honest}

\paragraph{Lines 322 to 388: the project cards.}

\begin{lstlisting}[firstnumber=322,caption={\file{css/site.css}, lines 322--388}]

.work-grid {
  display: grid;
  grid-template-columns: repeat(3, 1fr);
  gap: 14px;
}
.work-card {
  background: var(--s1);
  border: 1px solid var(--border);
  border-radius: var(--r);
  overflow: hidden;
  transition: transform 0.22s var(--ease), border-color 0.22s, box-shadow 0.22s;
}
.work-card:hover {
  transform: translateY(-4px);
  border-color: rgba(25,103,255,.35);
  box-shadow: 0 16px 48px rgba(0,0,0,.4);
}
.work-thumb {
  aspect-ratio: 16/9;
  position: relative;
  display: flex; align-items: center; justify-content: center;
  overflow: hidden;
}
.work-thumb-overlay {
  position: absolute; inset: 0;
  background: rgba(5,7,13,.2);
  display: flex; align-items: center; justify-content: center;
  opacity: 0;
  transition: opacity 0.22s;
}
.work-card:hover .work-thumb-overlay { opacity: 1; }
.work-play {
  width: 48px; height: 48px;
  border-radius: 50%;
  background: rgba(255,255,255,.15);
  backdrop-filter: blur(8px);
  display: flex; align-items: center; justify-content: center;
  border: 1.5px solid rgba(255,255,255,.3);
}
.work-play svg { width: 20px; height: 20px; color: #fff; margin-left: 2px; }
.work-badge {
  position: absolute; top: 10px; right: 10px;
  background: rgba(5,7,13,.75);
  backdrop-filter: blur(8px);
  border: 1px solid var(--border);
  border-radius: 20px;
  padding: 4px 10px;
  font-size: 11px; font-weight: 600;
  color: var(--green);
}
.wt1 { background: linear-gradient(135deg,#2a1000,#5a2800); }
.wt2 { background: linear-gradient(135deg,#2a1800,#5a3000); }
.wt3 { background: linear-gradient(135deg,#1a0a00,#3a1800); }
.wt4 { background: linear-gradient(135deg,#300800,#602000); }
.wt5 { background: linear-gradient(135deg,#200e00,#482800); }
.wt6 { background: linear-gradient(135deg,#281400,#503000); }
.work-info { padding: 16px; }
.work-client {
  font-size: 11px; color: var(--ink3); font-weight: 600;
  text-transform: uppercase; letter-spacing: 0.06em; margin-bottom: 4px;
}
.work-title {
  font-family: var(--font-d);
  font-size: 15px; font-weight: 600;
  color: var(--ink1); letter-spacing: -0.01em;
}
\end{lstlisting}

\code{overflow: hidden} on line 332 is what makes the rounded corners work: without
it the gradient thumbnail inside would square off the top two corners it is
supposed to follow.

\code{aspect-ratio: 16/9} on line 341 reserves a widescreen-shaped area regardless
of width. This is worth appreciating: the layout will not shift when a real image
eventually loads, because the space is already the right shape. It is the modern
replacement for a well-known hack involving percentage padding.

The overlay (lines 346 to 353) sits over the whole thumbnail at zero opacity and
fades to one on hover, revealing the play button beneath a 20\% black wash. The play
button (lines 354 to 361) is a frosted circle: white at 15\% opacity with an
8-pixel backdrop blur and a brighter rim.

Line 362 gives the play triangle \code{margin-left: 2px}. That is not an accident.
A triangle's visual centre of mass sits left of its geometric centre, so a
perfectly centred play icon looks slightly left. Nudging it two pixels right is a
standard correction, and noticing it is a good sign that whoever built this cared
about detail.

\begin{honest}{The third missed rebrand}
Line 337 sets the hover border to \code{rgba(25,103,255,.35)}, the agency blue, on
a card in an orange-branded site. Same cause, same one-line fix as line 65.
\end{honest}

Lines 373 to 378 are the six placeholder gradients, and they \emph{were} rebranded:
every one is a two-stop orange or brown blend, where the sibling's equivalents run
through blue, teal, violet and green. So the rebrand did reach these six literals
and missed five others, which tells you it was done by eye rather than by search.

Lines 380 to 388 style the two lines of text under each thumbnail: a small uppercase
client name in the dimmest ink, and the result in the display typeface at 15
pixels. This is the only place in the file where the display typeface does ordinary
work at a small size.

\paragraph{Lines 389 to 419: the metrics band.}

\begin{lstlisting}[firstnumber=389,caption={\file{css/site.css}, lines 389--419}]

/* ─── METRICS ─── */
.metrics-section {
  padding: 80px 0;
  border-top: 1px solid var(--borderS);
  border-bottom: 1px solid var(--borderS);
}
.metrics-grid {
  display: grid;
  grid-template-columns: repeat(4, 1fr);
  gap: 2px;
}
.metric-item {
  padding: 36px 24px;
  text-align: center;
  border-right: 1px solid var(--borderS);
}
.metric-item:last-child { border-right: none; }
.metric-num {
  font-family: var(--font-d);
  font-size: clamp(34px, 4vw, 56px);
  font-weight: 700;
  letter-spacing: -0.03em;
  line-height: 1;
  color: var(--ink1);
  margin-bottom: 8px;
}
.metric-label {
  font-size: 13px; color: var(--ink3); font-weight: 500;
  text-transform: uppercase; letter-spacing: 0.06em;
}
\end{lstlisting}

Four equal columns with a 2-pixel gap, and the separators drawn as a right border
on every item (line 404) with the last one removed (line 406).

Note that this is the only place in the file that draws dividing lines this way.
The sibling repository contains a second, cleverer technique for the same job (a
grid whose gaps show a coloured layer beneath), used there on its services grid. That
technique does not appear here, because this file has no equivalent grid, so the
inconsistency the sibling has within itself does not exist in this one.

Lines 407 to 415 style the big numbers: display typeface, scaling between 34 and 56
pixels, with \code{line-height: 1}. That last is deliberate. The page-wide default
of 1.5 is right for paragraphs and wrong for a large number, where it leaves a
visible gap above and below.

\paragraph{Lines 420 to 479: the case-study cards, and three dead rules.}

\begin{lstlisting}[firstnumber=420,caption={\file{css/site.css}, lines 420--479}]

/* ─── CASE STUDIES ─── */
.cases-section { padding: 100px 0; }
.case-card {
  display: grid;
  grid-template-columns: 1fr 1fr;
  gap: 48px;
  align-items: center;
  background: var(--s1);
  border: 1px solid var(--border);
  border-radius: 24px;
  padding: 48px;
  margin-bottom: 16px;
  transition: border-color 0.22s;
}
.case-card:hover { border-color: rgba(25,103,255,.3); }
.case-card.reverse { direction: rtl; }
.case-card.reverse > * { direction: ltr; }
.case-visual {
  border-radius: var(--r);
  overflow: hidden;
  aspect-ratio: 16/10;
  position: relative;
}
.cv1 { background: linear-gradient(135deg,#2a1000,#5a2800); }
.cv2 { background: linear-gradient(135deg,#300800,#602000); }
.case-visual-inner {
  position: absolute; inset: 0;
  display: flex; align-items: flex-end;
  padding: 20px;
  background: linear-gradient(to top, rgba(5,7,13,.8) 0%, transparent 60%);
}
.case-visual-metrics {
  display: flex; gap: 12px; flex-wrap: wrap;
}
.cvm {
  background: rgba(5,7,13,.7);
  backdrop-filter: blur(8px);
  border: 1px solid var(--border);
  border-radius: 20px;
  padding: 6px 12px;
  display: flex; flex-direction: column; gap: 1px;
}
.cvm-num {
  font-family: var(--font-d);
  font-size: 16px; font-weight: 700;
  color: var(--ink1); letter-spacing: -0.02em;
}
.cvm-lbl { font-size: 10px; color: var(--ink3); text-transform: uppercase; letter-spacing: 0.06em; }
.case-content { display: flex; flex-direction: column; gap: 20px; }
.case-client-name {
  font-size: 12px; font-weight: 700;
  color: var(--accent); text-transform: uppercase; letter-spacing: 0.08em;
}
.case-title {
  font-family: var(--font-d);
  font-size: clamp(20px, 2.5vw, 30px);
  font-weight: 700; letter-spacing: -0.02em;
  color: var(--ink1); line-height: 1.2;
}
\end{lstlisting}

Each case card is a two-column grid with a 48-pixel gap: a visual panel and a block
of text, vertically centred against each other.

Lines 436 and 437 are the dead rules identified in section 7.1.

\begin{honest}{Three rules that have never applied, and the mechanism they were
replaced by}
\code{.case-card.reverse} means ``an element carrying both the class
\code{case-card} and the class \code{reverse}''. The HTML at line 321 writes
\code{class="case-card case-card-reverse reveal"}. A hyphen does not create a second
class, so this selector matches nothing, and neither does its partner on line 437 or
the third instance in the responsive block at line 559.

What they were trying to do is elegant and worth understanding even though it is
dead. Setting \code{direction: rtl} on a grid container reverses the order its
children are placed in, and setting \code{direction: ltr} on the children puts their
text back the right way round. Two lines, and a two-column layout mirrors itself.

What actually happens instead is that \file{index.html} sets \code{order:2} and
\code{order:1} inline on the two children, which achieves the same result by a
different route. So the design is correct, three rules are dead, and nothing
reported it. This is the clearest illustration in either sibling of what ``silent
failure'' means: an entire mechanism can be broken while the page looks right,
because a second mechanism quietly does the same job.
\end{honest}

\begin{honest}{The fourth missed rebrand}
Line 435 sets the case card's hover border to \code{rgba(25,103,255,.3)}, the agency
blue again.
\end{honest}

Lines 446 to 451 lay the metric chips over the bottom of each visual, with a
gradient from 80\% opaque page colour at the bottom to transparent at 60\% height, so
the chips sit on a dark base and the top of the visual stays clear. Lines 455 to 468
style the chips themselves: frosted, rounded, a number over a tiny uppercase label.

\paragraph{Lines 480 to 524: the booking panel.}

\begin{lstlisting}[firstnumber=480,caption={\file{css/site.css}, lines 480--524}]

/* ─── BOOK A CALL ─── */
.book-section { padding: 80px 0 140px; }
.book-wrap {
  background: var(--s1);
  border: 1px solid var(--border);
  border-radius: 24px;
  padding: 72px 48px;
  text-align: center;
  position: relative;
  overflow: hidden;
}
.book-wrap::before {
  content: '';
  position: absolute;
  top: -100px; left: 50%;
  transform: translateX(-50%);
  width: 600px; height: 300px;
  background: radial-gradient(ellipse, rgba(25,103,255,.1) 0%, transparent 70%);
  pointer-events: none;
}
.book-eyebrow {
  font-size: 12px; font-weight: 700;
  text-transform: uppercase; letter-spacing: 0.1em;
  color: var(--accent); margin-bottom: 16px;
}
.book-title {
  font-family: var(--font-d);
  font-size: clamp(28px, 4vw, 48px);
  font-weight: 700; letter-spacing: -0.025em;
  color: var(--ink1); margin-bottom: 12px;
}
.book-sub {
  font-size: 16px; color: var(--ink3); margin-bottom: 36px;
}
.book-cta-row {
  display: flex; align-items: center; justify-content: center;
  gap: 14px; flex-wrap: wrap;
}
.book-detail {
  display: flex; align-items: center; justify-content: center;
  gap: 6px; font-size: 13px; color: var(--ink3); margin-top: 20px;
}
.book-detail svg { width: 15px; height: 15px; }
.book-detail span { color: var(--ink2); }
\end{lstlisting}

Lines 492 to 500 add a glow above the panel using a \code{::before}: a 600 by 300
ellipse positioned 100 pixels above the panel's top edge and centred with the
standard \code{left: 50\%} plus \code{translateX(-50\%)} idiom. Because the panel has
\code{overflow: hidden}, only the lower part of the glow is visible, which makes it
read as light spilling in from above rather than as a floating oval.

\code{pointer-events: none} on line 499 appears for the third time in the file, for
the same reason as before: a decorative layer must not intercept the click aimed at
the booking button.

\begin{honest}{The fifth and last missed rebrand}
Line 498 makes that glow \code{rgba(25,103,255,.1)}, the agency blue, on the most
important element of the page. It is subtle at 10\% opacity, and it means the halo
around the call to action is the wrong brand colour.

That completes the catalogue: lines 65, 255, 256, 337, 435 and 498, five rules and
six declarations, all inherited from the sibling and all missed by a rebrand that
correctly updated the tokens and the six thumbnail gradients.
\end{honest}

\paragraph{Lines 525 to 535: the footer.}

\begin{lstlisting}[firstnumber=525,caption={\file{css/site.css}, lines 525--535}]

/* ─── FOOTER ─── */
footer {
  border-top: 1px solid var(--borderS);
  padding: 32px 28px;
  text-align: center;
  font-size: 13px;
  color: var(--ink3);
}
footer a { color: var(--ink3); transition: color 0.15s; }
footer a:hover { color: var(--ink1); }
\end{lstlisting}

Deliberately unremarkable: one hairline, centred dim text, and links that brighten
on hover over 0.15 seconds.

\paragraph{Lines 536 to 575: the responsive rules and the accessibility block.}

\begin{lstlisting}[firstnumber=536,caption={\file{css/site.css}, lines 536--575}]

/* ─── RESPONSIVE ─── */
@media (max-width: 960px) {
  .hero {
    grid-template-columns: 1fr;
    min-height: auto;
    padding-top: 80px;
  }
  .hero-photo-side {
    height: 55vw;
    min-height: 280px;
  }
  .hero-content-side {
    padding: 40px 28px 60px;
  }
  .hero-mini-stats { gap: 20px; }
  .work-grid   { grid-template-columns: repeat(2, 1fr); }
  .metrics-grid { grid-template-columns: repeat(2, 1fr); }
  .metric-item { border-right: none; border-bottom: 1px solid var(--borderS); }
  .metric-item:nth-child(odd) { border-right: 1px solid var(--borderS); }
  .metric-item:nth-child(3),
  .metric-item:nth-child(4)   { border-bottom: none; }
  .case-card { grid-template-columns: 1fr; gap: 28px; padding: 32px; }
  .case-card.reverse { direction: ltr; }
}

@media (max-width: 600px) {
  .work-grid  { grid-template-columns: 1fr; }
  .book-wrap  { padding: 48px 24px; }
  .hero-float-badge:nth-child(4) { display: none; }
}

@media (prefers-reduced-motion: reduce) {
  *, *::before, *::after {
    animation-duration: .001ms !important;
    transition-duration: .001ms !important;
  }
  .reveal, .stagger > * { opacity: 1; transform: none; }
  .hero-float-badge { animation: none; }
}
\end{lstlisting}

\begin{jargon}{Media query}
A block of CSS applying only when a condition about the display is true.
\code{@media (max-width: 960px)} means ``apply everything here only when the window
is 960 pixels wide or narrower''. Because it comes later in the file, its rules
override the equivalents above.
\end{jargon}

The 960-pixel block does the real work of making this design usable on a phone. The
hero collapses from two columns to one (line 540) and stops demanding the full
window height (line 541). The photo column becomes \code{55vw} tall with a
280-pixel floor (lines 545 and 546), so the photograph stays a sensible shape
instead of filling a tall narrow screen.

Lines 553 to 557 fold the metrics band from four columns to two and rebuild its
dividing lines: right border removed from everything, bottom border added to
everything, right border restored on odd items so there is still a line down the
middle, bottom border removed from items 3 and 4 so the block does not end in a
stray line. That is four rules to reposition two lines, and it is the price of the
border technique.

Line 558 collapses the case cards to a single column, and line 559 is the third dead
rule.

The 600-pixel block collapses the work grid to one column, tightens the booking
panel, and hides the third floating badge, which is the one positioned at the top
of the photo column where a phone has least room.

The final block, lines 568 to 575, is the best thing in the file.

\begin{keyidea}{What \code{prefers-reduced-motion} is, and why this block is better
than its sibling's}
Operating systems let people ask for reduced motion, usually because animation
causes nausea, dizziness or migraine. The browser passes that preference to the
page and this block responds.

It does three things. It crushes every animation and transition to a thousandth of
a millisecond, effectively instant rather than disabled, a deliberate choice because
setting them to zero can stop some browsers firing the events that mark an animation
as finished. Line 573 then forces every revealed and staggered element to full
opacity, because those start invisible and would otherwise never appear once their
transitions were neutered. That second line is the detail most implementations get
wrong.

Line 574 is the addition this file makes and the sibling does not:
\code{.hero-float-badge \{ animation: none; \}} stops the three badges drifting
outright. The blanket duration rule above would already have all but stopped them,
so this is belt and braces on the one continuous, looping animation on the page,
which is exactly the kind a motion-sensitive visitor would notice most.
\end{keyidea}

\begin{honest}{The reduced-motion block cannot reach the counting numbers}
This block governs CSS only. The four metric numbers are animated by JavaScript,
counting up over 2.2 seconds, and \file{main.js} never checks this preference. A
visitor who has asked for reduced motion still gets four numbers whirring upward.

The fix is four lines in \file{main.js}, given in section 7.3. This is the clearest
defect in the project after the placeholder: not because it is severe, but because
the intent to handle it correctly is demonstrably here, and expressed more carefully
than in the sibling, and still was not carried across to the other file.
\end{honest}

\subsubsection*{What it connects to}

\textbf{Named by:} \file{new-design-reference/index.html}, line 16, and nothing
else.

\textbf{Names:} no other file. It refers to no images and imports nothing. It names
font \emph{families}, which the HTML is responsible for downloading.

\textbf{Coupled to \file{index.html}} in three ways that no tool checks: by class
name throughout, by child position for the floating badges (lines 217 to 219), and
by the shared contract that JavaScript adds the class \code{visible}. One of those
couplings is already broken, at lines 436, 437 and 559.

\textbf{Coupled to \file{js/main.js}} through exactly two class names,
\code{scrolled} and \code{visible}, both of which the JavaScript adds and this file
defines.

\subsubsection*{Concepts introduced}

Custom properties, the cascade, flexbox, grid, pseudo-elements, pseudo-classes,
keyframe animations, transitions, media queries, \code{clamp}, \code{calc},
\code{aspect-ratio}, \code{backdrop-filter}, \code{background-clip: text}, stacking
with \code{z-index}, and \code{prefers-reduced-motion}. Each is defined at first use
above and collected in the glossary.

\subsubsection*{How it breaks}

\begin{itemize}
\item \textbf{Delete this file and the page still works}, in the sense that every
  word remains readable. It looks like an unstyled document, and the two-column
  hero becomes a stack.
\item \textbf{Rename a class here without renaming it in the HTML} and that element
  silently loses its styling. Lines 436, 437 and 559 are a live example.
\item \textbf{Insert an element before the placeholder in the HTML} and all three
  badges jump, because lines 217 to 219 position them by child number.
\item \textbf{Remove \code{pointer-events: none} from either fade} (lines 189, 198)
  and those invisible layers begin swallowing clicks.
\item \textbf{Remove \code{position: relative} on line 137} and the badges jump to
  positions measured against the whole window.
\item \textbf{A missing semicolon} makes the browser discard the rest of that
  declaration block. Unlike a program, CSS never reports an error: it drops what it
  cannot parse and carries on, which is why a stylesheet bug shows up as something
  looking wrong rather than as a message.
\end{itemize}

\subsubsection*{Honest findings for this file}

Five rules still using the agency blue after the rebrand (lines 65, 255, 256, 337,
435, 498); a sixth inherited literal in the badge background (line 205); four
tokens defined and never read; three dead rules for a class the HTML does not use
(lines 436, 437, 559); a doubled flex layout in the navigation; smooth scrolling
implemented twice; no focus states on either button; and a reduced-motion block
that is the best in either sibling and cannot reach the one animation living
outside CSS.

% ===========================================================================
\subsection{\file{new-design-reference/js/main.js}}
% ===========================================================================

\textbf{JavaScript, 64 lines.} Everything on the page that reacts. It does three
things and nothing else: the navigation bar grows a frosted background once you
scroll, the four metric numbers count up from zero when they come into view, and
most blocks fade upward as they arrive.

\begin{provenance}{new-design-reference/js/main.js}
Hand-written, and its own header comment says as much on line 2: \code{Zero
external dependencies.} That claim is true and checkable: the file imports nothing,
requires nothing, and the repository has no \file{package.json} that could supply a
library.

It arrived complete in the first commit \code{1f7ed12} and has never been edited.
It is a condensed version of the sibling repository's file: same three behaviours,
same helper, same counter arithmetic, 15 lines shorter, with the comments trimmed
and one behaviour dropped (the sibling watches a fifth selector this page does not
have). Which was written first is not recorded and is not guessed here.
\textbf{Confidence: certain for hand-written and dependency-free; the relationship
to the sibling is stated as a close resemblance, not as a direction of copying.}
\end{provenance}

\subsubsection*{Why it exists}

Without this file the page is complete, readable and static. Two of its three jobs
are pure decoration and would not be missed. The third matters: the metric numbers
are written into the HTML as the character \code{0} with their real values hidden
in data attributes, so with this file gone the metrics band permanently advertises
zero of everything.

That damage is smaller here than in the sibling, because the three hero statistics
(\file{index.html} lines 114 to 127) state three of the same four numbers as plain
text. A visitor without JavaScript still sees ``17+ projects'', ``1.6M+ organic
views'' and ``10K+ users served'' in the hero, and four zeros further down. Odd,
but not silent.

\begin{keyidea}{The single most important line in this file is in another file}
\file{index.html} loads this script at line 393, immediately before
\code{</body>}. Everything below runs the instant it loads, with no wrapper telling
it to wait for the document. That only works because by then every element already
exists. Move the \code{<script>} tag into the \code{<head>} and this file would run
against an empty page, find nothing, attach nothing, and report no error at all.
\end{keyidea}

\subsubsection*{Line by line}

\paragraph{Lines 1 to 13: the header, the wrapper, and the navigation bar.}

\begin{lstlisting}[language=Java,firstnumber=1,caption={\file{js/main.js}, lines 1--13}]
/* salmanadnan.com · main.js
   Scroll-driven animations. Zero external dependencies. */

(function () {
  'use strict';

  /* ── Nav gloss on scroll ── */
  const nav = document.querySelector('.nav');
  if (nav) {
    const tick = () => nav.classList.toggle('scrolled', window.scrollY > 40);
    window.addEventListener('scroll', tick, { passive: true });
    tick();
  }
\end{lstlisting}

Lines 1 and 2 are a comment, using the same \code{/* */} form as CSS.

Line 4 opens a wrapper that looks like noise and is not.

\begin{jargon}{Immediately invoked function expression}
Writing \code{(function () \{ ... \})();} defines a function and runs it once,
straight away, without giving it a name. The point is not the running: it is the
boundary. Anything declared inside is invisible outside.

Without it, the names \code{nav}, \code{animateCounter} and \code{observe} declared
below would become global, visible to any other script on the page, and any other
script using the same name would collide with them. One stray global called
\code{observe} is exactly the kind of bug that takes an afternoon to find.
\end{jargon}

Line 5 is \code{'use strict';}, switching the language into a stricter mode for
everything inside the wrapper.

\begin{jargon}{Strict mode}
A setting making JavaScript refuse things it would otherwise silently allow. The
most valuable one: without it, assigning to a name you never declared quietly
creates a global variable, so a typo such as \code{navv = 3} succeeds and does
nothing useful. In strict mode it throws an error and you find the typo
immediately.
\end{jargon}

Line 8 finds the navigation bar. \code{document.querySelector} takes a CSS selector,
exactly the syntax the stylesheet uses, and returns the first matching element or
the special value \code{null} if there is none.

\begin{jargon}{const, and why not var}
\code{const} declares a name whose binding cannot later be pointed at something
else. It is the default choice in modern JavaScript, with \code{let} for names that
genuinely must change. The older \code{var} has confusing scoping rules and appears
nowhere in this file, which is a small sign of care.
\end{jargon}

Line 9 is a guard: everything happens only \code{if (nav)}, that is, only if
something was found. This is the pattern making the whole file tolerant of a
changing page.

\begin{jargon}{Truthy and falsy}
JavaScript treats certain values as false in a condition: \code{false}, \code{0}, an
empty string, \code{null} and \code{undefined}. Everything else counts as true. So
\code{if (nav)} reads as ``if an element was found'', because
\code{querySelector} returns \code{null} when it finds nothing.
\end{jargon}

Line 10 does in one line what the sibling spreads over three, using an arrow
function with an implicit return.

\begin{jargon}{Arrow function}
\code{() => expression} is a compact way to write a function. Written without
braces, as here, it returns the expression's value automatically. The return value
is discarded in this case; the compactness is the only benefit.
\end{jargon}

\code{classList.toggle('scrolled', condition)} adds the class when the condition is
true and removes it when false. The condition is \code{window.scrollY > 40}, meaning
the page has scrolled more than 40 pixels. The stylesheet at lines 101 to 107
defines what \code{scrolled} looks like, and the transition on line 99 makes the
change fade rather than snap.

Line 11 registers the function to run on every scroll event. \code{\{ passive: true
\}} is a performance promise to the browser: this handler will never cancel the
scroll. Without the promise the browser must wait for the handler to finish before
knowing whether the scroll may proceed, which can make scrolling feel sticky. It is
the correct flag here and it is often forgotten.

Line 12 calls the function once immediately, handling the case where the page loads
already scrolled: a refresh halfway down, or an arrival on a link ending in
\code{\#book}. Without it the bar would be transparent over content until the first
scroll.

\paragraph{Lines 14 to 33: the counter.}

\begin{lstlisting}[language=Java,firstnumber=14,caption={\file{js/main.js}, lines 14--33}]

  /* ── Counter animation ── */
  function animateCounter(el) {
    const raw    = el.dataset.target || '0';
    const suffix = el.dataset.suffix || '';
    const target = parseFloat(raw.replace(/[^0-9.]/g, ''));
    const isFloat = raw.includes('.');
    const dec     = isFloat ? (raw.split('.')[1] || '').length : 0;
    const dur     = 2200;
    const start   = performance.now();

    const step = (now) => {
      const p     = Math.min((now - start) / dur, 1);
      const eased = 1 - Math.pow(1 - p, 3);
      const val   = eased * target;
      el.textContent = (isFloat ? val.toFixed(dec) : Math.round(val).toLocaleString()) + suffix;
      if (p < 1) requestAnimationFrame(step);
    };
    requestAnimationFrame(step);
  }
\end{lstlisting}

This is the most substantial piece of logic in the repository and is worth taking
slowly.

Line 17 reads the target. \code{el.dataset.target} is how JavaScript reads the
attribute \code{data-target} written in the HTML: strip the \code{data-} prefix and
what remains is a property of \code{dataset}. The \code{|| '0'} supplies a default:
if the attribute is missing, \code{dataset.target} is \code{undefined}, which is
falsy, so the expression evaluates to \code{'0'}.

Line 18 does the same for the suffix, defaulting to an empty string.

Line 19 turns the text into a number. \code{raw.replace(/[\^{}0-9.]/g, '')} deletes
every character that is not a digit or a dot, and \code{parseFloat} reads what is
left as a number.

\begin{jargon}{Regular expression}
A pattern for matching text, written between slashes. \code{[0-9.]} means ``a digit
or a dot''; a leading caret negates it, so it means ``anything that is not a digit
or a dot''. The \code{g} after the closing slash means global: replace every
occurrence, not just the first.
\end{jargon}

Line 20 asks whether the original text contained a dot, which is how the code
distinguishes \code{1.6} from \code{17}. Line 21 counts the digits after the dot so
\code{1.6} displays to one decimal place throughout rather than as
\code{1.5999999}. \code{(raw.split('.')[1] || '')} takes everything after the dot,
falling back to an empty string, and \code{.length} counts it.

\begin{jargon}{Ternary operator}
\code{condition ? a : b} evaluates to \code{a} when the condition is true and
\code{b} otherwise. It is an \code{if} statement compressed into an expression, and
it appears twice in this file, on lines 21 and 29.
\end{jargon}

Line 22 fixes the duration at 2200 milliseconds. Line 23 records the start time with
\code{performance.now()}, a clock designed for measuring durations precisely.

Lines 25 to 31 define one frame, and this is the heart of it.

\code{const p = Math.min((now - start) / dur, 1)} computes progress from 0 to 1. The
\code{Math.min(..., 1)} caps it, so if the browser is busy and a frame arrives late,
progress never exceeds 1 and the number never overshoots. That cap is the difference
between a counter landing on 17 and one that occasionally flashes 19 before
settling.

\code{const eased = 1 - Math.pow(1 - p, 3)} is the easing. Linear progress feels
mechanical; this curve starts fast and slows dramatically near the target, which is
what makes it feel like settling.

\begin{keyidea}{What the easing formula actually does}
Take progress \code{p}, subtract it from 1, cube the result, subtract that from 1.
At \code{p = 0.5}, halfway through in time, the formula gives $1 - 0.5^3 = 0.875$:
the counter is already 87.5\% of the way there. The whole second half of the
animation covers the last eighth of the distance. That asymmetry is the entire
effect.
\end{keyidea}

Line 29 renders the number. If the target had a decimal point, \code{toFixed(dec)}
formats it to that many places. Otherwise \code{Math.round} makes it whole and
\code{toLocaleString()} inserts the thousands separators appropriate to the
visitor's region, so a large number reads as \code{1,200} for a British visitor and
\code{1.200} for a German one. The suffix is glued on the end.

Line 30 asks for the next frame if the animation is unfinished, and line 32 starts
the first.

\begin{jargon}{requestAnimationFrame}
A request to run a function just before the browser's next screen repaint, about 60
times a second. Calling it repeatedly from inside itself, as here, is the standard
way to run an animation: it stays in step with the display, and the browser pauses
it entirely when the tab is in the background, which a timer would not.
\end{jargon}

\begin{honest}{The counter ignores the reduced-motion preference}
This function animates for 2.2 seconds regardless of whether the visitor has asked
their operating system to reduce motion. The stylesheet honours that preference with
unusual care, going further than the sibling by explicitly stopping the floating
badges, so the intent is plainly there and simply was not carried into this file.

The fix is small, and this document states it precisely so nobody has to design it:

\begin{lstlisting}[language=Java,caption={A four-line fix. This code is a
suggestion in this document and is NOT present in the repository.}]
const reduce = window.matchMedia('(prefers-reduced-motion: reduce)').matches;
if (reduce) {
  el.textContent = (isFloat ? target.toFixed(dec)
                            : target.toLocaleString()) + suffix;
  return;
}
\end{lstlisting}

Placed after line 23, that shows the final number immediately and skips the
animation, which is the correct behaviour: the information is preserved and only the
movement is removed.
\end{honest}

\paragraph{Lines 34 to 48: the observer helper.}

\begin{lstlisting}[language=Java,firstnumber=34,caption={\file{js/main.js}, lines 34--48}]

  /* ── Intersection observer helper ── */
  function observe(selector, className, threshold, cb) {
    const els = document.querySelectorAll(selector);
    if (!els.length) return;
    const io = new IntersectionObserver((entries) => {
      entries.forEach((e) => {
        if (!e.isIntersecting) return;
        e.target.classList.add(className || 'visible');
        if (cb) cb(e.target);
        io.unobserve(e.target);
      });
    }, { threshold: threshold || 0.18 });
    els.forEach((el) => io.observe(el));
  }
\end{lstlisting}

This is the best-designed thing in the repository: a small general-purpose helper
that the three lines following it reuse for three different jobs.

\begin{jargon}{IntersectionObserver}
A browser facility that watches elements and reports when they enter or leave the
visible area. Before it existed, the only way to do this was to listen for every
scroll event and measure each element's position by hand, which ran hundreds of
times a second and made pages stutter. The observer does the work outside the main
thread and reports only when something changes.
\end{jargon}

Line 37 finds every matching element. \code{querySelectorAll} returns a list, which
may be empty, rather than \code{null}. Line 38 guards against an empty list and
returns early: \code{!els.length} reads as ``if the count is zero'', because
\code{0} is falsy and \code{!} inverts it.

Lines 39 to 46 create the observer. Its callback receives one entry per element
whose visibility changed, and for each it skips anything leaving rather than
entering, adds the class (defaulting to \code{visible}), runs the optional extra
callback, and stops watching that element.

That last step is the one that matters most.

\begin{keyidea}{Why it stops watching}
Every reveal in this design is meant to happen once. Without \code{unobserve}, an
element scrolling out of view and back would fire again, and for the metrics that
would mean the numbers restarting their count every time the visitor scrolled past.
Unobserving also releases the browser's bookkeeping, so what starts as a dozen
watched elements steadily becomes none.
\end{keyidea}

Line 46 sets the threshold, defaulting to 0.18: the callback fires once 18\% of the
element is visible rather than at the first pixel. Waiting for a fraction is what
makes the reveal feel connected to the scroll instead of triggering just off screen.

\begin{honest}{The defaults cannot express a falsy value}
\code{threshold || 0.18} substitutes the default whenever the argument is falsy, and
\code{0} is falsy. A caller asking for a threshold of \code{0}, meaning ``fire the
moment a single pixel appears'', would silently get \code{0.18}. The same objection
applies to \code{className || 'visible'}, though an empty class name is not a
sensible request.

No caller passes \code{0}, so nothing is broken today. It is a latent trap worth
recognising, and the modern remedy is the \code{??} operator, which substitutes only
for \code{null} and \code{undefined} and leaves \code{0} alone.
\end{honest}

\paragraph{Lines 49 to 52: the three things being watched.}

\begin{lstlisting}[language=Java,firstnumber=49,caption={\file{js/main.js}, lines 49--52}]

  observe('.reveal', 'visible', 0.15);
  observe('.stagger', 'visible', 0.15);
  observe('.metric-num[data-target]', 'counted', 0.5, animateCounter);
\end{lstlisting}

Three calls, and the page's entire reveal behaviour. The first two add the class
\code{visible} to anything carrying \code{reveal} or \code{stagger}, both at a 0.15
threshold, and the stylesheet does everything after that.

The third is different. Its selector \code{.metric-num[data-target]} matches
elements carrying the class \emph{and} the attribute, a small piece of defensiveness:
a metric element without a target is skipped rather than animated to \code{NaN}. Its
threshold is 0.5, so the numbers start counting only when the band is half visible.

Note the callback is passed as a bare function name, \code{animateCounter}, rather
than wrapped in another arrow function. That is cleaner than the sibling, which
writes \code{(el) => \{ animateCounter(el); \}} for the same effect.

\begin{honest}{The class \code{counted} is added and styled nowhere}
The third call passes \code{'counted'} rather than \code{'visible'}, so each metric
number gains a class no CSS rule and no other JavaScript ever reads. It is harmless
and could serve as a marker for a developer inspecting the page, but nothing uses
it. Note also that the metric numbers are children of \code{.metrics-grid}, which is
a \code{.stagger}, so they are already faded in by the second call: this class is
genuinely spare.
\end{honest}

\paragraph{Lines 53 to 64: smooth anchor scrolling, and the close.}

\begin{lstlisting}[language=Java,firstnumber=53,caption={\file{js/main.js}, lines 53--64}]

  /* ── Smooth anchor scroll ── */
  document.querySelectorAll('a[href^="#"]').forEach((a) => {
    a.addEventListener('click', (e) => {
      const t = document.querySelector(a.getAttribute('href'));
      if (!t) return;
      e.preventDefault();
      t.scrollIntoView({ behavior: 'smooth', block: 'start' });
    });
  });

})();
\end{lstlisting}

Line 55 finds every link whose destination starts with a hash.
\code{a[href\^{}="\#"]} is an attribute selector and \code{\^{}=} means ``starts
with''. On this page that matches three links: the navigation's ``Book a call'', the
hero's ``Book a call'', and the hero's ``See the work''.

Line 57 looks up the destination using the link's own \code{href} as a selector,
which works because \code{\#book} is valid CSS meaning ``the element with id book''.

Line 58 returns early if the destination does not exist, leaving the browser to
handle the click normally. Line 59 then cancels the browser's default jump, but only
after the destination is confirmed, which is the correct order: cancel first and a
broken link would do nothing at all instead of at least trying.

Line 60 performs the scroll, with \code{block: 'start'} aligning the destination to
the top of the window.

\begin{honest}{This whole block is redundant, and it carries a latent crash}
\file{css/site.css} line 28 already sets \code{scroll-behavior: smooth} on the
document, making every anchor link glide with no JavaScript at all. These ten lines
reproduce that.

They also introduce a failure the CSS version does not have. A link written
\code{href="\#"}, a common way to write a placeholder link, would reach line 57 as
\code{document.querySelector('\#')}. That is not a valid selector, so the browser
throws a \code{SyntaxError}, and because it happens inside a click handler it would
break that click and leave an error in the console.

No such link exists today, so this is latent rather than live. The honest summary is
that the safest change to this file would be to delete lines 53 to 62 entirely: it
would remove a redundancy and a hazard in one edit, and nothing visible would
change.
\end{honest}

Line 64 closes and immediately runs the wrapper opened on line 4.

\subsubsection*{What it connects to}

\textbf{Named by:} \file{new-design-reference/index.html}, line 393, and nothing
else.

\textbf{Names:} no file at all. It imports nothing and fetches nothing.

\textbf{Its entire interface with the rest of the project is six strings}, worth
listing because they are the contract nothing checks:

\begin{table}[H]\centering\small
\begin{tabular}{@{}lll@{}}
\toprule
String & Used at & Where it must also exist \\
\midrule
\code{.nav}        & line 8  & \file{index.html} line 23, \file{site.css} line 91 \\
\code{scrolled}    & line 10 & \file{site.css} line 101 \\
\code{.reveal}     & line 50 & \file{index.html} (five places), \file{site.css} line 75 \\
\code{.stagger}    & line 51 & \file{index.html} (two places), \file{site.css} line 82 \\
\code{visible}     & lines 50--51 & \file{site.css} lines 80 and 88 \\
\code{.metric-num} & line 52 & \file{index.html} (four places), \file{site.css} line 407 \\
\bottomrule
\end{tabular}
\caption{Every point of contact between this file and the other two. Break any one
string in one place and the page keeps working, minus one behaviour, with no error
reported anywhere.}
\end{table}

\subsubsection*{How it breaks}

\begin{itemize}
\item \textbf{Load it in the \code{<head>} instead of at the end of the body} and
  all three behaviours vanish silently, since every lookup runs immediately and
  finds nothing.
\item \textbf{Rename a class in the HTML} and the corresponding behaviour stops,
  silently, because of the early-return guards.
\item \textbf{Delete this file} and the page is static, with four zeros in the
  metrics band. The hero statistics still read correctly, which limits the damage.
\item \textbf{Add a link with \code{href="\#"}} and clicking it throws a
  \code{SyntaxError} from line 57.
\item \textbf{A visitor who has asked for reduced motion} still gets the counting
  animation, because this file never checks.
\item \textbf{An ancient browser without \code{IntersectionObserver}} would throw on
  line 39 and lose all three reveal behaviours, which would leave every
  \code{.reveal} element permanently invisible were it not for the CSS
  reduced-motion block. Every browser in current use supports it, so this is a
  historical note rather than a live concern.
\end{itemize}

\subsubsection*{Honest findings for this file}

The counter ignores \code{prefers-reduced-motion}; the smooth-scroll block
duplicates a CSS feature and carries a latent \code{SyntaxError}; the
\code{counted} class is inert; and the \code{||} defaults on lines 42 and 46 cannot
express a falsy argument. Against those: the file is genuinely well written, with
consistent guards, a sensible shared helper, correct use of \code{unobserve} and
\code{passive}, no dependencies, and slightly tighter than its sibling in two
places.

% ===========================================================================
\subsection{\file{new-design-reference/netlify.toml}}
% ===========================================================================

\textbf{TOML, 11 lines.} Instructions for a hosting service called Netlify: how to
publish the folder, and four security headers to send with every response.

\begin{provenance}{new-design-reference/netlify.toml}
Hand-written, added in the first commit \code{1f7ed12} and never edited. It is
\textbf{byte for byte identical} to the sibling repository's copy, which was
verified by comparing the two files directly: \code{diff} reports no differences at
all. Every other file in this pair differs somewhere; this one does not.

That is itself the evidence for how it got here: it is the one file in the design
that carried no site-specific content, so it was reused without needing a single
edit. The structure follows Netlify's own documented configuration format.
\textbf{Confidence: certain for hand-written and for the identity with the sibling,
both checked directly; no specific template is claimed.}
\end{provenance}

\subsubsection*{Why it exists}

Today, for no reason at all: nothing reads it.

The honest reading of the history is that the site was first expected to be hosted
on Netlify, a common choice for a static page like this, and the plan changed two
commits later. Commit \code{692e162} on 2026-07-28 is titled in part ``VPS rsync
deploy'', which is the moment Netlify stopped being the destination. This file was
never removed.

\begin{jargon}{TOML}
A file format for configuration designed to be obvious to read. A name in square
brackets starts a section; \code{name = "value"} sets a setting. Double square
brackets, as on line 5, mean ``another entry in a list of these'', so a file can
carry several header rules.
\end{jargon}

\subsubsection*{Line by line}

\begin{lstlisting}[firstnumber=1,caption={\file{netlify.toml}, lines 1--11}]
[build]
  publish = "."
  command = ""

[[headers]]
  for = "/*"
  [headers.values]
    X-Frame-Options = "DENY"
    X-Content-Type-Options = "nosniff"
    Referrer-Policy = "strict-origin-when-cross-origin"
    Permissions-Policy = "camera=(), microphone=(), geolocation=()"
\end{lstlisting}

Lines 1 to 3 describe how to build the site. \code{publish = "."} means publish this
directory itself, and \code{command = ""} means there is no build step to run first.
Those two lines are the entire ``no build step'' decision, stated to a hosting
provider.

Line 5 opens a headers rule and line 6 applies it to \code{/*}, every path on the
site.

\begin{jargon}{HTTP header}
An extra line of information a web server sends alongside a file, which the visitor
never sees. Some describe the file; the four below are instructions to the browser
about what it is allowed to do with the page.
\end{jargon}

Line 8, \code{X-Frame-Options = "DENY"}, forbids any other site from embedding this
page inside a frame. That defends against clickjacking, where an attacker loads your
real page invisibly over their own and tricks a visitor into clicking your buttons
while believing they are clicking something else.

Line 9, \code{X-Content-Type-Options = "nosniff"}, tells the browser to believe the
declared type of every file rather than guessing from its contents. Guessing is a
historical browser behaviour that can be manipulated into treating an uploaded image
as a script.

Line 10, \code{Referrer-Policy = "strict-origin-when-cross-origin"}, limits what is
disclosed when a visitor follows a link away. Within this site the full address is
passed; to another site only the domain; to an insecure site nothing at all. This is
the modern default and a sensible privacy choice.

Line 11, \code{Permissions-Policy}, switches off three device capabilities for this
page entirely: camera, microphone and location. The empty parentheses mean ``nobody,
not even this page''. The page is static text and gradients and needs none of them,
so turning them off means a future compromise cannot quietly request them.

\begin{keyidea}{These four headers are good practice and they are not in effect}
Every one is a sensible defensive choice, and the file they live in is read by a
service this project does not use. The page is served by Caddy on the agency VPS,
which never sees this file.

Moving these four lines into the Caddy configuration on the server would take a few
minutes and make them real. As things stand, the repository looks like it has
thought about security headers and the live site has none of them.
\end{keyidea}

\subsubsection*{What it connects to}

Nothing in this repository refers to it. Netlify appears nowhere else in the tree
and no workflow reads it. The deploy workflow does not even exclude it, so if that
workflow were restored this file would be copied to the VPS where Caddy would
ignore it.

\subsubsection*{Concepts introduced}

TOML as a configuration format, and the four HTTP response headers above. No
concept here appears anywhere else in the repository, which is another way of
saying this file is disconnected from everything.

\subsubsection*{How it breaks}

It cannot break, because nothing reads it. The realistic failure is the opposite: a
future maintainer reads this file, concludes the site sends these headers, and
builds on an assumption that is false.

\begin{honest}{An unused configuration file is worse than no configuration file}
This is a small file and a real trap. It states four security guarantees the live
site does not provide, and nothing in the repository indicates it is inert. Either
delete it, or move its contents to the Caddy configuration and delete it. Leaving it
is the only option with no upside.

The identical file sits in the sibling repository with the identical problem, so
fixing one should mean fixing both.
\end{honest}

% ===========================================================================
\subsection{\file{new-design-reference/assets/ASSETS.md}}
% ===========================================================================

\textbf{Markdown, 32 lines.} A checklist of the eleven files somebody would need to
produce before this design could be considered finished, plus three notes about
placeholder content to replace.

\begin{provenance}{new-design-reference/assets/ASSETS.md}
Hand-written, added in the first commit \code{1f7ed12} alongside the design it
describes, and never edited. Written by the same person in the same sitting as the
HTML, which is evident from its references to the exact class names used there
(\code{.hero-photo-placeholder}, \code{wt1} through \code{wt6}, \code{cv1},
\code{cv2}) and from the fact that line 9 quotes, character for character, the
\code{<img>} tag already written into the HTML's banner comment at line 34.
\textbf{Confidence: high, from the commit history and the internal references.}
\end{provenance}

\subsubsection*{Why it exists}

It is a handover note. The design was built with placeholder visuals, and this file
records precisely what would replace each one and what HTML change each replacement
needs. As an artefact it is genuinely good practice: it converts ``this is not
finished'' from a feeling into a list.

It is also the document that makes the state of this project unambiguous. The
placeholder in the hero could be argued to be a design choice; a checklist listing
the photograph as \emph{required before launch} cannot.

\begin{jargon}{Markdown}
A way of writing formatted text in a plain file, using punctuation instead of menus:
\code{\#} starts a heading, backticks mark code, and pipe characters draw a table. It
is designed to be readable as plain text and is the standard format for notes kept
alongside code.
\end{jargon}

\subsubsection*{Line by line}

\begin{lstlisting}[firstnumber=1,caption={\file{ASSETS.md}, lines 1--19}]
# Assets checklist: salmanadnan.com

Drop real assets here, then update the HTML as noted.

## Required before launch

| File | Size | Where used | HTML change needed |
|---|---|---|---|
| `salman.jpg` | 1200px tall min, 2:3 ratio | Hero photo left side | Replace `.hero-photo-placeholder` div with `<img src="/assets/salman.jpg" alt="Salman Adnan" style="width:100%;height:100%;object-fit:cover;object-position:top;">` |
| `thumbs/work-1.jpg` | 800x450px | Work grid card 1 | Add `style="background-image:url('/assets/thumbs/work-1.jpg')"` to `.wt1` |
| `thumbs/work-2.jpg` | 800x450px | Work grid card 2 | Same pattern for `.wt2` through `.wt6` |
| `thumbs/work-3.jpg` | 800x450px | Work grid card 3 | |
| `thumbs/work-4.jpg` | 800x450px | Work grid card 4 | |
| `thumbs/work-5.jpg` | 800x450px | Work grid card 5 | |
| `thumbs/work-6.jpg` | 800x450px | Work grid card 6 | |
| `case-1.jpg` | 1200x750px | Case study 1 visual | Add `background-image` to `.cv1` |
| `case-2.jpg` | 1200x750px | Case study 2 visual | Add `background-image` to `.cv2` |
| `favicon.svg` | 32x32 | Browser tab | Add `<link rel="icon" href="/assets/favicon.svg">` in `<head>` |
| `og-image.png` | 1200x630px | Social share | Update `og:image` meta tag |
\end{lstlisting}

Line 1 is the title and line 3 the instruction: put the real files in this folder,
then change the HTML as the table says.

Lines 7 and 8 open a table, line 8 being the separator row Markdown requires between
a table's headings and its body.

The eleven rows are the substance, and reading them tells you exactly what the
design is pretending to have.

\textbf{The photograph} is first, and it is the one that matters. Line 9 specifies
it precisely: at least 1200 pixels tall, in a 2:3 portrait ratio, and it gives the
complete replacement tag including \code{object-fit: cover} so the image fills the
frame without distorting, and \code{object-position: top} so a face stays in shot
when the frame crops. Whoever wrote this knew exactly what was needed and did not
have the file.

\textbf{Six project thumbnails} at 800 by 450, which is exactly the 16 by 9 shape the
stylesheet reserves with \code{aspect-ratio} at line 341. Line 11 gives the pattern
for the remaining five rather than repeating it, which is why rows 12 to 15 have an
empty final column.

\textbf{Two case-study images} at 1200 by 750, matching the \code{16/10} aspect ratio
set at \file{site.css} line 441. That the two ratios differ, and that each listed
asset matches its own, is a small sign of care.

\textbf{A favicon} and \textbf{a social preview image}, each with the exact tag
needed to install it.

\begin{honest}{None of the eleven exist, and the live site has at least one of them}
The \file{assets/} folder contains exactly one file: this checklist. Every asset it
calls ``required before launch'' is absent, so the design was never finished to the
standard its own author set.

The comparison that makes this sharp: fetching the live salmanadnan.com on
2026-08-11 found \code{<meta property="og:image" content=
"https://salmanadnan.com/assets/og.png">} along with explicit width and height
tags. So the social preview image listed here as required does exist, for the build
that replaced this one. This design was superseded by a version that finished the
checklist.
\end{honest}

\begin{lstlisting}[firstnumber=20,caption={\file{ASSETS.md}, lines 20--32}]

## Numbers to update

Search for `data-target` in `index.html` (metrics band) and replace all four values with real numbers.
Search for the three `hero-mini-num` divs in the hero and update those too.

## Booking link

Search for `YOUR_CAL_LINK` in `index.html` and replace with the actual booking URL.

## Client names

Three `<!-- UPDATE -->` comments in the work grid and case study sections mark placeholder client entries.
\end{lstlisting}

Line 23 is the metrics instruction, and line 24 is the one worth noticing.

\begin{keyidea}{This file knows about the duplication and tells you to fix it in
both places}
Line 23 says to search for \code{data-target} in the metrics band. Line 24 then says
to search for the three \code{hero-mini-num} divs and update those too.

That is the author documenting, at the time of writing, that the same three numbers
live in two places in two formats, as section 7.1 described. It was a recognised
cost, accepted deliberately, and written down. That is a materially different thing
from an accident, and it is why this document treats the duplication as a design
trade rather than a defect.

The sibling repository has no equivalent note because it has no hero statistics, and
consequently no fallback when JavaScript does not run.
\end{keyidea}

Line 28 tells the reader to search for \code{YOUR\_CAL\_LINK}.

\begin{honest}{Two of this file's three instructions are stale}
Searching \file{index.html} for \code{YOUR\_CAL\_LINK} finds nothing. The literal
placeholder was replaced on 2026-07-28 by commit \code{014406e}, which set the link
to \code{cal.com/salman-adnan/meeting}, a real personal calendar. So somebody
following this instruction searches, finds nothing, and cannot tell whether the job
is done or the instruction is wrong. Here the job \emph{is} done, unlike in the
sibling repository where the same stale instruction hides a link still pointing at
another project's calendar.

Line 32 is stale in a different way. It refers to ``three \code{<!-{}- UPDATE -{}->}
comments in the work grid and case study sections'' marking placeholder client
entries. There are no such comments anywhere in \file{index.html}: its
\code{UPDATE} notes live in the banner comments at the top of each section and none
of them is about client names. Either they were removed or they were never written.

Only the metrics instruction on lines 23 and 24 is verifiably still accurate, and
the \code{UPDATE} comment at \file{index.html} line 256 agrees with it.
\end{honest}

\subsubsection*{What it connects to}

It describes \file{new-design-reference/index.html} and
\file{new-design-reference/css/site.css} by name and by class, and it is the only
file in the repository documenting the project's intent. Nothing reads it: it is for
people.

The deploy workflow explicitly excludes it from publication (line 82 of the
workflow), which is the right instinct, and that exclusion no longer matches its
path. See section 7.6.

\subsubsection*{Concepts introduced}

Markdown tables; \code{object-fit} and \code{object-position}, the two properties
that control how an image fills a frame it does not exactly match; and the standard
sizes for a favicon (32 by 32) and an Open Graph preview image (1200 by 630).

\subsubsection*{How it breaks}

As a document it cannot break. As a source of truth it already has: two of its three
instructions describe a state of the code that no longer exists or never did, which
is the characteristic failure of documentation kept beside code and not updated with
it.

\subsubsection*{Honest findings for this file}

Two stale instructions (lines 28 and 32); eleven assets specified in useful detail
and never delivered; and, to its credit, one instruction (lines 23 and 24) that is
still exactly right and documents a known duplication rather than hiding it.

% ===========================================================================
\subsection{\file{.github/workflows-disabled/deploy.yml.disabled}}
% ===========================================================================

\textbf{YAML, 103 lines.} The automation that used to publish this site: it watched
for pushes, checked for a marker in the commit message, copied the files to the
agency's server over an encrypted connection, then checked that the site still
answered. It is currently switched off.

\begin{provenance}{.github/workflows-disabled/deploy.yml.disabled}
Hand-written, added in commit \code{05210af} on 2026-07-28 at the path
\file{.github/workflows/deploy.yml}, rewritten in \code{692e162} the same day when
the deploy target changed from Netlify to the VPS, fixed four times on 2026-07-29,
and finally renamed to its present path in \code{0dd839e}.

It is \textbf{nearly identical to the sibling repository's workflow}: comparing the
two files directly shows differences on exactly six lines, all of them the site's
own name or path (the header comment, the serving path, the workflow name, the
\code{rsync} destination and the verification URL). Every mechanism, every retry
loop and every flag is the same.

This repository is where the evidence for those mechanisms lives. Its history
contains two commits explaining the retry loops by name (\code{6b1aa25},
\code{6fcffee}); the sibling has the loops and no such commits. The
\code{[deploy]} convention itself is workspace-wide and predates both.
\textbf{Confidence: certain for hand-written and for the rename history, both
recorded in git; certain for the six-line difference, checked directly.}
\end{provenance}

\subsubsection*{Why it exists}

To remove a manual step. Without it, publishing means somebody connecting to the
server and copying files by hand, which is slow and easy to get wrong. With it,
publishing is a consequence of writing a commit message in a particular way.

\begin{jargon}{YAML}
A file format for configuration where structure is expressed by indentation rather
than brackets. Two spaces of indent mean ``this belongs to the line above''. It is
unforgiving about indentation, which is its main drawback: one misplaced space
changes the meaning.
\end{jargon}

\begin{jargon}{GitHub Actions, and a workflow}
A service running jobs on GitHub's computers when something happens in a
repository. The jobs are described in a workflow file in \file{.github/workflows/}.
Each job runs on a fresh, empty machine, which is why the file below installs its
own SSH key every time.
\end{jargon}

\subsubsection*{Line by line}

\paragraph{Lines 1 to 19: the header comment and the trigger.}

\begin{lstlisting}[firstnumber=1,caption={\file{deploy.yml.disabled}, lines 1--19}]
# salmanadnan.com: deploy to VPS when commit subject starts with [deploy].
#
# A push to main always checks the marker. It DEPLOYS only when the commit
# subject line STARTS with [deploy], matching every other agency repository.
# Mentioning the marker anywhere else in the message does nothing.
#
# The site is served by Caddy from /root/salmanadnan.com on the VPS.
# Caddy picks up the new files immediately, no restart needed.
#
# Required repository secrets:
#   VPS_HOST     46.202.194.85
#   VPS_USER     deploy
#   VPS_SSH_KEY  private key for that user
name: Deploy salmanadnan.com

on:
  push:
    branches: [main]

\end{lstlisting}

Lines 1 to 13 are a comment block, and it is the best documentation in the
repository. In YAML a \code{\#} starts a comment. It states the marker rule twice
and in capitals, which tells you it had caught somebody out; it names the serving
path and notes that Caddy needs no restart; and it lists the three secrets required.

\begin{honest}{The server's address is written in the comment explaining it is a
secret}
Line 11 records \code{VPS\_HOST 46.202.194.85} in plain text immediately under the
heading ``Required repository secrets''. Whatever the value of keeping the address
in GitHub's secret store, writing it here defeats it.

Proportion matters: an IP address is not a credential, this repository is private,
and the machine is publicly reachable anyway. This document confirmed independently
that salmanadnan.com resolves to exactly this address today, which incidentally
makes the comment currently accurate. The genuine cost is that it is permanent in
the git history, so it becomes misleading rather than merely redundant if the server
is ever replaced. The private key, which is what actually matters, is correctly
absent.
\end{honest}

Line 14 names the workflow, the label shown in GitHub's interface.

Lines 16 to 18 set the trigger: every push to \code{main}, and nothing else. Not on
pull requests, not on a schedule, not manually. That last absence is worth noting:
there is no \code{workflow\_dispatch}, so there is no way to publish without making
a commit.

\paragraph{Lines 20 to 38: the marker check.}

\begin{lstlisting}[firstnumber=20,caption={\file{deploy.yml.disabled}, lines 20--38}]
jobs:
  check-deploy:
    name: Check Deploy Marker
    runs-on: ubuntu-latest
    outputs:
      should_deploy: ${{ steps.check.outputs.should_deploy }}
    steps:
      - id: check
        env:
          COMMIT_MSG: ${{ github.event.head_commit.message }}
        run: |
          SUBJECT="${COMMIT_MSG%%$'\n'*}"
          if [[ "$SUBJECT" == "[deploy]"* ]]; then
            echo "should_deploy=true" >> "$GITHUB_OUTPUT"
            echo "Deploy marker found on the subject line."
          else
            echo "should_deploy=false" >> "$GITHUB_OUTPUT"
            echo "No deploy marker: skipping deploy."
          fi
\end{lstlisting}

The first job exists only to answer a yes-or-no question, and runs on a fresh Ubuntu
machine to do it. Lines 24 and 25 declare that it produces an output called
\code{should\_deploy}, which the second job will read.

Line 28 passes the commit message in as an environment variable rather than pasting
it into the script text. That distinction is a genuine security measure.

\begin{keyidea}{Why the commit message goes through an environment variable}
If line 31 had been written \code{SUBJECT="\$\{\{
github.event.head\_commit.message \}\}"}, GitHub would paste the commit message
directly into the script before the shell ever saw it. Anyone who could write a
commit message could then write shell commands into it and have them run on the
machine holding the deploy key.

Passing it as an environment variable and reading it as \code{\$COMMIT\_MSG} means
the message is data, never code. This is a well-known class of vulnerability in
GitHub Actions, and this file gets it right.
\end{keyidea}

Line 31 extracts the subject line. \code{\$\{COMMIT\_MSG\%\%\$'\textbackslash n'*\}}
is shell syntax meaning ``delete the longest match of a newline followed by
anything, from the end'', leaving everything before the first newline.

Line 32 is the test. \code{"\$SUBJECT" == "[deploy]"*} compares with a trailing
wildcard, which is why the marker must \emph{begin} the subject.

Lines 33 and 36 write the answer to \code{\$GITHUB\_OUTPUT}, a file GitHub reads to
collect a step's outputs. Lines 34 and 37 print a human-readable explanation into
the log, a small kindness to whoever is wondering why nothing deployed.

\paragraph{Lines 39 to 63: starting the deploy, and the SSH setup.}

\begin{lstlisting}[firstnumber=39,caption={\file{deploy.yml.disabled}, lines 39--63}]

  deploy:
    name: Deploy to VPS
    runs-on: ubuntu-latest
    needs: check-deploy
    if: needs.check-deploy.outputs.should_deploy == 'true'
    steps:
      - uses: actions/checkout@v4

      - name: Authorize SSH
        env:
          KEY: ${{ secrets.VPS_SSH_KEY }}
          HOST: ${{ secrets.VPS_HOST }}
        run: |
          mkdir -p ~/.ssh && chmod 700 ~/.ssh
          printf '%s\n' "$KEY" > ~/.ssh/deploy_key
          chmod 600 ~/.ssh/deploy_key
          for i in 1 2 3 4 5; do
            if ssh-keyscan -T 10 -H "$HOST" >> ~/.ssh/known_hosts 2>/tmp/keyscan.err; then
              [ -s ~/.ssh/known_hosts ] && break
            fi
            echo "ssh-keyscan attempt $i failed: $(cat /tmp/keyscan.err 2>/dev/null)"
            sleep 5
          done
          [ -s ~/.ssh/known_hosts ] || { echo "::error::could not fetch the host key"; exit 1; }
\end{lstlisting}

Line 43 makes this job wait for the first, and line 44 runs it only if the answer was
yes. Everything below is skipped entirely on an ordinary commit.

Line 46 checks out the repository onto the machine using GitHub's own action.
Without this the machine would be empty.

Lines 53 to 55 install the private key: create the \file{.ssh} folder, write the key
into it, and set permissions so only the owner can read it. \code{chmod 600} is not
optional; SSH refuses to use a key file other users can read, and its error message
when it does is famously unhelpful.

\begin{jargon}{SSH key pair, and the host key}
Connecting securely involves two separate keys and confusing them is common. The
\emph{deploy key} here is the private half of a pair proving \emph{we} are allowed
in; its public half sits on the server. The \emph{host key} is the server's own
identity, which the client checks to be sure it is talking to the right machine and
not an impostor. Lines 53 to 55 install the first; lines 56 to 63 fetch the second.
\end{jargon}

Lines 56 to 62 are the retry loop, and this repository records exactly why it exists.
\code{ssh-keyscan} asks the server for its host key with a 10-second timeout,
appending the result to \file{known\_hosts} and any complaint to a temporary file.
\code{[ -s ~/.ssh/known\_hosts ]} then tests that the file is non-empty, because
\code{ssh-keyscan} can succeed as a command while producing nothing at all. That
distinction is the subtle part: checking the exit code alone would not have been
enough.

\begin{keyidea}{Two commits name this problem outright}
Commit \code{6b1aa25} is titled ``Retry after transient keyscan failure'' and
\code{6fcffee} ``Retry, runner IP lottery''. Read together they are a complete
diagnosis: GitHub's runners come from a wide pool of addresses, some of them were
being refused or timing out, and the same push succeeded on a second attempt.

This is the clearest example in the whole pair of repositories of a challenge with a
scar. The sibling has the identical retry loop and no commit explaining it, so
anybody reading only that repository would see defensive code with no stated cause.
\end{keyidea}

Line 63 fails the whole job with a message if five attempts produced nothing.
\code{::error::} is GitHub's syntax for marking a log line as an error so it surfaces
in the interface.

\begin{keyidea}{Why fetching the host key each time is a real trade-off}
Fetching a server's identity at the moment you connect, and trusting whatever comes
back, means you are not verifying anything: an attacker positioned between the
runner and the server could supply their own. The secure alternative is to store the
expected host key as a secret and compare.

That is a real weakness and it is a widespread convention, and another project in
this workspace, \code{passbolt-digitalise}, does it the stricter way by keeping a
\file{.github/known\_hosts} file in the repository. Two approaches to the same
problem inside one workspace is worth knowing about.
\end{keyidea}

\paragraph{Lines 64 to 83: copying the files.}

\begin{lstlisting}[firstnumber=64,caption={\file{deploy.yml.disabled}, lines 64--83}]

      - name: Sync files
        env:
          HOST: ${{ secrets.VPS_HOST }}
          USER: ${{ secrets.VPS_USER }}
        run: |
          retry() {
            for attempt in 1 2 3; do
              "$@" && return 0
              echo "attempt $attempt failed: $*"
              [ "$attempt" -lt 3 ] && sleep $(( attempt * 15 ))
            done
            return 1
          }
          retry rsync -az --no-owner --no-group --omit-dir-times --delete \
            -e "ssh -i ~/.ssh/deploy_key -o BatchMode=yes -o ConnectTimeout=25" \
            --exclude '.git' \
            --exclude '.github' \
            --exclude 'assets/ASSETS.md' \
            ./ "$USER@$HOST:/root/salmanadnan.com/"
\end{lstlisting}

Lines 70 to 77 define a small shell function running whatever it is given up to
three times, waiting longer between attempts: 15 seconds, then 30. \code{"\$@"}
means ``all the arguments, run as a command''. This is the second retry loop, and
like the first it is evidence of a step that failed in practice.

Line 78 is the copy itself. Each flag earns its place:

\begin{table}[H]\centering\small
\begin{tabular}{@{}ll@{}}
\toprule
Flag & What it does \\
\midrule
\code{-a} & archive mode: recurse, and preserve timestamps, permissions and links \\
\code{-z} & compress during transfer \\
\code{--no-owner} & do not try to set file ownership on the far end \\
\code{--no-group} & do not try to set group ownership either \\
\code{--omit-dir-times} & do not try to set directory timestamps \\
\code{--delete} & remove files on the server that no longer exist here \\
\bottomrule
\end{tabular}
\caption{The \code{rsync} flags, and why each is present.}
\end{table}

The three \code{--no} and \code{--omit} flags all work around the same thing: the
\code{deploy} user does not own the target directory and cannot change ownership or
directory timestamps. Without them \code{rsync} reports errors and fails. These
flags are the surviving trace of the problem commit \code{431ebbf} was fixing.

\code{--delete} deserves respect: it makes the server an exact mirror, so removing a
file here removes it there. That is usually what you want, and it is also the flag
that turns a mistaken source path into a wiped site.

Line 79 tells \code{rsync} how to connect: use this key, never prompt for a password
(\code{BatchMode=yes}, essential because there is no human to answer), and give up
after 25 seconds.

Lines 80 to 82 exclude three things: the git history, the workflow folder, and the
internal checklist.

Line 83 is the source and destination. \code{./} means the current directory's
\emph{contents}, and the trailing slash is significant: without it \code{rsync}
would copy the directory itself into the target, producing a nested folder.

\begin{honest}{This line would now take down the live site}
\code{./} was the right source when \file{index.html} sat at the repository root.
Since commit \code{0dd839e} moved everything one level down, \code{./} contains a
single folder called \file{new-design-reference}. Restoring this workflow unchanged
would copy that folder to the server and, because \code{--delete} is set, remove the
real site's files in the same operation.

This is not hypothetical harm. Section 1.2 established that salmanadnan.com is live,
answering with \code{200}, and resolves to the very address on line 11. The correct
source after the move would be \code{new-design-reference/}.

The same commit that moved the files disabled this workflow, so the two are
consistent; the trap exists only for somebody who re-enables one without reading the
other.
\end{honest}

\begin{honest}{The exclusion for the checklist no longer matches}
\code{--exclude 'assets/ASSETS.md'} was correct when that path was correct. The file
now lives at \file{new-design-reference/assets/ASSETS.md}, which this pattern does
not match, so a restored workflow would publish the internal checklist, including
its instruction to drop a photograph in, to the public internet.
\end{honest}

\paragraph{Lines 84 to 103: verifying, and cleaning up.}

\begin{lstlisting}[firstnumber=84,caption={\file{deploy.yml.disabled}, lines 84--103}]

      - name: Verify
        env:
          HOST: ${{ secrets.VPS_HOST }}
          USER: ${{ secrets.VPS_USER }}
        run: |
          for i in $(seq 1 12); do
            code=$(ssh -i ~/.ssh/deploy_key -o BatchMode=yes "$USER@$HOST" \
              'curl -sf -o /dev/null -w "%{http_code}" --max-time 10 https://salmanadnan.com/' || echo 000)
            case "$code" in
              2*|3*) echo "Site answering with $code after ${i} check(s)."; exit 0 ;;
            esac
            echo "waiting, got $code ($i/12)"
            sleep 5
          done
          echo "::warning::site did not respond 2xx/3xx; check Caddy on the VPS"

      - name: Clean up key
        if: always()
        run: rm -f ~/.ssh/deploy_key
\end{lstlisting}

The verification step asks, up to twelve times over about a minute, whether the site
is answering. It does this by connecting to the server and running \code{curl}
\emph{there} rather than from the runner, which tests the server's own view of
itself.

The \code{curl} flags: \code{-s} quiet, \code{-f} treat an error status as a
failure, \code{-o /dev/null} throw the page away,
\code{-w "\%\{http\_code\}"} print just the status number, and a 10-second limit.
The \code{|| echo 000} supplies a fallback so a failed connection yields the harmless
string \code{000} rather than an empty one.

Line 93 accepts any status beginning with 2 or 3, meaning success or redirect, and
exits successfully. That tolerance is well judged and this document can confirm it
matters: \code{http://salmanadnan.com/} currently answers \code{308}, a permanent
redirect to the secure version, so a check accepting only \code{200} could fail on a
perfectly healthy site.

\begin{honest}{A failed verification is only a warning}
Line 99 uses \code{::warning::} rather than \code{::error::}, and the step does not
exit with a failure code. So if the site never came back, the workflow would still
finish green with a warning buried in the log.

Given that the preceding step ran \code{rsync --delete} against a live site, this is
the one moment you would want a loud red failure and an email. It is the clearest
correctness weakness in the workflow, and it is one word to fix.
\end{honest}

Lines 101 to 103 delete the private key from the runner. \code{if: always()} makes
this run even when an earlier step failed, which is the point: a key must not be left
behind because something else went wrong. In practice GitHub destroys the machine
anyway, so this is belt and braces, and it is the right instinct.

\subsubsection*{What it connects to}

\textbf{Reads:} three GitHub secrets (\code{VPS\_SSH\_KEY}, \code{VPS\_HOST},
\code{VPS\_USER}) and the commit message of the push that triggered it.

\textbf{Writes:} the contents of \file{/root/salmanadnan.com/} on the VPS at
\code{46.202.194.85}, which Caddy serves at \code{https://salmanadnan.com/}.

\textbf{Excludes:} \file{.git}, \file{.github} and \file{assets/ASSETS.md}.

\textbf{Runs when:} never, currently. It sits in \file{.github/workflows-disabled/}
with the extension \file{.disabled}, and GitHub only reads
\file{.github/workflows/*.yml}.

\subsubsection*{Concepts introduced}

YAML; GitHub Actions jobs, steps, outputs and secrets; environment variables as a
defence against injection; SSH key pairs and host keys; \code{rsync} and its
archive and delete semantics; retry loops with growing backoff; and the
\code{[deploy]} marker convention shared across this workspace.

\subsubsection*{How it breaks}

\begin{itemize}
\item \textbf{Restored unchanged}, it copies the wrong directory and, with
  \code{--delete}, takes down the live salmanadnan.com.
\item \textbf{A missing or wrong \code{VPS\_SSH\_KEY}} fails at the \code{rsync}
  step after three attempts. Commit \code{cfa65d7} fixed this.
\item \textbf{A \code{deploy} user without permission to traverse into
  \code{/root}} fails the same way. Commit \code{431ebbf} fixed this, and it is the
  least obvious of the failures: the user can hold the key, connect successfully, and
  still be unable to reach the target directory.
\item \textbf{GitHub's runner being refused by the server} fails the host-key fetch.
  Commits \code{6b1aa25} and \code{6fcffee} are both this.
\item \textbf{A commit whose subject does not start with \code{[deploy]}} does
  nothing at all, by design, and the log says so.
\item \textbf{The site failing to come back} produces a warning and a green tick,
  which is the failure mode worth fixing first.
\end{itemize}

\subsubsection*{Honest findings for this file}

The server address is recorded in a comment describing it as a secret; the
\code{rsync} source path and the \file{ASSETS.md} exclusion are both stale after the
move, and here the consequence is a live site going down rather than a parked one;
the host key is trusted on first sight rather than pinned, unlike
\code{passbolt-digitalise} in the same workspace; and a failed post-deploy check
produces a warning rather than a failure. Against those: the commit message is
handled safely as data rather than as code, both retry loops are well built and
documented by the commits that caused them, the status check correctly tolerates
redirects, and the key is cleaned up unconditionally.


% ===========================================================================
\section{Generated and third-party files}
% ===========================================================================

\textbf{There are none.} All six tracked files were written by hand. No lock file,
no minified bundle, no vendored dependency, no binary asset, and no generator
banner. Checked rather than assumed:

\begin{lstlisting}[language=bash,caption={The check that was run, and its result}]
$ git grep -lIE '(DO NOT EDIT|auto-?generated|@generated|Generated by)'
(no output)
\end{lstlisting}

So the tier 2 count in Appendix A is zero and every file is walked line by line.

% ===========================================================================
\section{Configuration and secrets}
% ===========================================================================

Nothing is configurable at run time: no settings file, no environment file, no
\file{.env.example}. To change what the page says, you edit the HTML.

Three values behave like configuration and are written into the source:

\begin{table}[H]\centering\small
\begin{tabular}{@{}lll@{}}
\toprule
Value & Where & What it controls \\
\midrule
\code{https://cal.com/salman-adnan/meeting} & \file{index.html} line 362 & the booking calendar \\
\code{salman@digitalise.agency} & \file{index.html} line 386 & the footer contact address \\
The nineteen tokens & \file{css/site.css} lines 5--25 & the whole colour scheme \\
\bottomrule
\end{tabular}
\caption{Values that behave like settings but are written into the source.}
\end{table}

\subsection{Secrets}

The repository contains \textbf{no passwords, keys or tokens}, and should not: a
page running entirely in a visitor's browser can keep no secret.

Three are referred to by the workflow and stored in GitHub rather than here:
\code{VPS\_HOST}, \code{VPS\_USER} and \code{VPS\_SSH\_KEY}.

\begin{honest}{The server's IP address is written in a comment}
Line 11 of the workflow records \code{VPS\_HOST 46.202.194.85} in a comment
directly under the heading ``Required repository secrets''. Storing the value in a
comment defeats storing it as a secret.

Proportion matters: an IP address is not a credential, this repository is private,
and the machine is reachable from the internet anyway; this document confirmed
independently that salmanadnan.com resolves to exactly that address. The genuine
cost is that it is now permanent in the git history, so it becomes misleading
rather than merely redundant if the server is ever replaced. The private key,
which is what actually matters, is correctly absent.
\end{honest}

% ===========================================================================
\section{Dependencies}
% ===========================================================================

\textbf{No build-time dependencies at all}: no \file{package.json}, no
\file{requirements.txt}, nothing to install.

Two run-time dependencies, both services on other people's computers, both reached
by the visitor's browser rather than by the project:

\begin{table}[H]\centering\small
\begin{tabular}{@{}lll@{}}
\toprule
Dependency & Reached from & What it is for \\
\midrule
Google Fonts & \file{index.html} lines 13--15 & downloads the Space Grotesk and Inter typefaces \\
cal.com & \file{index.html} line 362 & shows the calendar when ``Book a call'' is pressed \\
\bottomrule
\end{tabular}
\caption{The two external services the page depends on.}
\end{table}

\begin{honest}{Fonts are loaded from Google, which the agency's main site
deliberately does not do}
The agency's main website (\code{digitalise-website} in this workspace) self-hosts
its fonts. This page fetches them from \code{fonts.googleapis.com} on every visit,
so every visitor's browser contacts Google, which is a privacy and consent question
in some jurisdictions, and the page's appearance depends on a third party staying
up.

It does degrade gracefully: both font definitions list fallbacks
(\code{ui-sans-serif}, \code{system-ui}, \code{sans-serif}), so text still appears
in whatever the device provides. The layout shifts; nothing disappears.
\end{honest}

% ===========================================================================
\section{How to run it}
% ===========================================================================

Nothing to install, nothing to start. Two ways work, and only one looks right:

\begin{lstlisting}[language=bash,caption={The way that works properly}]
cd new-design-reference
python3 -m http.server 8000
# then open http://localhost:8000/ in a browser
\end{lstlisting}

Opening the file directly instead renders the page completely unstyled: black text
on white, everything in one column.

\begin{keyidea}{Why opening the file directly looks broken}
The page asks for its stylesheet at \code{/css/site.css}. The leading slash means
``from the very top of this site''. Opened off a disk, the top is the root of your
hard drive, so the browser looks for \code{/css/site.css} there and finds nothing.

Writing \code{css/site.css} without the slash would have worked both ways. The
slashed form was chosen, correct for the live site and inconvenient for local
inspection. This one character is the most likely thing to make a newcomer think
the project is broken when it is not.
\end{keyidea}

\subsection{What you should see}

A two-column hero filling the window: on the left, an orange gradient with a giant
letter \code{S} and the words ``Drop Salman's photo here'', with three small
frosted badges drifting gently over it; on the right, the name ``Salman Adnan''
with the surname in orange, a tagline, two buttons and three small statistics.
Below that, in order: six project cards, four numbers that count up as they arrive,
two large case-study cards, a booking panel and a one-line footer.

% ===========================================================================
\section{The agent-facing surface}
% ===========================================================================

Several projects in this workspace ship their own Claude Code commands, skills and
hooks. \textbf{This project ships none.} Checked rather than assumed:

\begin{lstlisting}[language=bash,caption={The check, and its result}]
$ find . -maxdepth 2 \( -name 'CLAUDE.md' -o -name 'skills.md' \)
(no output)
$ find .claude -type f 2>/dev/null
(no output)
\end{lstlisting}

\begin{honest}{No README, no CLAUDE.md, and the important fact lives in a commit
message}
Somebody cloning this repository is given six files and no statement of what it is,
whether it is live, or why everything sits inside a folder called
\file{new-design-reference}. The only explanation anywhere is the message on commit
\code{0dd839e}, invisible unless you know to run \code{git log}.

For a repository whose central fact is ``this is parked, the live site is
elsewhere'', that belongs in a README where it cannot be missed. The same gap
exists in the sibling \code{ai-digitalise}, so it is a habit rather than an
oversight, and it would take five minutes to close in both.
\end{honest}

The two PDFs in \file{docs/explainer/} are produced by a workspace-level command,
\code{/explain-projects}, which lives outside this repository and owns only that
folder.

% ===========================================================================
\section{Deployment}
% ===========================================================================

\begin{figure}[H]\centering
\begin{tikzpicture}[node distance=9mm and 14mm]
  \node[box]                (push) {A commit is pushed\\to \code{main}};
  \node[dec, right=of push] (chk)  {subject starts\\with \code{[deploy]}?};
  \node[box, right=of chk]  (ssh)  {Fetch the host key,\\write the SSH key};
  \node[box, below=of ssh]  (sync) {\code{rsync} the files\\to the VPS};
  \node[box, left=of sync]  (ver)  {Ask the VPS to fetch\\the live URL};
  \node[box, left=of ver]   (stop) {Stop, do nothing};
  \draw[flow] (push) -- (chk);
  \draw[flow] (chk)  -- node[lbl]{yes} (ssh);
  \draw[flow] (chk)  |- node[lbl,pos=0.3]{no} (stop);
  \draw[flow] (ssh)  -- (sync);
  \draw[flow] (sync) -- (ver);
\end{tikzpicture}
\caption{The publishing pipeline as written. Drawn for reference: it is switched
off, twice over, and has been since 2026-07-29.}
\end{figure}

Section 7.6 walks it line by line. In short: a push to \code{main} checks whether
the commit's first line begins with the literal text \code{[deploy]}. If so, the
files are copied to the VPS with \code{rsync} and the workflow then asks the VPS to
fetch the live address to confirm the site answers. If not, nothing happens.

\begin{keyidea}{The \code{[deploy]} marker must start the subject line}
Every repository in this workspace shares this convention. Writing ``fix the header
and \code{[deploy]}'' does nothing at all: the check is
\code{"\$SUBJECT" == "[deploy]"*}, which tests the \emph{beginning} of the string.
\end{keyidea}

\subsection{What would happen if it were restored}

The \code{rsync} command copies \code{./}, the repository root, to
\code{/root/salmanadnan.com/} on the server. When it was written,
\file{index.html} was at that root; since \code{0dd839e} it is one level down. So
a restored workflow would copy a folder called \file{new-design-reference} to the
server, and because \code{--delete} is set it would remove the real site's files
in the same operation. The live salmanadnan.com would go down.

The same commit that moved the files disabled the workflow, so the two changes are
consistent. The trap exists only for somebody who re-enables one without reading
the other.

\begin{honest}{An exclusion is stale as well}
The \code{rsync} command excludes \code{assets/ASSETS.md} so the internal
checklist is not published. After the move that file is at
\file{new-design-reference/assets/ASSETS.md}, which the pattern no longer matches,
so a restored workflow would publish the checklist to the public internet.
\end{honest}

% ===========================================================================
\section{Honest findings}
% ===========================================================================

Each is argued at the point in the walkthrough where it arises. Collected here so
it can be read as a list.

\subsection{The design was never finished}

\begin{enumerate}
\item \textbf{The hero photo is a visible placeholder.} Half the opening screen
  reads ``Drop Salman's photo here / assets/salman.jpg'' over a gradient. On a
  personal brand page, whose entire proposition is a person, the missing element is
  the person. This is the finding that matters most.
\item \textbf{None of the eleven assets in \file{ASSETS.md} exist.} The
  \file{assets/} folder contains only the checklist. No photograph, no project
  thumbnails, no case-study screenshots, no favicon and no social preview image.
\item \textbf{Six ``UPDATE'' comments and one ``PLACEHOLDER'' comment remain} in
  \file{index.html}, some describing work that was done and some describing work
  that was not.
\end{enumerate}

\subsection{The rebrand to orange was left half done}

\begin{honest}{Five places still use the agency blue on an orange-branded site}
Commit \code{05210af} is titled ``Differentiate to orange brand''. The token block
was duly changed: \code{--accent} became \code{\#ff6719}. But five rules further
down still hard-code the agency blue \code{rgba(25,103,255,...)} rather than
reading the token:

\begin{table}[H]\centering\small
\begin{tabular}{@{}llp{5.4cm}@{}}
\toprule
Line & Rule & What a visitor sees \\
\midrule
65  & \code{.btn-primary:hover} & an orange button casting a blue shadow \\
255--256 & \code{.hero-tag} & a blue-tinted pill containing orange text \\
337 & \code{.work-card:hover} & a project card outlined in blue on hover \\
435 & \code{.case-card:hover} & the same on the case-study cards \\
498 & \code{.book-wrap::before} & a blue glow above the booking panel \\
\bottomrule
\end{tabular}
\end{table}

The blue shadow under the orange button (line 65) is the most visible: the two
colours are close to complementary, so it reads as a mistake rather than a choice.

None of these is a crash and all five are one-line fixes. Together they are a
clean illustration of what the token system is for and what happens when a colour
is written literally instead: the rebrand changed the tokens and could not reach
the literals.
\end{honest}

\subsection{Dead code}

\begin{enumerate}\setcounter{enumi}{3}
\item \textbf{Two stylesheet rules for the reversed case-study card have never
  applied.} \file{site.css} lines 436 and 437 target \code{.case-card.reverse};
  \file{index.html} line 321 writes \code{class="case-card case-card-reverse"}.
  A hyphen is not a second class, so the selector never matches. A third rule in
  the responsive block (line 559) is dead for the same reason. The intended layout
  happens anyway, by inline \code{order} properties on lines 322 and 340, so
  nothing looks wrong and three rules do nothing.
\item \textbf{Four tokens are defined and never read}: \code{--s2},
  \code{--accent2}, \code{--rsm} and \code{--shadow}. Counted mechanically; each
  returns zero uses. \code{--accent2} is the agency blue, kept as a secondary
  colour and then never used as one.
\item \textbf{The class \code{counted} is added by the JavaScript and styled
  nowhere.}
\item \textbf{\file{netlify.toml} configures a host that is not used.} The deploy
  path is \code{rsync} to a VPS running Caddy; Netlify appears nowhere else.
\end{enumerate}

\subsection{Content and consistency}

\begin{enumerate}\setcounter{enumi}{7}
\item \textbf{The three hero statistics duplicate three of the four in the metrics
  band}, written as plain text in one place and as data attributes in the other. So
  the same numbers exist twice in two formats, and \file{ASSETS.md} line 24 knows
  it: it tells the reader to update both.
\item \textbf{Two of the three instructions in \file{ASSETS.md} are stale.} It
  says to search for \code{YOUR\_CAL\_LINK}, which no longer exists, and refers to
  three \code{<!-- UPDATE -->} comments in the work grid and case sections, which
  do not exist either.
\item \textbf{The six project cards show a play button on hover and play
  nothing.} Unlike the sibling repository, these cards do not set
  \code{cursor: pointer}, so the promise is weaker here.
\end{enumerate}

\subsection{Accessibility}

\begin{enumerate}\setcounter{enumi}{10}
\item \textbf{The counting numbers ignore the reduced-motion preference.} The
  stylesheet honours it carefully, and goes further than the sibling by explicitly
  stopping the floating badges. The counter is animated in JavaScript over 2.2
  seconds and never checks. The fix is four lines, given in section 7.3.
\item \textbf{Neither button has a focus state}, so keyboard navigation relies on
  whatever outline the browser draws, often nearly invisible on a dark background.
  These are the page's only interactive elements.
\item \textbf{The metric values are invisible without JavaScript.} The real
  numbers live in data attributes and the visible text is \code{0}, so a search
  engine or a visitor with JavaScript disabled sees a page claiming zero of
  everything.
\item \textbf{The navigation logo is not a link.} It is a \code{<span>}, so
  there is no way back to the top of the page by clicking it. On a single-page
  site this is defensible and it differs from the sibling, where the logo is a
  link.
\end{enumerate}

% ===========================================================================
\section{The project's story}
% ===========================================================================

Parts 1 to 4 and part 8 are fact, taken from the repository and its history. Parts
5, 6 and 7 are advice and judgement, labelled where they begin.

\subsection{Deliverables}

\begin{itemize}
\item A complete, self-contained one-page personal brand website: 395 lines of
  HTML, 575 of CSS, 64 of JavaScript, no dependencies to install.
\item A hosting configuration for Netlify setting four security response headers.
\item A GitHub Actions workflow publishing to the agency VPS over SSH, gated on a
  \code{[deploy]} marker, with retries on both the host-key fetch and the file
  copy, and a post-deploy check that the site answers.
\item An asset checklist naming the eleven files needed to finish the design.
\end{itemize}

What it did not produce is a live site. salmanadnan.com serves a different build.

\subsection{Timeline and effort}

Ten commits across two calendar days.

\begin{table}[H]\centering\small
\begin{tabular}{@{}llp{7.0cm}@{}}
\toprule
Date & Commit & Subject \\
\midrule
2026-07-28 & \code{1f7ed12} & Launch salmanadnan.com, a personal brand site \\
2026-07-28 & \code{014406e} & Update booking link to cal.com/salman-adnan/meeting \\
2026-07-28 & \code{05210af} & \code{[deploy]} Differentiate to orange brand, add deploy workflow \\
2026-07-28 & \code{692e162} & \code{[deploy]} Replace all placeholder content with real Daztify project data, orange theme, VPS rsync deploy \\
2026-07-29 & \code{066917a} & \code{[deploy]} Trigger deploy after VPS secrets added \\
2026-07-29 & \code{cfa65d7} & \code{[deploy]} Fix VPS deploy key \\
2026-07-29 & \code{431ebbf} & \code{[deploy]} Fix /root traverse permission for deploy user \\
2026-07-29 & \code{6b1aa25} & \code{[deploy]} Retry after transient keyscan failure \\
2026-07-29 & \code{6fcffee} & \code{[deploy]} Retry, runner IP lottery \\
2026-07-29 & \code{0dd839e} & Move new design into new-design-reference and disable auto-deploy \\
\bottomrule
\end{tabular}
\caption{The complete commit history. Subjects are reproduced as written except
that \code{1f7ed12} and \code{6fcffee} use a long dash where this table uses a
comma; that character is not printable in this document's typeface.}
\end{table}

The shape is identical to the sibling's and slightly worse: four commits on day
one building the thing, then five consecutive commits on day two fighting the
deploy, then one commit switching it off. Nothing in the history suggests the
design was difficult. Everything suggests publishing it was.

\textbf{This is calendar time, not effort.} Git records when commits happened, not
hours spent, and no document here records effort, so none is claimed.

\subsection{Challenges}

Each is claimed because the repository carries a scar for it.

\textbf{1. SSH access to the VPS, across three commits.} \code{066917a} (the
workflow existed before the secrets did), \code{cfa65d7} (the key was wrong), and
\code{431ebbf} (the key was right but the \code{deploy} user could not traverse
into \code{/root} to reach the target folder). The last is a genuinely unobvious
Unix permission problem and its workaround survives in three \code{rsync} flags.

\textbf{2. GitHub's runners being refused, across two more commits.}
\code{6b1aa25} is titled ``Retry after transient keyscan failure'' and
\code{6fcffee} ``Retry, runner IP lottery''. That second title is the diagnosis
stated outright: GitHub's runners come from a wide pool of addresses and some were
being refused. This repository is where the evidence lives; the sibling has the
same two retry loops and no commit explaining them.

\textbf{3. Deciding what the brand was.} Commit \code{05210af} exists solely to
differentiate this site from the agency's blue. As section 14.2 shows, the change
reached the tokens and missed five literals, which is the trace of a decision made
after the design was written rather than before.

\textbf{4. Being overtaken.} The final commit is a retreat, not a fix. The page was
finished and working when something else became the live site. That is the largest
event recorded here and it was not technical.

\subsection{Outcomes}

\begin{itemize}
\item \textbf{Built, and unfinished.} It renders and is responsive, and the single
  most prominent element on the page is still a placeholder asking for a
  photograph.
\item \textbf{Not live, deliberately.} Parked under \file{new-design-reference/}
  with the workflow disabled twice over.
\item \textbf{Superseded.} salmanadnan.com answers from the same VPS with a
  different build that has a proper title, a correct canonical link and a social
  preview image, none of which this version has.
\item \textbf{One thing did survive.} The live site uses the same booking link,
  \code{cal.com/salman-adnan/meeting}, set here in commit \code{014406e}.
\item \textbf{Dormant} since 2026-07-29.
\end{itemize}

\subsection{What to do (judgement)}

\emph{From here on this is advice, not a record of project decisions.}

\begin{itemize}
\item \textbf{Read \code{git log} before touching anything.} The only statement of
  this repository's status lives in a commit message.
\item \textbf{Serve the folder over HTTP to view it}, never open the file directly.
\item \textbf{Read it beside \code{ai-digitalise}.} The two are the same skeleton
  with different content, and the differences between them are where the
  interesting decisions are.
\item \textbf{Keep the \code{[deploy]} convention and both retry loops} if the
  workflow is ever restored. They were paid for over five commits.
\item \textbf{Treat the hero as the reusable asset.} The two-column layout with
  fading edges and drifting frosted badges is the best-designed thing in either
  sibling, and it is about eighty lines of CSS.
\end{itemize}

\subsection{What not to do (judgement)}

\begin{itemize}
\item \textbf{Do not re-enable the workflow without changing the \code{rsync}
  source path.} It would copy a folder rather than a page, with \code{--delete}
  set, and take the live salmanadnan.com down.
\item \textbf{Do not publish this design without the photograph.} The placeholder
  is the largest element on the page and it is addressed to the developer.
\item \textbf{Do not trust \file{ASSETS.md}.} Two of its three instructions refer
  to things that are no longer in the code, or never were.
\item \textbf{Do not fix the five blue literals by changing the token.} The token
  is correct; the literals are wrong. Changing \code{--accent} back to blue would
  make five rules consistent and the entire site wrong.
\item \textbf{Do not assume this repository is what visitors see.} It is not.
\end{itemize}

\subsection{If you were starting over (judgement)}

\emph{Entirely opinion, argued from what the code shows.}

\textbf{Keep:} the no-build-step decision, the token block, the three-behaviour
JavaScript, the reduced-motion block (which is better here than in the sibling,
since it also stops the floating badges), and the two-column hero. For a page this
size every one of those is right, and a framework would have added a hundred times
the machinery for nothing a visitor could perceive.

\textbf{Do differently:} get the photograph first, because a personal brand page
without a photograph of the person is not a draft of the design, it is a draft of
the idea. Write a README on day one. Use a relative stylesheet path. Never write a
colour literally when a token exists, which would have made the rebrand a one-line
change instead of a six-line one with five lines missed. Write the real numbers in
the HTML and let the counter animate towards them, so the page never claims zero.
And state the relationship to \code{ai-digitalise} in both repositories.

\textbf{Drop:} \file{netlify.toml}, and the four unused tokens.

\textbf{Built today it would look almost identical.} The only structural change
worth making is moving the four security headers out of \file{netlify.toml} and
into the Caddy configuration on the VPS, where they would actually take effect.

\subsection{Open threads}

\begin{itemize}
\item \textbf{Eleven assets never delivered}, the photograph above all.
\item \textbf{Five blue literals} on an orange-branded site.
\item \textbf{Three dead stylesheet rules} for a class the HTML does not use.
\item \textbf{Seven stale comments} in \file{index.html}.
\item \textbf{The live ai.digitalise.agency canonical link} points at this domain,
  which is a bug in the portfolio repository that produced both live sites, and is
  recorded in this document because this is where the evidence for it was found.
\item \textbf{The relationship to \code{ai-digitalise} is undocumented} in both
  repositories.
\end{itemize}

% ===========================================================================
\section{Glossary}
% ===========================================================================

\begin{description}[leftmargin=0pt,style=nextline]
\item[ARIA] Attributes beginning \code{aria-} describing a page to screen readers.
  They change nothing visually.
\item[Attribute] Extra information inside an opening tag, as a name and a value.
\item[Breakpoint] A screen width at which the layout changes. This project has
  two: 960 and 600 pixels.
\item[Canonical URL] A declaration of where content authoritatively lives, used by
  search engines to avoid indexing it twice.
\item[Class] A label attached to elements so CSS and JavaScript can find them.
\item[clamp] A CSS function taking a minimum, a preferred value and a maximum.
\item[Commit] One recorded change to a repository, with a message explaining it.
\item[CSS] Cascading Style Sheets: the file controlling appearance.
\item[Custom property] A named value in CSS, defined \code{--name} and read with
  \code{var(--name)}.
\item[Data attribute] An attribute beginning \code{data-} that HTML ignores and
  JavaScript can read.
\item[Element] One thing on the page, produced by a tag.
\item[Flexbox] A layout system arranging children in a row or a column.
\item[Grid] A layout system arranging children in rows and columns at once.
\item[Hex colour] A colour written as six characters after a hash.
\item[HTML] HyperText Markup Language: content and structure.
\item[IntersectionObserver] A browser facility reporting when an element scrolls
  into view.
\item[JavaScript] The language running in the browser after the page loads.
\item[Media query] CSS applying only when a condition about the screen is true.
\item[order] A CSS property setting the position of an item inside a flex or grid
  container, independent of its position in the HTML.
\item[prefers-reduced-motion] A system setting asking pages to avoid animation.
\item[Pseudo-element] An extra box CSS creates with no tag in the HTML, written
  \code{::before} or \code{::after}.
\item[rsync] A program copying files between computers, transferring only what
  differs.
\item[Selector] The part of a CSS rule saying which elements it applies to.
\item[SSH] Secure Shell: the standard way to connect to another computer securely.
\item[SVG] A picture described as shapes and coordinates rather than pixels.
\item[Tag] The marker creating an element, written between angle brackets.
\item[Token] In this project, one of the nineteen named values at the top of the
  stylesheet.
\item[Viewport] The visible area of the page in the browser window.
\item[VPS] Virtual Private Server: a rented computer controlled remotely.
\item[Web server] A program answering requests for files over the network.
\item[Workflow] A file describing a job GitHub runs automatically.
\end{description}

% ===========================================================================
\appendix
\section{Appendix A: the file ledger}
% ===========================================================================

Every path git tracks, with its tier, its line count, and the section that
explains it. This is what makes the coverage claim checkable rather than merely
asserted: run \code{git ls-files} and compare.

\subsection*{\file{.github/}}
\begin{table}[H]\centering\small
\begin{tabular}{@{}llrl@{}}
\toprule
Path & Coverage & Lines & Section \\
\midrule
\file{.github/workflows-disabled/deploy.yml.disabled} & explained in full & 103 & 7.6 \\
\bottomrule
\end{tabular}
\caption{Coverage ledger: \file{.github/}. 1 file, 103 lines.}
\end{table}

\subsection*{\file{new-design-reference/}}
\begin{table}[H]\centering\small
\begin{tabular}{@{}llrl@{}}
\toprule
Path & Coverage & Lines & Section \\
\midrule
\file{new-design-reference/index.html}        & explained in full & 395 & 7.1 \\
\file{new-design-reference/css/site.css}      & explained in full & 575 & 7.2 \\
\file{new-design-reference/js/main.js}        & explained in full &  64 & 7.3 \\
\file{new-design-reference/netlify.toml}      & explained in full &  11 & 7.4 \\
\file{new-design-reference/assets/ASSETS.md}  & explained in full &  32 & 7.5 \\
\bottomrule
\end{tabular}
\caption{Coverage ledger: \file{new-design-reference/}. 5 files, 1{,}077 lines.}
\end{table}

\subsection*{Untracked directories (tier 3)}

\textbf{There are none.} \code{git status --porcelain --ignored} produces no
output at all: no ignored files, no build output, no dependency folder, and not
even a \file{.gitignore}. Every file on disk is tracked and every tracked file is
explained above.

\subsection*{The arithmetic}

\code{git ls-files} lists \textbf{6} paths. All \textbf{6} are tier 1, explained
in full, line by line, in section 7. \textbf{0} files are tier 2: there is no
generated, vendored, minified or binary content, checked with a search for
generator banners that returned nothing. \textbf{0} directories are tier 3,
because nothing on disk is untracked. Total lines accounted for:
\textbf{1{,}180}, the sum reported by
\code{git ls-files -z | xargs -0 wc -l}.

$6 + 0 = 6$, and $103 + 1{,}077 = 1{,}180$. The ledger reconciles.

\end{document}
